\input{preamble}

% D'accord, commençons ici.
%
\begin{document}

\title{Faisceaux sur les champs algébriques}


\maketitle

\phantomsection
\label{section-phantom}

\tableofcontents

\section{Introduction}
\label{section-introduction}

\noindent
Il existe une multitude de façons d'envisager les faisceaux sur les champs algébriques.
Dans ce chapitre, nous présentons une approche particulièrement bien
adaptée à nos fondements pour les champs algébriques. Chaque fois que nous introduirons
un type de faisceaux, nous préciserons sa relation exacte avec les
notions analogues de la littérature.
Le but de ce chapitre est d'énoncer les résultats
qui sont soit évidents, soit faciles à démontrer,
et de réserver les constructions plus délicates pour plus tard.

\medskip\noindent
En fait, développer une théorie complète des faisceaux étales
constructibles, ou encore une étude adéquate des catégories dérivées
de complexes de $\mathcal{O}$-modules dont les faisceaux de cohomologie sont
quasi-cohérents, exige un travail considérable ; voir
\cite{olsson_sheaves}. Nous y reviendrons dans
Cohomologie des champs, Section \ref{stacks-cohomology-section-introduction}.

\medskip\noindent
Dans la littérature et les articles de recherche consacrés aux faisceaux sur les champs
algébriques, le site lisse-étale d'un champ algébrique joue souvent un rôle majeur.
Il s'agit toutefois d'un objet problématique : un morphisme de
champs algébriques n'induit pas, en général, de morphisme de topos lisses-étales. Nous avons
donc choisi d'éviter aussi longtemps que possible toute mention du site
lisse-étale. Les arguments qui utilisent traditionnellement ce site
seront remplacés par des arguments faisant appel à un recouvrement de {\v C}ech
dans le site $\mathcal{X}_{smooth}$ défini ci-dessous.

\medskip\noindent
Certaines notations, conventions et certains termes de ce chapitre sont peu commodes
et pourront sembler contraires aux usages au lecteur expérimenté. C'est intentionnel.
On en trouvera une explication dans Quotients, Section
\ref{quot-section-conventions}.




\section{Conventions}
\label{section-conventions}

\noindent
Les conventions employées dans ce chapitre sont les mêmes que dans le
chapitre sur les champs algébriques ; voir
Champs algébriques, Section \ref{algebraic-section-conventions}.
Nous les répétons ici par commodité.

\medskip\noindent
Nous travaillons dans un grand site fppf convenable $\Sch_{fppf}$, comme dans
Topologies, Définition \ref{topologies-definition-big-fppf-site}.
Ainsi, sauf indication contraire explicite, tous les schémas seront des objets
de $\Sch_{fppf}$. Nous indiquerons ailleurs ce qui change lorsqu'on remplace le
grand site fppf (insérer ici une référence future).

\medskip\noindent
Nous travaillerons toujours relativement à une base $S$ appartenant à $\Sch_{fppf}$.
Nous utiliserons alors le grand site fppf $(\Sch/S)_{fppf}$ ; voir
Topologies, Définition \ref{topologies-definition-big-small-fppf}.
On retrouve le cas absolu en prenant
$S = \Spec(\mathbf{Z})$.





\section{Préfaisceaux}
\label{section-presheaves}

\noindent
Dans cette section, nous définissons les préfaisceaux sur les catégories fibrées en
groupoïdes au-dessus de $(\Sch/S)_{fppf}$, mais l'essentiel de la discussion vaut
pour les catégories au-dessus d'une catégorie de base quelconque. Cette section sert aussi à
introduire les notations que nous emploierons par la suite.

\begin{definition}
\label{definition-presheaves}
Soit $p : \mathcal{X} \to (\Sch/S)_{fppf}$ une catégorie fibrée en
groupoïdes.
\begin{enumerate}
\item Un {\it préfaisceau sur $\mathcal{X}$} est un préfaisceau sur la
catégorie sous-jacente à $\mathcal{X}$.
\item Un {\it morphisme de préfaisceaux sur $\mathcal{X}$} est un morphisme de
préfaisceaux sur la catégorie sous-jacente à $\mathcal{X}$.
\end{enumerate}
Nous notons $\textit{PSh}(\mathcal{X})$ la catégorie des préfaisceaux sur
$\mathcal{X}$.
\end{definition}

\noindent
Cela définit les préfaisceaux d'ensembles. On peut bien entendu parler aussi de
préfaisceaux d'ensembles pointés, de groupes abéliens, de groupes, de monoïdes, d'anneaux,
de modules sur un anneau fixé, d'algèbres de Lie sur un corps fixé, etc.
La catégorie des {\it préfaisceaux abéliens}, c'est-à-dire des préfaisceaux de groupes
abéliens, est notée $\textit{PAb}(\mathcal{X})$.

\medskip\noindent
Soit $f : \mathcal{X} \to \mathcal{Y}$ un $1$-morphisme de catégories
fibrées en groupoïdes au-dessus de $(\Sch/S)_{fppf}$. Rappelons que cela
signifie simplement que $f$ est un foncteur au-dessus de $(\Sch/S)_{fppf}$.
Les résultats de
Sites, Section \ref{sites-section-more-functoriality-PSh},
nous fournissent une paire de foncteurs adjoints\footnote{Ces foncteurs
seront notés $f^{-1}$ et $f_*$ une fois le
Lemme \ref{lemma-functoriality-sheaves}
démontré.}
\begin{equation}
\label{equation-pushforward-pullback}
f^p : \textit{PSh}(\mathcal{Y}) \longrightarrow \textit{PSh}(\mathcal{X})
\quad\text{et}\quad
{}_pf : \textit{PSh}(\mathcal{X}) \longrightarrow \textit{PSh}(\mathcal{Y}).
\end{equation}
L'adjonction s'écrit
$$
\Mor_{\textit{PSh}(\mathcal{X})}(f^p\mathcal{G}, \mathcal{F})
=
\Mor_{\textit{PSh}(\mathcal{Y})}(\mathcal{G}, {}_pf\mathcal{F})
$$
où $\mathcal{F} \in \Ob(\textit{PSh}(\mathcal{X}))$ et
$\mathcal{G} \in \Ob(\textit{PSh}(\mathcal{Y}))$. Nous appelons
$f^p\mathcal{G}$ l'{\it image réciproque} de $\mathcal{G}$. Il résulte
des définitions que
$$
f^p\mathcal{G}(x) = \mathcal{G}(f(x))
$$
pour tout $x \in \Ob(\mathcal{X})$. Le préfaisceau ${}_pf\mathcal{F}$
est appelé l'{\it image directe} de $\mathcal{F}$. Il est décrit
par la formule
$$
({}_pf\mathcal{F})(y) = \lim_{f(x) \to y} \mathcal{F}(x).
$$
Le reste de cette section devrait probablement être déplacé dans le chapitre
sur les sites et peut, en tout cas, être omis en première lecture.

\begin{lemma}
\label{lemma-1-morphisms-presheaves}
Soient $f : \mathcal{X} \to \mathcal{Y}$ et $g : \mathcal{Y} \to \mathcal{Z}$
des $1$-morphismes de catégories fibrées en groupoïdes au-dessus de
$(\Sch/S)_{fppf}$. Alors $(g \circ f)^p = f^p \circ g^p$ et
il existe un isomorphisme canonique
${}_p(g \circ f) \to {}_pg \circ {}_pf$
compatible aux adjonctions de $(f^p, {}_pf)$, $(g^p, {}_pg)$ et
$((g \circ f)^p, {}_p(g \circ f))$.
\end{lemma}

\begin{proof}
Soit $\mathcal{H}$ un préfaisceau sur $\mathcal{Z}$. Alors
$(g \circ f)^p\mathcal{H} = f^p (g^p\mathcal{H})$ est donné
par les égalités
$$
(g \circ f)^p\mathcal{H}(x) = \mathcal{H}((g \circ f)(x))
= \mathcal{H}(g(f(x))) = f^p (g^p\mathcal{H})(x).
$$
Nous omettons de vérifier la compatibilité avec les applications de restriction.

\medskip\noindent
Définissons ensuite la transformation ${}_p(g \circ f) \to {}_pg \circ {}_pf$.
Soit $\mathcal{F}$ un préfaisceau sur $\mathcal{X}$.
Si $z$ est un objet de $\mathcal{Z}$, on obtient une
catégorie $\mathcal{J}$ de quadruplets
$(x, f(x) \to y, y, g(y) \to z)$ et une catégorie $\mathcal{I}$
de couples $(x, g(f(x)) \to z)$. Il existe un foncteur canonique
$\mathcal{J} \to \mathcal{I}$ qui envoie l'objet
$(x, \alpha : f(x) \to y, y, \beta : g(y) \to z)$ sur
$(x, \beta \circ f(\alpha) : g(f(x)) \to z)$. Il fournit la flèche dans
\begin{align*}
({}_p(g \circ f)\mathcal{F})(z) & =
\lim_{g(f(x)) \to z} \mathcal{F}(x) \\
& = \lim_\mathcal{I} \mathcal{F} \\
& \to \lim_\mathcal{J} \mathcal{F} \\
& = \lim_{g(y) \to z}
\Big(\lim_{f(x) \to y} \mathcal{F}(x)\Big) \\
& =
({}_pg \circ {}_pf\mathcal{F})(x)
\end{align*}
par
Catégories, Lemme \ref{categories-lemma-functorial-limit}.
Nous omettons de vérifier la compatibilité avec les applications de restriction.
Au lieu de cette construction directe, on peut définir
${}_p(g \circ f) \cong {}_pg \circ {}_pf$
comme l'unique application compatible aux propriétés d'adjonction. Cette méthode
présente aussi l'avantage de dispenser de démontrer la compatibilité.

\medskip\noindent
La compatibilité aux adjonctions de $(f^p, {}_pf)$, $(g^p, {}_pg)$ et
$((g \circ f)^p, {}_p(g \circ f))$ signifie que, pour les préfaisceaux
$\mathcal{H}$ et $\mathcal{F}$ ci-dessus, le diagramme suivant est commutatif :
$$
\xymatrix{
\Mor_{\textit{PSh}(\mathcal{X})}(f^pg^p\mathcal{H}, \mathcal{F})
\ar@{=}[r] \ar@{=}[d] &
\Mor_{\textit{PSh}(\mathcal{Y})}(g^p\mathcal{H}, {}_pf\mathcal{F})
\ar@{=}[r] &
\Mor_{\textit{PSh}(\mathcal{Y})}(\mathcal{H}, {}_pg{}_pf\mathcal{F})
\\
\Mor_{\textit{PSh}(\mathcal{X})}((g \circ f)^p\mathcal{G}, \mathcal{F})
\ar@{=}[rr] & &
\Mor_{\textit{PSh}(\mathcal{Y})}(\mathcal{G}, {}_p(g \circ f)\mathcal{F})
\ar[u]
}
$$
Démonstration omise.
\end{proof}

\begin{lemma}
\label{lemma-2-morphisms-presheaves}
Soient $f, g : \mathcal{X} \to \mathcal{Y}$ des $1$-morphismes de catégories
fibrées en groupoïdes au-dessus de $(\Sch/S)_{fppf}$. Soit $t : f \to g$
un $2$-morphisme de catégories fibrées en groupoïdes au-dessus de
$(\Sch/S)_{fppf}$. À $t$ sont associés des
isomorphismes canoniques de foncteurs
$$
t^p : g^p \longrightarrow f^p
\quad\text{et}\quad
{}_pt : {}_pf \longrightarrow {}_pg
$$
compatibles aux adjonctions de $(f^p, {}_pf)$ et
$(g^p, {}_pg)$, ainsi qu'aux
compositions verticale et horizontale des $2$-morphismes.
\end{lemma}

\begin{proof}
Soit $\mathcal{G}$ un préfaisceau sur $\mathcal{Y}$. Alors
$t^p : g^p\mathcal{G} \to f^p\mathcal{G}$ est donné par la famille
d'applications
$$
g^p\mathcal{G}(x) = \mathcal{G}(g(x))
\xrightarrow{\mathcal{G}(t_x)}
\mathcal{G}(f(x)) = f^p\mathcal{G}(x)
$$
paramétrée par $x \in \Ob(\mathcal{X})$. Cela a un sens puisque
$t_x : f(x) \to g(x)$ et que $\mathcal{G}$ est un foncteur contravariant.
Nous omettons de vérifier la compatibilité avec les applications
de restriction.

\medskip\noindent
Pour définir la transformation ${}_pt$, posons, pour $y \in \Ob(\mathcal{Y})$,
que ${}_y^f\mathcal{I}$, resp.\ ${}_y^g\mathcal{I}$, est la catégorie
des couples $(x, \psi : f(x) \to y)$, resp.\ $(x, \psi : g(x) \to y)$ ; voir
Sites, Section \ref{sites-section-more-functoriality-PSh}.
Remarquons que $t$ définit un foncteur
${}_yt : {}_y^g\mathcal{I} \to {}_y^f\mathcal{I}$
donné par la règle
$$
(x, g(x) \to y) \longmapsto (x, f(x) \xrightarrow{t_x} g(x) \to y).
$$
Pour un préfaisceau $\mathcal{F}$ sur $\mathcal{X}$, remarquons que le composé
de ${}_yt$ avec $\mathcal{F} : {}_y^f\mathcal{I}^{opp} \to \textit{Sets}$,
$(x, f(x) \to y) \mapsto \mathcal{F}(x)$, est égal à
$\mathcal{F} : {}_y^g\mathcal{I}^{opp} \to \textit{Sets}$. Ainsi, d'après
Catégories, Lemme \ref{categories-lemma-functorial-limit},
nous obtenons pour tout $y \in \Ob(\mathcal{Y})$ une application canonique
$$
({}_pf\mathcal{F})(y) = \lim_{{}_y^f\mathcal{I}} \mathcal{F}
\longrightarrow
\lim_{{}_y^g\mathcal{I}} \mathcal{F} = ({}_pg\mathcal{F})(y)
$$
Nous omettons de vérifier la compatibilité avec les applications
de restriction. Au lieu de cette construction directe, on peut définir
${}_pt$ comme l'unique application compatible aux propriétés d'adjonction
des couples $(f^p, {}_pf)$ et $(g^p, {}_pg)$ (voir ci-dessous). Cette méthode
présente aussi l'avantage de dispenser de démontrer la compatibilité.

\medskip\noindent
La compatibilité aux adjonctions de $(f^p, {}_pf)$ et $(g^p, {}_pg)$ signifie
que, pour les préfaisceaux $\mathcal{G}$ et $\mathcal{F}$ ci-dessus, nous avons
un diagramme commutatif
$$
\xymatrix{
\Mor_{\textit{PSh}(\mathcal{X})}(f^p\mathcal{G}, \mathcal{F})
\ar@{=}[r] \ar[d]_{- \circ t^p} &
\Mor_{\textit{PSh}(\mathcal{Y})}(\mathcal{G}, {}_pf\mathcal{F})
\ar[d]^{{}_pt \circ -} \\
\Mor_{\textit{PSh}(\mathcal{X})}(g^p\mathcal{G}, \mathcal{F})
\ar@{=}[r] &
\Mor_{\textit{PSh}(\mathcal{Y})}(\mathcal{G}, {}_pg\mathcal{F})
}
$$
Démonstration omise. Indication : reprendre la démonstration de
Sites, Lemme \ref{sites-lemma-adjoints-pu},
et déduire la compatibilité de la description explicite des
applications horizontales et verticales du diagramme.

\medskip\noindent
Nous omettons de vérifier la compatibilité avec les compositions verticale et
horizontale. Indication : pour $t^p$, la démonstration est immédiate, et
la compatibilité à l'adjonction permet d'en déduire le résultat pour les applications
${}_pt$.
\end{proof}







\section{Faisceaux}
\label{section-sheaves}

\noindent
Commençons par une observation importante et triviale
(surtout pour les lecteurs que les questions de théorie des ensembles
ne préoccupent pas).

\medskip\noindent
Considérons un grand site fppf $\Sch_{fppf}$ comme dans
Topologies, Définition \ref{topologies-definition-big-fppf-site},
et notons $\Sch_\alpha$ sa catégorie sous-jacente.
Outre qu'elle est la catégorie sous-jacente à un site fppf,
la catégorie $\Sch_\alpha$ peut aussi servir de catégorie
sous-jacente à un grand site de Zariski, un grand site étale, un grand site lisse
et un grand site syntomique ; voir
Topologies, Remarque \ref{topologies-remark-choice-sites}.
Nous notons ces sites $\Sch_{Zar}$, $\Sch_\etale$,
$\Sch_{smooth}$ et $\Sch_{syntomic}$.
Dans cette situation, puisque nous avons défini
le grand site de Zariski $(\Sch/S)_{Zar}$ de $S$,
le grand site étale $(\Sch/S)_\etale$ de $S$,
le grand site lisse $(\Sch/S)_{smooth}$ de $S$,
le grand site syntomique $(\Sch/S)_{syntomic}$ de $S$ et
le grand site fppf $(\Sch/S)_{fppf}$ de $S$
comme les localisations (voir
Sites, Section \ref{sites-section-localize})
$\Sch_{Zar}/S$, $\Sch_\etale/S$,
$\Sch_{smooth}/S$, $\Sch_{syntomic}/S$ et
$\Sch_{fppf}/S$
de ces grands sites (absolus), tous ont la
même catégorie sous-jacente, à savoir $\Sch_\alpha/S$.

\medskip\noindent
Il s'ensuit que, si l'on se donne une catégorie
$p : \mathcal{X} \to (\Sch/S)_{fppf}$ fibrée en groupoïdes, alors
$\mathcal{X}$ hérite de topologies de Zariski, étale, lisse, syntomique et
fppf ; voir
Champs, Définition \ref{stacks-definition-topology-inherited}.

\begin{definition}
\label{definition-inherited-topologies}
Soit $\mathcal{X}$ une catégorie fibrée en groupoïdes au-dessus de
$(\Sch/S)_{fppf}$.
\begin{enumerate}
\item Le {\it site de Zariski associé}, noté $\mathcal{X}_{Zar}$,
est la structure de site sur $\mathcal{X}$ héritée de
$(\Sch/S)_{Zar}$.
\item Le {\it site étale associé}, noté $\mathcal{X}_\etale$,
est la structure de site sur $\mathcal{X}$ héritée de
$(\Sch/S)_\etale$.
\item Le {\it site lisse associé}, noté $\mathcal{X}_{smooth}$,
est la structure de site sur $\mathcal{X}$ héritée de
$(\Sch/S)_{smooth}$.
\item Le {\it site syntomique associé}, noté $\mathcal{X}_{syntomic}$,
est la structure de site sur $\mathcal{X}$ héritée de
$(\Sch/S)_{syntomic}$.
\item Le {\it site fppf associé}, noté $\mathcal{X}_{fppf}$,
est la structure de site sur $\mathcal{X}$ héritée de
$(\Sch/S)_{fppf}$.
\end{enumerate}
\end{definition}

\noindent
La discussion précédente donne un sens à cette définition. Si $\mathcal{X}$
est un champ algébrique, la littérature appelle $\mathcal{X}_{fppf}$ (ou un
site qui lui est équivalent) le {\it grand site fppf} de $\mathcal{X}$, et de même
pour les autres. Nous emploierons parfois cette terminologie afin de
distinguer cette construction d'autres constructions.

\begin{remark}
\label{remark-ambiguity}
Nous n'employons cette notation que lorsque le symbole $\mathcal{X}$ désigne une
catégorie fibrée en groupoïdes, et non un schéma, un espace algébrique, etc.
Nous éviterons ainsi toute confusion avec le petit site étale d'un
schéma ou d'un espace algébrique, que l'on note $X_\etale$ (auquel cas
la lettre capitale est romaine et non calligraphique).
\end{remark}

\noindent
Ces topologies étant définies, nous pouvons préciser ce qu'est
un faisceau sur $\mathcal{X}$, c'est-à-dire définir les topos correspondants.

\begin{definition}
\label{definition-sheaves}
Soit $\mathcal{X}$ une catégorie fibrée en groupoïdes au-dessus de
$(\Sch/S)_{fppf}$. Soit $\mathcal{F}$ un préfaisceau sur $\mathcal{X}$.
\begin{enumerate}
\item Nous disons que $\mathcal{F}$ est un {\it faisceau de Zariski}, ou un
{\it faisceau pour la topologie de Zariski}, si $\mathcal{F}$
est un faisceau sur le site de Zariski associé $\mathcal{X}_{Zar}$.
\item Nous disons que $\mathcal{F}$ est un {\it faisceau étale}, ou un
{\it faisceau pour la topologie étale}, si $\mathcal{F}$
est un faisceau sur le site étale associé $\mathcal{X}_\etale$.
\item Nous disons que $\mathcal{F}$ est un {\it faisceau lisse}, ou un
{\it faisceau pour la topologie lisse}, si $\mathcal{F}$
est un faisceau sur le site lisse associé $\mathcal{X}_{smooth}$.
\item Nous disons que $\mathcal{F}$ est un {\it faisceau syntomique}, ou un
{\it faisceau pour la topologie syntomique}, si $\mathcal{F}$
est un faisceau sur le site syntomique associé $\mathcal{X}_{syntomic}$.
\item Nous disons que $\mathcal{F}$ est un {\it faisceau fppf}, ou un {\it faisceau},
ou encore un {\it faisceau pour la topologie fppf}, si $\mathcal{F}$
est un faisceau sur le site fppf associé $\mathcal{X}_{fppf}$.
\end{enumerate}
Un morphisme de faisceaux est simplement un morphisme de préfaisceaux. Nous notons
ces catégories de faisceaux
$\Sh(\mathcal{X}_{Zar})$,
$\Sh(\mathcal{X}_\etale)$,
$\Sh(\mathcal{X}_{smooth})$,
$\Sh(\mathcal{X}_{syntomic})$ et
$\Sh(\mathcal{X}_{fppf})$.
\end{definition}

\noindent
On peut bien entendu parler aussi de faisceaux d'ensembles pointés, de groupes abéliens,
de groupes, de monoïdes, d'anneaux, de modules sur un anneau fixé et d'algèbres de Lie sur
un corps fixé, etc. La catégorie des {\it faisceaux abéliens}, c'est-à-dire des faisceaux
de groupes abéliens, est notée $\textit{Ab}(\mathcal{X}_{fppf})$,
et de même pour les autres topologies. Si $\mathcal{X}$ est un champ
algébrique, alors $\Sh(\mathcal{X}_{fppf})$ est équivalente (aux
questions ensemblistes près) à ce que la littérature appelle
la {\it catégorie des faisceaux sur le grand site fppf de $\mathcal{X}$}. Il en va de même
pour les autres topologies. Nous emploierons parfois cette terminologie afin de
distinguer cette construction d'autres constructions.

\medskip\noindent
Puisque les topologies sont énumérées par finesse croissante, nous avons
les inclusions strictement pleines suivantes :
$$
\Sh(\mathcal{X}_{fppf}) \subset
\Sh(\mathcal{X}_{syntomic}) \subset
\Sh(\mathcal{X}_{smooth}) \subset
\Sh(\mathcal{X}_\etale) \subset
\Sh(\mathcal{X}_{Zar}) \subset \textit{PSh}(\mathcal{X})
$$
Nous écrivons parfois
$\Sh(\mathcal{X}_{fppf}) = \Sh(\mathcal{X})$
and
$\textit{Ab}(\mathcal{X}_{fppf}) = \textit{Ab}(\mathcal{X})$
conformément à notre convention selon laquelle un faisceau sur $\mathcal{X}$
est un faisceau fppf sur $\mathcal{X}$.

\medskip\noindent
Dans ce cadre, la fonctorialité de ces topos est immédiate et
elle est, de plus, compatible aux foncteurs d'inclusion ci-dessus.

\begin{lemma}
\label{lemma-functoriality-sheaves}
Soit $f : \mathcal{X} \to \mathcal{Y}$ un $1$-morphisme de catégories
fibrées en groupoïdes au-dessus de $(\Sch/S)_{fppf}$. Soit
$\tau \in \{Zar, \etale, smooth, syntomic, fppf\}$.
Les foncteurs ${}_pf$ et $f^p$ de (\ref{equation-pushforward-pullback})
transforment les faisceaux pour $\tau$ en faisceaux pour $\tau$ et définissent un morphisme
de topos
$f : \Sh(\mathcal{X}_\tau) \to \Sh(\mathcal{Y}_\tau)$.
\end{lemma}

\begin{proof}
Cela résulte immédiatement de
Champs, Lemme \ref{stacks-lemma-topology-inherited-functorial}.
\end{proof}

\noindent
Autrement dit, l'image directe et l'image réciproque des préfaisceaux définies dans la
Section \ref{section-presheaves}
donnent aussi l'{\it image directe} et l'{\it image réciproque} des faisceaux pour $\tau$.
Compte tenu de ce qui précède, nous pouvons écrire $f^p = f^{-1}$
et ${}_pf = f_*$ sans risque de confusion.

\begin{definition}
\label{definition-morphism}
Soit $f : \mathcal{X} \to \mathcal{Y}$ un morphisme de catégories
fibrées en groupoïdes au-dessus de $(\Sch/S)_{fppf}$. Nous notons
$$
f = (f^{-1}, f_*) :
\Sh(\mathcal{X}_{fppf})
\longrightarrow
\Sh(\mathcal{Y}_{fppf})
$$
le {\it morphisme de topos fppf associé} construit ci-dessus.
De même pour les topos de Zariski, étale, lisse et syntomique associés.
\end{definition}

\noindent
Comme expliqué dans
Sites, Section \ref{sites-section-sheaves-algebraic-structures},
la même formule (appliquée au faisceau d'ensembles sous-jacent) définit
l'image directe et l'image réciproque des faisceaux (pour l'une de nos topologies)
d'ensembles pointés, de groupes abéliens, de groupes, de monoïdes, d'anneaux, de modules
sur un anneau fixé, d'algèbres de Lie sur un corps fixé, etc.








\section{Calcul de l'image directe}
\label{section-pushforward}

\noindent
Soit $f : \mathcal{X} \to \mathcal{Y}$ un $1$-morphisme de catégories
fibrées en groupoïdes au-dessus de $(\Sch/S)_{fppf}$. Soit $\mathcal{F}$
un préfaisceau sur $\mathcal{X}$. Soit $y \in \Ob(\mathcal{Y})$.
On peut calculer $f_*\mathcal{F}(y)$ comme suit. Supposons que
$y$ soit au-dessus du schéma $V$ et, au moyen du lemme de $2$-Yoneda, considérons
$y$ comme un $1$-morphisme. Considérons la projection
$$
\text{pr} :
(\Sch/V)_{fppf} \times_{y, \mathcal{Y}} \mathcal{X}
\longrightarrow
\mathcal{X}
$$
On a alors une identification canonique
\begin{equation}
\label{equation-pushforward}
f_*\mathcal{F}(y) = \Gamma\Big(
(\Sch/V)_{fppf} \times_{y, \mathcal{Y}} \mathcal{X},
\ \text{pr}^{-1}\mathcal{F}\Big)
\end{equation}
En effet, les objets du $2$-produit fibré sont les triplets
$(h : U \to V, x, f(x) \to h^*y)$. En omettant $h$ dans la
notation, cela équivaut à la donnée d'un objet
$x$ de $\mathcal{X}$ et d'un morphisme $\alpha : f(x) \to y$ de $\mathcal{Y}$.
Comme $f_*\mathcal{F}(y) = \lim_{f(x) \to y} \mathcal{F}(x)$ par définition,
l'égalité en résulte.

\medskip\noindent
Nous en déduisons le résultat de « changement de base » suivant pour les
images directes. Il est trivial et repose sur le fait que
nous utilisons des « grands » sites.

\begin{lemma}
\label{lemma-base-change}
Soit $S$ un schéma. Soit
$$
\xymatrix{
\mathcal{Y}' \times_\mathcal{Y} \mathcal{X} \ar[r]_{g'} \ar[d]_{f'} &
\mathcal{X} \ar[d]^f \\
\mathcal{Y}' \ar[r]^g & \mathcal{Y}
}
$$
un diagramme $2$-cartésien de catégories fibrées en groupoïdes au-dessus de $S$.
On a alors un isomorphisme canonique
$$
g^{-1}f_*\mathcal{F} \longrightarrow f'_*(g')^{-1}\mathcal{F}
$$
fonctoriel en le préfaisceau $\mathcal{F}$ sur $\mathcal{X}$.
\end{lemma}

\begin{proof}
Pour tout objet $y'$ de $\mathcal{Y}'$ au-dessus de $V$,
il existe une équivalence
$$
(\Sch/V)_{fppf} \times_{g(y'), \mathcal{Y}} \mathcal{X}
=
(\Sch/V)_{fppf} \times_{y', \mathcal{Y}'}
(\mathcal{Y}' \times_\mathcal{Y} \mathcal{X})
$$
D'où, par (\ref{equation-pushforward}), une bijection
$g^{-1}f_*\mathcal{F}(y') \to f'_*(g')^{-1}\mathcal{F}(y')$.
Nous omettons de vérifier la compatibilité avec les applications
de restriction.
\end{proof}

\noindent
Dans le cas d'un morphisme représentable de catégories fibrées en groupoïdes,
la formule (\ref{equation-pushforward}) se simplifie. Nous conseillons au
lecteur d'omettre le reste de cette section.

\begin{lemma}
\label{lemma-representable}
Soit $f : \mathcal{X} \to \mathcal{Y}$ un $1$-morphisme de catégories
fibrées en groupoïdes au-dessus de $(\Sch/S)_{fppf}$. Les conditions suivantes sont
équivalentes :
\begin{enumerate}
\item $f$ est représentable, et
\item pour tout $y \in \Ob(\mathcal{Y})$, le foncteur
$\mathcal{X}^{opp} \to \textit{Sets}$,
$x \mapsto \Mor_\mathcal{Y}(f(x), y)$
est représentable.
\end{enumerate}
\end{lemma}

\begin{proof}
Selon la discussion de
Champs algébriques, Section \ref{algebraic-section-representable-morphism},
$f$ est représentable si et seulement si, pour tout
$y \in \Ob(\mathcal{Y})$
au-dessus de $U$, le $2$-produit fibré
$(\Sch/U)_{fppf} \times_{y, \mathcal{Y}} \mathcal{X}$
est représentable, c'est-à-dire de la forme $(\Sch/V_y)_{fppf}$ pour un
schéma $V_y$ au-dessus de $U$. Les objets de ce $2$-produit fibré sont les triplets
$(h : V \to U, x, \alpha : f(x) \to h^*y)$ où $\alpha$ est
au-dessus de $\text{id}_V$. En omettant $h$ dans la notation, cela
équivaut à la donnée d'un objet $x$ de $\mathcal{X}$ et d'un
morphisme $f(x) \to y$. Ainsi, le $2$-produit fibré est
représenté par $V_y$ et $f(x_y) \to y$, où $x_y$ est un objet
de $\mathcal{X}$ au-dessus de $V_y$, si et seulement si le foncteur de (2) est représenté
par $x_y$, l'objet universel étant une application $f(x_y) \to y$.
\end{proof}

\noindent
Let
$$
\xymatrix{
\mathcal{X} \ar[rr]_f \ar[rd]_p & &  \mathcal{Y} \ar[ld]^q \\
& (\Sch/S)_{fppf}
}
$$
un $1$-morphisme de catégories fibrées en groupoïdes. Supposons $f$
représentable. Pour tout $y \in \Ob(\mathcal{Y})$, choisissons
un objet $u(y) \in \Ob(\mathcal{X})$ représentant le foncteur
$x \mapsto \Mor_\mathcal{Y}(f(x), y)$ of
Lemma \ref{lemma-representable}
(ce qui est possible par l'axiome du choix).
Par construction, ces objets sont munis de morphismes canoniques
$f(u(y)) \to y$.
Pour tout morphisme $\beta : y' \to y$ de $\mathcal{Y}$, on obtient un unique
morphisme $u(\beta) : u(y') \to u(y)$ de $\mathcal{X}$ tel que le diagramme
$$
\xymatrix{
f(u(y')) \ar[d] \ar[rr]_{f(u(\beta))} & & f(u(y)) \ar[d] \\
y' \ar[rr] & & y
}
$$
soit commutatif. Autrement dit, $u : \mathcal{Y} \to \mathcal{X}$ est un foncteur.
On peut en fait en dire un peu plus. Supposons en effet que
$V' = q(y')$, $V = q(y)$, $U' = p(u(y'))$ and $U = p(u(y))$. Then
$$
\xymatrix{
U' \ar[rr]_{p(u(\beta))} \ar[d] & & U \ar[d] \\
V' \ar[rr]^{q(\beta)} & & V
}
$$
soit un carré cartésien. Cela tient à ce que $U' \to U$ représente
le changement de base
$(\Sch/V')_{fppf} \times_{y', \mathcal{Y}} \mathcal{X} \to
(\Sch/V)_{fppf} \times_{y, \mathcal{Y}} \mathcal{X}$
de $V' \to V$.

\begin{lemma}
\label{lemma-representable-pushforward}
Soit $f : \mathcal{X} \to \mathcal{Y}$ un $1$-morphisme représentable de
catégories fibrées en groupoïdes au-dessus de $(\Sch/S)_{fppf}$. Soit
$\tau \in \{Zar, \etale, smooth, syntomic, fppf\}$.
Alors le foncteur $u : \mathcal{Y}_\tau \to \mathcal{X}_\tau$ est continu
et définit un morphisme de sites $\mathcal{X}_\tau \to \mathcal{Y}_\tau$
qui induit le même morphisme de topos
$\Sh(\mathcal{X}_\tau) \to \Sh(\mathcal{Y}_\tau)$
que le morphisme $f$ construit dans le
Lemme \ref{lemma-functoriality-sheaves}.
De plus, $f_*\mathcal{F}(y) = \mathcal{F}(u(y))$ pour tout préfaisceau
$\mathcal{F}$ sur $\mathcal{X}$.
\end{lemma}

\begin{proof}
Soit $\{y_i \to y\}$ un recouvrement pour $\tau$ dans $\mathcal{Y}$. Par définition,
cela signifie simplement que $\{q(y_i) \to q(y)\}$ est un recouvrement pour $\tau$ de
schémas. La dernière remarque précédant le lemme montre que
$\{p(u(y_i)) \to p(u(y))\}$ est le changement de base du recouvrement pour $\tau$
$\{q(y_i) \to q(y)\}$ par $p(u(y)) \to q(y)$ ; c'est donc lui-même un
recouvrement pour $\tau$ par les axiomes d'un site. Ainsi $\{u(y_i) \to u(y)\}$
est un recouvrement pour $\tau$ de $\mathcal{X}$. Cela démontre que $u$ est
continu.

\medskip\noindent
Employons les notations $u_p, u_s, u^p, u^s$ de
Sites, Sections \ref{sites-section-functoriality-PSh} et
\ref{sites-section-continuous-functors}.
Si nous démontrons la dernière assertion du lemme, alors
$f_* = u^p = u^s$ (par la continuité de $u$ établie ci-dessus) ; l'adjonction donne donc
$f^{-1} = u_s$, ce qui prouve que $u_s$ est exact, puis que $u$ détermine
un morphisme de sites ; l'égalité sera elle aussi claire.
Pour voir que $f_*\mathcal{F}(y) = \mathcal{F}(u(y))$, remarquons que, par
définition,
$$
f_*\mathcal{F}(y) = ({}_pf\mathcal{F})(y) =
\lim_{f(x) \to y} \mathcal{F}(x).
$$
Puisque $u(y)$ est un objet final de la catégorie sur laquelle la limite est
prise, on conclut.
\end{proof}





\section{Le faisceau structural}
\label{section-structure-sheaf}

\noindent
Soit $\tau \in \{Zar, \etale, smooth, syntomic, fppf\}$.
Soit $p : \mathcal{X} \to (\Sch/S)_{fppf}$ une catégorie
fibrée en groupoïdes. La $2$-catégorie des catégories fibrées en groupoïdes au-dessus de
$(\Sch/S)_{fppf}$ possède un objet final, à savoir
$\text{id} : (\Sch/S)_{fppf} \to (\Sch/S)_{fppf}$
et $p$ est un $1$-morphisme de $\mathcal{X}$ vers cet objet final.
Ainsi, tout préfaisceau $\mathcal{G}$ sur $(\Sch/S)_{fppf}$ fournit un
préfaisceau $p^{-1}\mathcal{G}$ sur $\mathcal{X}$ défini par la règle
$p^{-1}\mathcal{G}(x) = \mathcal{G}(p(x))$. De plus, la discussion de la
Section \ref{section-sheaves}
montre que $p^{-1}\mathcal{G}$ est un faisceau pour $\tau$ dès que
$\mathcal{G}$ en est un.

\medskip\noindent
Rappelons que le site $(\Sch/S)_{fppf}$ est un site annelé,
de faisceau structural $\mathcal{O}$ défini par la règle
$$
(\Sch/S)^{opp} \longrightarrow \textit{Rings},
\quad
U/S \longmapsto \Gamma(U, \mathcal{O}_U)
$$
voir
Descente, Définition \ref{descent-definition-structure-sheaf}.

\begin{definition}
\label{definition-structure-sheaf}
Soit $p : \mathcal{X} \to (\Sch/S)_{fppf}$ une catégorie
fibrée en groupoïdes. Le
{\it faisceau structural de $\mathcal{X}$} est le faisceau d'anneaux
$\mathcal{O}_\mathcal{X} = p^{-1}\mathcal{O}$.
\end{definition}

\noindent
Pour un objet $x$ de $\mathcal{X}$ au-dessus de $U$, nous avons
$\mathcal{O}_\mathcal{X}(x) = \mathcal{O}(U) = \Gamma(U, \mathcal{O}_U)$.
Il va sans dire que $\mathcal{O}_\mathcal{X}$ est aussi un faisceau de Zariski, étale,
lisse et syntomique ; chacun des sites
$\mathcal{X}_{Zar}$, $\mathcal{X}_\etale$, $\mathcal{X}_{smooth}$,
$\mathcal{X}_{syntomic}$ et $\mathcal{X}_{fppf}$ est donc un site annelé.
Cette construction est elle aussi fonctorielle.

\begin{lemma}
\label{lemma-functoriality-structure-sheaf}
Soit $f : \mathcal{X} \to \mathcal{Y}$ un $1$-morphisme de catégories
fibrées en groupoïdes au-dessus de $(\Sch/S)_{fppf}$. Soit
$\tau \in \{Zar, \etale, smooth, syntomic, fppf\}$.
Il existe une identification canonique
$f^{-1}\mathcal{O}_\mathcal{Y} = \mathcal{O}_\mathcal{X}$
qui fait de
$f : \Sh(\mathcal{X}_\tau) \to \Sh(\mathcal{Y}_\tau)$
un morphisme de topos annelés.
\end{lemma}

\begin{proof}
Notons $p : \mathcal{X} \to (\Sch/S)_{fppf}$ et
$q : \mathcal{Y} \to (\Sch/S)_{fppf}$ les foncteurs structuraux.
Alors $p = q \circ f$, donc $p^{-1} = f^{-1} \circ q^{-1}$ d'après le
Lemme \ref{lemma-1-morphisms-presheaves}. Puisque
$\mathcal{O}_\mathcal{X} = p^{-1}\mathcal{O}$ and
$\mathcal{O}_\mathcal{Y} = q^{-1}\mathcal{O}$
le résultat en découle.
\end{proof}

\begin{remark}
\label{remark-flat}
Dans la situation du
Lemme \ref{lemma-functoriality-structure-sheaf},
le morphisme de topos annelés
$f : \Sh(\mathcal{X}_\tau) \to \Sh(\mathcal{Y}_\tau)$
est plat, comme le montre l'égalité
$f^{-1}\mathcal{O}_\mathcal{X} = \mathcal{O}_\mathcal{Y}$.
Cela est quelque peu contre-intuitif, car, par exemple, une immersion
fermée de champs algébriques n'est généralement pas plate (comme morphisme de
champs algébriques).
Toutefois, il se produit exactement la même chose pour une immersion
fermée $i : X \to Y$ de schémas : dans ce cas, le morphisme associé
de grands sites pour $\tau$
$i : (\Sch/X)_\tau \to (\Sch/Y)_\tau$
est lui aussi plat.
\end{remark}




\section{Faisceaux de modules}
\label{section-modules}

\noindent
Puisque nous disposons d'un faisceau structural, nous disposons de modules.

\begin{definition}
\label{definition-modules}
Soit $\mathcal{X}$ une catégorie fibrée en groupoïdes au-dessus de
$(\Sch/S)_{fppf}$.
\begin{enumerate}
\item Un {\it préfaisceau de modules sur $\mathcal{X}$} est un
préfaisceau de $\mathcal{O}_\mathcal{X}$-modules. La catégorie des
préfaisceaux de modules est notée $\textit{PMod}(\mathcal{O}_\mathcal{X})$.
\item Nous disons qu'un préfaisceau de modules $\mathcal{F}$ est un
{\it $\mathcal{O}_\mathcal{X}$-module}, ou plus précisément un
{\it faisceau de $\mathcal{O}_\mathcal{X}$-modules}, si $\mathcal{F}$
est un faisceau fppf. La catégorie des $\mathcal{O}_\mathcal{X}$-modules
est notée $\textit{Mod}(\mathcal{O}_\mathcal{X})$.
\end{enumerate}
\end{definition}

\noindent
Dans la littérature, ces (pré)faisceaux de modules apparaissent comme des {\it (pré)faisceaux
de $\mathcal{O}_\mathcal{X}$-modules sur le grand site fppf de $\mathcal{X}$}.
Nous emploierons parfois cette terminologie pour distinguer ces
catégories d'autres catégories. Nous rencontrerons aussi des préfaisceaux de modules qui
sont des faisceaux pour les topologies de Zariski, étale, lisse ou syntomique
(sans être nécessairement des faisceaux fppf). Au besoin, nous les noterons
$\textit{Mod}(\mathcal{X}_\etale, \mathcal{O}_\mathcal{X})$
et de même pour les autres topologies.

\medskip\noindent
Étudions maintenant la fonctorialité, d'abord pour les préfaisceaux de modules. Soit
$$
\xymatrix{
\mathcal{X} \ar[rr]_f \ar[rd]_p & &  \mathcal{Y} \ar[ld]^q \\
& (\Sch/S)_{fppf}
}
$$
un $1$-morphisme de catégories fibrées en groupoïdes.
Les foncteurs $f^{-1}$ et $f_*$ sur les préfaisceaux abéliens s'étendent en des foncteurs
\begin{equation}
\label{equation-functoriality-presheaves-modules}
f^{-1} :
\textit{PMod}(\mathcal{O}_\mathcal{Y})
\longrightarrow
\textit{PMod}(\mathcal{O}_\mathcal{X})
\quad\text{et}\quad
f_* :
\textit{PMod}(\mathcal{O}_\mathcal{X})
\longrightarrow
\textit{PMod}(\mathcal{O}_\mathcal{Y})
\end{equation}
Cela est immédiat pour $f^{-1}$, car
$f^{-1}\mathcal{G}(x) = \mathcal{G}(f(x))$, qui est un module sur
$\mathcal{O}_\mathcal{Y}(f(x)) = \mathcal{O}(q(f(x))) = \mathcal{O}(p(x)) =
\mathcal{O}_\mathcal{X}(x)$. On peut aussi le déduire de ce que
$f^{-1}\mathcal{O}_\mathcal{Y} = \mathcal{O}_\mathcal{X}$
et de ce que $f^{-1}$ commute aux limites (sur les préfaisceaux).
Comme $f_*$ est un adjoint à droite, il commute à toutes les limites
(sur les préfaisceaux), en particulier aux produits. Nous pouvons donc étendre
$f_*$ en un foncteur sur les préfaisceaux de modules comme dans la démonstration de
Modules sur les sites, Lemme \ref{sites-modules-lemma-pushforward-module}.
Nous affirmons que les foncteurs (\ref{equation-functoriality-presheaves-modules})
forment une paire de foncteurs adjoints :
$$
\Mor_{\textit{PMod}(\mathcal{O}_\mathcal{X})}(
f^{-1}\mathcal{G}, \mathcal{F})
=
\Mor_{\textit{PMod}(\mathcal{O}_\mathcal{Y})}(
\mathcal{G}, f_*\mathcal{F}).
$$
Comme $f^{-1}\mathcal{O}_\mathcal{Y} = \mathcal{O}_\mathcal{X}$,
cela résulte de
Modules sur les sites, Lemme \ref{sites-modules-lemma-adjoint-push-pull-modules},
en munissant $\mathcal{X}$ et $\mathcal{Y}$ de la topologie
chaotique.

\medskip\noindent
Étudions ensuite la fonctorialité pour les modules, c'est-à-dire pour les faisceaux de modules
de la topologie fppf. Notons encore $f$ le morphisme induit de topos
annelés ; voir le
Lemme \ref{lemma-functoriality-structure-sheaf}
(pour l'instant appliqué aux topologies fppf). Remarquons que les foncteurs
$f^{-1}$ and $f_*$ of (\ref{equation-functoriality-presheaves-modules})
préservent les sous-catégories de faisceaux de modules ; voir le
Lemme \ref{lemma-functoriality-sheaves}.
Il s'ensuit immédiatement que
\begin{equation}
\label{equation-functoriality-sheaves-modules}
f^{-1} :
\textit{Mod}(\mathcal{O}_\mathcal{Y})
\longrightarrow
\textit{Mod}(\mathcal{O}_\mathcal{X})
\quad\text{et}\quad
f_* :
\textit{Mod}(\mathcal{O}_\mathcal{X})
\longrightarrow
\textit{Mod}(\mathcal{O}_\mathcal{Y})
\end{equation}
forment une paire de foncteurs adjoints :
$$
\Mor_{\textit{Mod}(\mathcal{O}_\mathcal{X})}(
f^{-1}\mathcal{G}, \mathcal{F})
=
\Mor_{\textit{Mod}(\mathcal{O}_\mathcal{Y})}(
\mathcal{G}, f_*\mathcal{F}).
$$
D'après l'unicité des adjoints, nous concluons que
$f^* = f^{-1}$, où $f^*$ est défini dans
Modules sur les sites, Section \ref{sites-modules-section-functoriality-modules},
pour le morphisme de topos annelés $f$ ci-dessus. Nous aurions bien entendu pu
le voir directement, car
$f^*(-) = f^{-1}(-) \otimes_{f^{-1}\mathcal{O}_\mathcal{Y}}
\mathcal{O}_\mathcal{X}$, et puisque
$f^{-1}\mathcal{O}_\mathcal{Y} = \mathcal{O}_\mathcal{X}$.

\medskip\noindent
Il en va de même pour les faisceaux de modules des topologies de Zariski, étale, lisse et
syntomique.



\section{Catégories représentables}
\label{section-representable}

\noindent
Dans cette courte section, nous comparons nos définitions à la situation
où les champs algébriques considérés sont représentables.

\begin{lemma}
\label{lemma-compare-with-scheme}
Soit $S$ un schéma. Soit $\mathcal{X}$ une catégorie fibrée
en groupoïdes au-dessus de $(\Sch/S)$. Supposons que $\mathcal{X}$ soit représentable
par un schéma $X$. Pour $\tau \in \{Zar,\linebreak[0] \etale,\linebreak[0]
smooth,\linebreak[0] syntomic,\linebreak[0] fppf\}$
il existe une équivalence canonique
$$
(\mathcal{X}_\tau, \mathcal{O}_\mathcal{X}) =
((\Sch/X)_\tau, \mathcal{O}_X)
$$
de sites annelés.
\end{lemma}

\begin{proof}
Il suffit de choisir une équivalence
$(\Sch/X)_\tau \to \mathcal{X}$ de catégories fibrées en groupoïdes
au-dessus de $(\Sch/S)_{fppf}$ et d'utiliser la fonctorialité de
la construction $\mathcal{X} \leadsto \mathcal{X}_\tau$.
\end{proof}

\begin{lemma}
\label{lemma-compare-with-morphism-of-schemes}
Soit $S$ un schéma. Soit $f : \mathcal{X} \to \mathcal{Y}$ un morphisme
de catégories fibrées en groupoïdes au-dessus de $S$.
Supposons que $\mathcal{X}$ et $\mathcal{Y}$ soient représentables par des schémas
$X$ et $Y$. Soit $f : X \to Y$ le morphisme de schémas correspondant
à $f$. Pour $\tau \in \{Zar,\linebreak[0] \etale,\linebreak[0]
smooth,\linebreak[0] syntomic,\linebreak[0] fppf\}$
le morphisme de topos annelés
$f : (\Sh(\mathcal{X}_\tau), \mathcal{O}_\mathcal{X}) \to
(\Sh(\mathcal{Y}_\tau), \mathcal{O}_\mathcal{Y})$
coïncide avec le morphisme de topos annelés
$f : (\Sh((\Sch/X)_\tau), \mathcal{O}_X) \to 
(\Sh((\Sch/Y)_\tau), \mathcal{O}_Y)$ au moyen des identifications du
Lemme \ref{lemma-compare-with-scheme}.
\end{lemma}

\begin{proof}
Il suffit de dérouler les définitions.
\end{proof}




\section{Restriction}
\label{section-restriction}


\noindent
Une observation triviale mais utile est que la localisation
d'une catégorie fibrée en groupoïdes en un objet
est équivalente au grand site du schéma au-dessus duquel il se trouve.

\begin{lemma}
\label{lemma-localizing}
Soit $p : \mathcal{X} \to (\Sch/S)_{fppf}$ une catégorie fibrée
en groupoïdes. Soit $\tau \in \{Zar, \etale, smooth, syntomic, fppf\}$.
Soit $x \in \Ob(\mathcal{X})$ au-dessus de $U = p(x)$.
Le foncteur $p$ induit une équivalence de sites
$\mathcal{X}_\tau/x \to (\Sch/U)_\tau$.
\end{lemma}

\begin{proof}
Cas particulier de Champs, Lemme \ref{stacks-lemma-localizing}.
\end{proof}

\noindent
Nous utilisons le lemme précédent pour définir l'image réciproque et la restriction
d'un (pré)faisceau à un schéma.

\begin{definition}
\label{definition-pullback}
Soit $p : \mathcal{X} \to (\Sch/S)_{fppf}$ une catégorie fibrée
en groupoïdes. Soit $x \in \Ob(\mathcal{X})$ au-dessus de $U = p(x)$.
Soit $\mathcal{F}$ un préfaisceau sur $\mathcal{X}$.
\begin{enumerate}
\item L'{\it image réciproque $x^{-1}\mathcal{F}$ de $\mathcal{F}$} est la
restriction $\mathcal{F}|_{(\mathcal{X}/x)}$, considérée comme un préfaisceau sur
$(\Sch/U)_{fppf}$ au moyen de l'équivalence
$\mathcal{X}/x \to (\Sch/U)_{fppf}$ du
Lemme \ref{lemma-localizing}.
\item La {\it restriction de $\mathcal{F}$ à $U_\etale$}
est $x^{-1}\mathcal{F}|_{U_\etale}$, notée abusivement
$\mathcal{F}|_{U_\etale}$.
\end{enumerate}
\end{definition}

\noindent
Cette notation a un sens, car le lemme de $2$-Yoneda associe à l'objet $x$, voir
Champs algébriques, Section \ref{algebraic-section-2-yoneda},
un $1$-morphisme $x : (\Sch/U)_{fppf} \to \mathcal{X}/x$
quasi-inverse de $p : \mathcal{X}/x \to (\Sch/U)_{fppf}$.
Ainsi, $x^{-1}\mathcal{F}$ est bien l'image réciproque de $\mathcal{F}$ par ce
$1$-morphisme. En particulier, d'après ce qui précède, si $\mathcal{F}$
est un faisceau (ou un faisceau de Zariski, étale, lisse ou syntomique), alors
$x^{-1}\mathcal{F}$ est un faisceau sur $(\Sch/U)_{fppf}$ (ou sur
$(\Sch/U)_{Zar}$, $(\Sch/U)_\etale$,
$(\Sch/U)_{smooth}$, $(\Sch/U)_{syntomic}$).

\medskip\noindent
Soit $p : \mathcal{X} \to (\Sch/S)_{fppf}$ une catégorie fibrée
en groupoïdes. Soit $\varphi : x \to y$ un morphisme de $\mathcal{X}$
au-dessus du morphisme de schémas $a : U \to V$.
Rappelons que $a$ induit un morphisme de petits sites
étales $a_{small} : U_\etale \to V_\etale$ ; voir
Cohomologie étale, Section \ref{etale-cohomology-section-functoriality}.
Soit $\mathcal{F}$ un préfaisceau sur $\mathcal{X}$.
Soient $\mathcal{F}|_{U_\etale}$ et
$\mathcal{F}|_{V_\etale}$ les restrictions de $\mathcal{F}$
par $x$ et $y$. Il existe une application naturelle de {\it comparaison}
\begin{equation}
\label{equation-comparison-push}
c_\varphi :
\mathcal{F}|_{V_\etale}
\longrightarrow
a_{small, *}(\mathcal{F}|_{U_\etale})
\end{equation}
de préfaisceaux sur $U_\etale$. En effet, si $V' \to V$ est étale,
posons $U' = V' \times_V U$ et définissons $c_\varphi$ sur les sections au-dessus de $V'$
par
$$
\xymatrix{
a_{small, *}(\mathcal{F}|_{U_\etale})(V') &
\mathcal{F}|_{U_\etale}(U') \ar@{=}[l] &
\mathcal{F}(x') \ar@{=}[l] \\
\mathcal{F}|_{V_\etale}(V') \ar@{=}[rr] \ar[u]^{c_\varphi} &
&
\mathcal{F}(y') \ar[u]_{\mathcal{F}(\varphi')}
}
$$
Here $\varphi' : x' \to y'$ is a morphism of $\mathcal{X}$
fitting into a commutative diagram
$$
\vcenter{
\xymatrix{
x' \ar[r] \ar[d]_{\varphi'} & x \ar[d]^\varphi \\
y' \ar[r] & y
}
}
\quad\text{au-dessus de}\quad
\vcenter{
\xymatrix{
U' \ar[r] \ar[d] & U \ar[d]^a \\
V' \ar[r] & V
}
}
$$
L'existence et l'unicité de $\varphi'$ résultent des axiomes
d'une catégorie fibrée en groupoïdes.
Nous omettons de vérifier que $c_\varphi$ ainsi défini est bien un morphisme
de préfaisceaux (c'est-à-dire compatible avec les applications de restriction) et qu'il
est fonctoriel en $\mathcal{F}$. Si $\mathcal{F}$ est un faisceau pour la
topologie étale, nous obtenons une application de {\it comparaison}
\begin{equation}
\label{equation-comparison}
c_\varphi : a_{small}^{-1}(\mathcal{F}|_{V_\etale})
\longrightarrow
\mathcal{F}|_{U_\etale}
\end{equation}
que l'on note également $c_\varphi$, comme indiqué (c'est l'abus de notation usuel
qui consiste à ne pas distinguer des applications adjointes).

\begin{lemma}
\label{lemma-comparison}
Soit $\mathcal{F}$ un faisceau étale sur $\mathcal{X} \to (\Sch/S)_{fppf}$.
\begin{enumerate}
\item Si $\varphi : x \to y$ et $\psi : y \to z$
sont des morphismes de $\mathcal{X}$ au-dessus de $a : U \to V$ et
$b : V \to W$, alors le composé
$$
a_{small}^{-1}(b_{small}^{-1} (\mathcal{F}|_{W_\etale}))
\xrightarrow{a_{small}^{-1}c_\psi}
a_{small}^{-1}(\mathcal{F}|_{V_\etale})
\xrightarrow{c_\varphi}
\mathcal{F}|_{U_\etale}
$$
est égal à $c_{\psi \circ \varphi}$ au moyen de l'identification
$$
(b \circ a)_{small}^{-1}(\mathcal{F}|_{W_\etale}) =
a_{small}^{-1}(b_{small}^{-1} (\mathcal{F}|_{W_\etale})).
$$
\item Si $\varphi : x \to y$ est au-dessus d'un morphisme étale de schémas
$a : U \to V$, alors (\ref{equation-comparison}) est un isomorphisme.
\item Supposons que $f : \mathcal{Y} \to \mathcal{X}$ soit un $1$-morphisme de
catégories fibrées en groupoïdes au-dessus de $(\Sch/S)_{fppf}$, et que $y$ soit
un objet de $\mathcal{Y}$ au-dessus du schéma $U$, d'image
$x = f(y)$. Il existe alors une identification canonique
$f^{-1}\mathcal{F}|_{U_\etale} = \mathcal{F}|_{U_\etale}$.
\item De plus, étant donné $\psi : y' \to y$ dans $\mathcal{Y}$ au-dessus de
$a : U' \to U$, l'application de comparaison
$c_\psi : a_{small}^{-1}(f^{-1}\mathcal{F}|_{U_\etale}) \to
f^{-1}\mathcal{F}|_{U'_\etale}$ est égale à
l'application de comparaison $c_{f(\psi)} : a_{small}^{-1}\mathcal{F}|_{U_\etale}
\to \mathcal{F}|_{U'_\etale}$ au moyen des identifications de (3).
\end{enumerate}
\end{lemma}

\begin{proof}
La vérification de ces propriétés est omise.
\end{proof}

\noindent
Passons à la restriction des (pré)faisceaux de modules.

\begin{lemma}
\label{lemma-localizing-structure-sheaf}
Soit $p : \mathcal{X} \to (\Sch/S)_{fppf}$ une catégorie fibrée
en groupoïdes. Soit $\tau \in \{Zar, \etale, smooth, syntomic, fppf\}$.
Soit $x \in \Ob(\mathcal{X})$ au-dessus de $U = p(x)$.
L'équivalence du
Lemme \ref{lemma-localizing}
s'étend en une équivalence de sites annelés
$(\mathcal{X}_\tau/x, \mathcal{O}_\mathcal{X}|_x) \to
((\Sch/U)_\tau, \mathcal{O})$.
\end{lemma}

\begin{proof}
Cela résulte immédiatement de la construction des faisceaux structuraux.
\end{proof}

\noindent
Soit $\mathcal{X}$ une catégorie fibrée en groupoïdes au-dessus de $(\Sch/S)_{fppf}$.
Soit $\mathcal{F}$ un (pré)faisceau de modules sur $\mathcal{X}$ comme dans la
Définition \ref{definition-modules}.
Soit $x$ un objet de $\mathcal{X}$ au-dessus de $U$. Alors le
Lemme \ref{lemma-localizing-structure-sheaf}
garantit que la restriction
$x^{-1}\mathcal{F}$ est un (pré)faisceau de modules sur $(\Sch/U)_{fppf}$.
Dans ce cas, nous écrirons parfois $x^*\mathcal{F} = x^{-1}\mathcal{F}$.
De même, si $\mathcal{F}$ est un faisceau pour la topologie de Zariski, étale, lisse
ou syntomique, alors $x^{-1}\mathcal{F}$ en est un aussi. De plus, la
restriction
$\mathcal{F}|_{U_\etale} = x^{-1}\mathcal{F}|_{U_\etale}$
à $U$ est un préfaisceau de $\mathcal{O}_{U_\etale}$-modules.
Si $\mathcal{F}$ est un faisceau pour la topologie étale, alors
$\mathcal{F}|_{U_\etale}$ est un faisceau de modules. De plus,
si $\varphi : x \to y$ est un morphisme de $\mathcal{X}$ au-dessus de
$a : U \to V$, alors l'application de comparaison (\ref{equation-comparison})
est compatible à $a_{small}^\sharp$ (voir
Descente, Remarque \ref{descent-remark-change-topologies-ringed})
et induit une application de {\it comparaison}
\begin{equation}
\label{equation-comparison-modules}
c_\varphi : a_{small}^*(\mathcal{F}|_{V_\etale})
\longrightarrow
\mathcal{F}|_{U_\etale}
\end{equation}
de $\mathcal{O}_{U_\etale}$-modules.
Remarquons que les propriétés (1), (2), (3) et (4) du
Lemme \ref{lemma-comparison}
valent également pour les faisceaux étales de modules.
Nous l'utiliserons dans la suite sans autre mention.

\begin{lemma}
\label{lemma-enough-points}
Soit $p : \mathcal{X} \to (\Sch/S)_{fppf}$ une catégorie fibrée
en groupoïdes. Soit $\tau \in \{Zar, \etale, smooth, syntomic, fppf\}$.
Le site $\mathcal{X}_\tau$ a suffisamment de points.
\end{lemma}

\begin{proof}
D'après
Sites, Lemme \ref{sites-lemma-enough-points-local},
il faut montrer qu'il existe une famille d'objets $x$ de $\mathcal{X}$
telle que $\mathcal{X}_\tau/x$ ait suffisamment de points et que les faisceaux
$h_x^\#$ recouvrent l'objet final de la catégorie des faisceaux.
D'après le
Lemme \ref{lemma-localizing}
et
Cohomologie étale, Lemme \ref{etale-cohomology-lemma-points-fppf},
$\mathcal{X}_\tau/x$ a suffisamment de points pour tout objet
$x$, ce qui conclut.
\end{proof}







\section{Restriction aux espaces algébriques}
\label{section-restriction-algebraic-spaces}

\noindent
Dans cette section, nous considérons les faisceaux sur des catégories représentables par des
espaces algébriques. Le lemme suivant est l'analogue de
Topologies, Lemma \ref{topologies-lemma-at-the-bottom-etale}
pour les espaces algébriques.

\begin{lemma}
\label{lemma-compare}
Soit $S$ un schéma. Soit $\mathcal{X} \to (\Sch/S)_{fppf}$ une catégorie
fibrée en groupoïdes. Supposons que $\mathcal{X}$ soit représentable par un espace
algébrique $F$. Il existe alors un foncteur continu et cocontinu
$
F_\etale \to \mathcal{X}_\etale
$
qui induit un morphisme de sites annelés
$$
\pi_F :
(\mathcal{X}_\etale, \mathcal{O}_\mathcal{X})
\longrightarrow
(F_\etale, \mathcal{O}_F)
$$
et un morphisme de topos annelés
$$
i_F :
(\Sh(F_\etale), \mathcal{O}_F)
\longrightarrow
(\Sh(\mathcal{X}_\etale), \mathcal{O}_\mathcal{X})
$$
tel que $\pi_F \circ i_F = \text{id}$. De plus, $\pi_{F, *} = i_F^{-1}$.
\end{lemma}

\begin{proof}
Choisissons une équivalence $j : \mathcal{S}_F \to \mathcal{X}$ ; voir
Champs algébriques, Sections \ref{algebraic-section-split} et
\ref{algebraic-section-representable-by-algebraic-spaces}.
Un objet de $F_\etale$ est un schéma $U$ muni d'un
morphisme étale $\varphi : U \to F$. Alors $\varphi$ est un objet
de $\mathcal{S}_F$ au-dessus de $U$. Ainsi $j(\varphi)$ est un objet de
$\mathcal{X}$ au-dessus de $U$. De cette façon, $j$ induit un foncteur
$u : F_\etale \to \mathcal{X}$. Il est clair que
$u$ est continu et cocontinu pour la topologie étale sur
$\mathcal{X}$. Puisque $j$ est une équivalence, le foncteur $u$ est pleinement
fidèle. En outre, les produits fibrés et les égalisateurs existent dans $F_\etale$
et $u$ y commute, car ils se calculent au niveau
des schémas sous-jacents dans $F_\etale$. Ainsi,
Sites, Lemmas \ref{sites-lemma-when-shriek},
\ref{sites-lemma-preserve-equalizers} et
\ref{sites-lemma-back-and-forth}
s'appliquent. En particulier, $u$ définit un morphisme de topos
$i_F : \Sh(F_\etale) \to \Sh(\mathcal{X}_\etale)$
et il existe un adjoint à gauche $i_{F, !}$ de $i_F^{-1}$ qui commute
aux produits fibrés et aux égalisateurs.

\medskip\noindent
Nous affirmons que $i_{F, !}$ est exact. Si tel est le cas, nous pouvons définir
$\pi_F$ par les règles $\pi_F^{-1} = i_{F, !}$ et $\pi_{F, *} = i_F^{-1}$,
et tout est clair. Pour démontrer l'assertion, remarquons que nous savons déjà
que $i_{F, !}$
est exact à droite et préserve les produits fibrés. Il suffit donc de montrer
que $i_{F, !}* = *$, où $*$ désigne l'objet final de la catégorie
des faisceaux d'ensembles. Soit $U$ un schéma et soit
$\varphi : U \to F$ surjectif et étale. Posons $R = U \times_F U$.
Alors
$$
\xymatrix{
h_R \ar@<1ex>[r] \ar@<-1ex>[r] & h_U \ar[r] & {*}
}
$$
est un diagramme coégalisateur dans $\Sh(F_\etale)$. En utilisant
l'exactitude à droite de $i_{F, !}$, l'égalité $i_{F, !} = (u_p\ )^\#$ et le
Sites, Lemma \ref{sites-lemma-pullback-representable-presheaf}
nous voyons que
$$
\xymatrix{
h_{u(R)} \ar@<1ex>[r] \ar@<-1ex>[r] & h_{u(U)} \ar[r] & i_{F, !}{*}
}
$$
est un diagramme coégalisateur dans $\Sh(\mathcal{X}_\etale)$. Puisque
$j$ est une équivalence et que $F = U/R$, il s'ensuit que
le coégalisateur dans $\Sh(\mathcal{X}_\etale)$ des
deux applications $h_{u(R)} \to h_{u(U)}$ est $*$. Nous omettons de démontrer que
ces morphismes sont compatibles aux faisceaux structuraux.
\end{proof}

\begin{remark}
\label{remark-compare}
Les constructions du Lemme \ref{lemma-compare} sont compatibles à la localisation
étale. En voici une formulation précise. Soit $S$ un schéma.
Soit $f : \mathcal{X} \to \mathcal{Y}$ un morphisme
de catégories fibrées en groupoïdes au-dessus de $(\Sch/S)_{fppf}$. Supposons que
$\mathcal{X}$ et $\mathcal{Y}$ soient représentables par des espaces algébriques $F$ et $G$,
et que le morphisme induit $f : F \to G$ d'espaces algébriques soit étale.
Notons $f_{small} : F_\etale \to G_\etale$
le morphisme correspondant de topos annelés. Alors
$$
\xymatrix{
(\Sh(F_\etale), \mathcal{O}_F) \ar[rr]_{f_{small}} \ar[d]_{i_F} & &
(\Sh(G_\etale), \mathcal{O}_G) \ar[d]^{i_G} \\
(\Sh(\mathcal{X}_\etale), \mathcal{O}_\mathcal{X})
\ar[d]_{\pi_F} \ar[rr]_f & &
(\Sh(\mathcal{Y}_\etale), \mathcal{O}_\mathcal{Y}) \ar[d]^{\pi_G} \\
(\Sh(F_\etale), \mathcal{O}_F) \ar[rr]^{f_{small}} & &
(\Sh(G_\etale), \mathcal{O}_G)
}
$$
est un diagramme commutatif de topos annelés. Nous omettons les détails.
\end{remark}

\noindent
Supposons que $\mathcal{X}$ soit un champ algébrique représenté par
l'espace algébrique $F$.
Soit $j : \mathcal{S}_F \to \mathcal{X}$ une équivalence, et notons
$u : F_\etale \to \mathcal{X}_\etale$ le
foncteur de la démonstration du Lemme \ref{lemma-compare} ci-dessus.
Pour tout faisceau $\mathcal{F}$ sur $\mathcal{X}_\etale$, nous avons
$$
\pi_{F, *}\mathcal{F}(U) = i_F^{-1}\mathcal{F}(U) = \mathcal{F}(u(U)).
$$
C'est pourquoi nous considérons souvent $i_F^{-1}$ comme un {\it foncteur de restriction},
par analogie avec la
Définition \ref{definition-pullback}
et avec la restriction d'un faisceau du grand site étale d'un
schéma au petit site étale de ce schéma. Nous employons souvent la notation
\begin{equation}
\label{equation-restrict}
\mathcal{F}|_{F_\etale} = i_F^{-1}\mathcal{F} = \pi_{F, *}\mathcal{F}
\end{equation}
dans cette situation.

\begin{lemma}
\label{lemma-compare-morphism}
Soit $S$ un schéma. Soit $f : \mathcal{X} \to \mathcal{Y}$ un morphisme
de catégories fibrées en groupoïdes au-dessus de $(\Sch/S)_{fppf}$. Supposons que
$\mathcal{X}$ et $\mathcal{Y}$ soient représentables par des espaces algébriques $F$ et $G$.
Notons $f : F \to G$ le morphisme induit d'espaces algébriques, et
$f_{small} : F_\etale \to G_\etale$
le morphisme correspondant de topos annelés. Alors
$$
\xymatrix{
(\Sh(\mathcal{X}_\etale), \mathcal{O}_\mathcal{X})
\ar[d]_{\pi_F} \ar[rr]_f & &
(\Sh(\mathcal{Y}_\etale), \mathcal{O}_\mathcal{Y}) \ar[d]^{\pi_G} \\
(\Sh(F_\etale), \mathcal{O}_F) \ar[rr]^{f_{small}} & &
(\Sh(G_\etale), \mathcal{O}_G)
}
$$
est un diagramme commutatif de topos annelés.
\end{lemma}

\begin{proof}
Cela ressemble à
Topologies, Lemme \ref{topologies-lemma-morphism-big-small-etale} (3),
mais une petite difficulté vient de ce que $F \to G$ n'est pas nécessairement
représentable par des schémas. En particulier, nous n'obtenons pas un diagramme commutatif
de sites annelés, mais seulement un diagramme commutatif de topos annelés.

\medskip\noindent
Avant d'aborder la démonstration proprement dite, choisissons des équivalences
$j : \mathcal{S}_F \to \mathcal{X}$ et
$j' : \mathcal{S}_G \to \mathcal{Y}$ qui induisent des foncteurs
$u : F_\etale \to \mathcal{X}$ et
$u' : G_\etale \to \mathcal{Y}$ comme dans la démonstration du
Lemme \ref{lemma-compare}. Grâce à la $2$-fonctorialité des
faisceaux sur les catégories fibrées en groupoïdes au-dessus de $\Sch_{fppf}$
(voir la discussion de la Section \ref{section-presheaves}),
nous pouvons supposer que $\mathcal{X} = \mathcal{S}_F$ et
$\mathcal{Y} = \mathcal{S}_G$, et que $f : \mathcal{S}_F \to \mathcal{S}_G$
est le foncteur associé au morphisme $f : F \to G$. Par conséquent,
nous omettrons $u$ et $u'$ de la notation : étant donné un objet
$U \to F$ de $F_\etale$, nous notons $U/F$
l'objet correspondant de $\mathcal{X}$. De même pour $G$.

\medskip\noindent
Soit $\mathcal{G}$ un faisceau sur $\mathcal{X}_\etale$.
Pour démontrer (2), calculons $\pi_{G, *}f_*\mathcal{G}$ et
$f_{small, *}\pi_{F, *}\mathcal{G}$. Pour cela, soit $V \to G$ un objet
de $G_\etale$. Alors
$$
\pi_{G, *}f_*\mathcal{G}(V) = f_*\mathcal{G}(V/G) =
\Gamma\Big(
(\Sch/V)_{fppf} \times_{\mathcal{Y}} \mathcal{X},
\ \text{pr}^{-1}\mathcal{G}\Big)
$$
voir (\ref{equation-pushforward}). Le produit fibré figurant dans la formule est
$$
(\Sch/V)_{fppf} \times_{\mathcal{Y}} \mathcal{X} =
(\Sch/V)_{fppf} \times_{\mathcal{S}_G} \mathcal{S}_F =
\mathcal{S}_{V \times_G F}
$$
c'est-à-dire la catégorie scindée fibrée en groupoïdes associée à
l'espace algébrique $V \times_G F$. Et $\text{pr}^{-1}\mathcal{G}$ est un
faisceau sur $\mathcal{S}_{V \times_G F}$ pour la topologie étale.

\medskip\noindent
En particulier, si $V \times_G F$ est représentable, c'est-à-dire s'il est un schéma,
alors $\pi_{G, *}f_*\mathcal{G}(V) = \mathcal{G}(V \times_G F/F)$, et
également
$$
f_{small, *}\pi_{F, *}\mathcal{G}(V) =
\pi_{F, *}\mathcal{G}(V \times_G F) =
\mathcal{G}(V \times_G F/F)
$$
ce qui démontre l'égalité voulue dans ce cas particulier.

\medskip\noindent
En général, choisissons un schéma $U$ et un morphisme étale surjectif
$U \to V \times_G F$. Posons $R = U \times_{V \times_G F} U$. Alors
$U/V \times_G F$ et $R/V \times_G F$ sont des objets de la catégorie
produit fibré ci-dessus. Puisque $\text{pr}^{-1}\mathcal{G}$ est un
faisceau pour la topologie étale sur $\mathcal{S}_{V \times_G F}$,
le diagramme
$$
\xymatrix{
\Gamma\Big(
(\Sch/V)_{fppf} \times_{\mathcal{Y}} \mathcal{X},
\ \text{pr}^{-1}\mathcal{G}\Big)
\ar[r] &
\text{pr}^{-1}\mathcal{G}(U/V \times_G F) \ar@<1ex>[r] \ar@<-1ex>[r] &
\text{pr}^{-1}\mathcal{G}(R/V \times_G F)
}
$$
est un diagramme égalisateur. Remarquons que
$\text{pr}^{-1}\mathcal{G}(U/V \times_G F) = \mathcal{G}(U/F)$
and $\text{pr}^{-1}\mathcal{G}(R/V \times_G F) = \mathcal{G}(R/F)$
par définition des images réciproques. De plus, d'après les résultats de
Propriétés des espaces, Section \ref{spaces-properties-section-etale-site}
(en particulier,
Propriétés des espaces,
Remarque \ref{spaces-properties-remark-explain-equivalence} et
Lemme \ref{spaces-properties-lemma-functoriality-etale-site}),
nous voyons qu'il existe un diagramme égalisateur
$$
\xymatrix{
f_{small, *}\pi_{F, *}\mathcal{G}(V)
\ar[r] &
\pi_{F, *}\mathcal{G}(U/F) \ar@<1ex>[r] \ar@<-1ex>[r] &
\pi_{F, *}\mathcal{G}(R/F)
}
$$
Comme nous avons aussi $\pi_{F, *}\mathcal{G}(U/F) = \mathcal{G}(U/F)$
et $\pi_{F, *}\mathcal{G}(U/F) = \mathcal{G}(U/F)$,
nous obtenons une identification canonique
$f_{small, *}\pi_{F, *}\mathcal{G}(V) = \pi_{G, *}f_*\mathcal{G}(V)$.
Nous omettons de démontrer la compatibilité avec les applications de restriction
et la fonctorialité en $\mathcal{G}$.
\end{proof}

\noindent
Soient $f : \mathcal{X} \to \mathcal{Y}$ et $f : F \to G$ comme dans la
seconde partie du lemme précédent. Le lemme et
(\ref{equation-restrict}) donnent
\begin{equation}
\label{equation-compare-big-small}
(f_*\mathcal{F})|_{G_\etale} =
f_{small, *}(\mathcal{F}|_{F_\etale})
\end{equation}
pour tout faisceau $\mathcal{F}$ sur $\mathcal{X}_\etale$.
De plus, si $\mathcal{F}$ est un faisceau de $\mathcal{O}$-modules, alors
(\ref{equation-compare-big-small}) est un isomorphisme de
$\mathcal{O}_G$-modules sur $G_\etale$.

\medskip\noindent
Enfin, supposons donné un diagramme $2$-commutatif
$$
\xymatrix{
\mathcal{U} \ar[r]^a \ar[dr]_f \drtwocell<\omit>{<-2>\varphi} &
\mathcal{V} \ar[d]^g \\
& \mathcal{X}
}
$$
de $1$-morphismes de catégories fibrées en groupoïdes au-dessus de $(\Sch/S)_{fppf}$,
que $\mathcal{F}$ soit un faisceau sur $\mathcal{X}_\etale$,
et que $\mathcal{U}$ et $\mathcal{V}$ soient représentables par les espaces algébriques
$U$ et $V$. Nous obtenons alors une application de comparaison
\begin{equation}
\label{equation-comparison-algebraic-spaces}
c_\varphi : a_{small}^{-1}(g^{-1}\mathcal{F}|_{V_\etale})
\longrightarrow
f^{-1}\mathcal{F}|_{U_\etale}
\end{equation}
où $a : U \to V$ désigne le morphisme d'espaces algébriques correspondant
à $a$. C'est l'analogue de (\ref{equation-comparison}). Nous définissons
$c_\varphi$ comme l'adjointe de l'application
$$
g^{-1}\mathcal{F}|_{V_\etale}
\longrightarrow
a_{small, *}(f^{-1}\mathcal{F}|_{U_\etale}) =
(a_*f^{-1}\mathcal{F})|_{V_\etale}
$$
(égalité donnée par (\ref{equation-compare-big-small}))
qui est la restriction à $V$, au sens de (\ref{equation-restrict}), de l'application
$$
g^{-1}\mathcal{F} \to a_*a^{-1}g^{-1}\mathcal{F} = a_*f^{-1}\mathcal{F}
$$
où la dernière égalité utilise la $2$-commutativité du diagramme ci-dessus.
Si $\mathcal{F}$ est un faisceau de $\mathcal{O}_\mathcal{X}$-modules,
$c_\varphi$ induit une application de {\it comparaison}
\begin{equation}
\label{equation-comparison-algebraic-spaces-modules}
c_\varphi : a_{small}^*(g^*\mathcal{F}|_{V_\etale})
\longrightarrow
f^*\mathcal{F}|_{U_\etale}
\end{equation}
de $\mathcal{O}_{U_\etale}$-modules. C'est l'analogue de
(\ref{equation-comparison-modules}).
Remarquons que les propriétés (1), (2), (3) et (4) du
Lemme \ref{lemma-comparison}
valent aussi dans ce cadre.












\section{Modules quasi-cohérents}
\label{section-quasi-coherent}

\noindent
Nous pouvons maintenant appliquer la définition générale d'un module
quasi-cohérent à la situation étudiée dans ce chapitre.

\begin{definition}
\label{definition-quasi-coherent}
Soit $p : \mathcal{X} \to (\Sch/S)_{fppf}$ une catégorie fibrée
en groupoïdes. Un {\it module quasi-cohérent sur $\mathcal{X}$}, ou un
{\it $\mathcal{O}_\mathcal{X}$-module quasi-cohérent}, est un
module quasi-cohérent sur le site annelé
$(\mathcal{X}_{fppf}, \mathcal{O}_\mathcal{X})$ au sens de
Modules sur les sites, Définition \ref{sites-modules-definition-site-local}.
La catégorie des faisceaux quasi-cohérents sur $\mathcal{X}$
est notée $\QCoh(\mathcal{O}_\mathcal{X})$.
\end{definition}

\noindent
Si $\mathcal{X}$ est un champ algébrique, cette définition coïncide avec toutes les
définitions de la littérature, au sens où $\QCoh(\mathcal{O}_\mathcal{X})$
est équivalente (aux questions ensemblistes près) à toute variante de cette catégorie
définie dans la littérature. Par exemple, nous comparerons notre définition à
celle de \cite[Définition 6.1]{olsson_sheaves} dans
Cohomologie des champs, Lemme \ref{lemma-quasi-coherent}.
Nous verrons également plus loin d'autres constructions de cette catégorie.

\medskip\noindent
En général (comme pour les morphismes de schémas), l'image directe
d'un faisceau quasi-cohérent par un $1$-morphisme n'est pas quasi-cohérente.
L'image réciproque, en revanche, préserve la quasi-cohérence.

\begin{lemma}
\label{lemma-pullback-quasi-coherent}
Soit $f : \mathcal{X} \to \mathcal{Y}$ un $1$-morphisme de catégories
fibrées en groupoïdes au-dessus de $(\Sch/S)_{fppf}$.
Le foncteur d'image réciproque
$f^* = f^{-1} : \textit{Mod}(\mathcal{O}_\mathcal{Y}) \to
\textit{Mod}(\mathcal{O}_\mathcal{X})$
préserve les faisceaux quasi-cohérents.
\end{lemma}

\begin{proof}
C'est un fait général ; voir
Modules sur les sites, Lemme \ref{sites-modules-lemma-local-pullback}.
\end{proof}

\noindent
Les faisceaux quasi-cohérents admettent une caractérisation très simple
en termes de leurs images réciproques. Voir aussi le
Lemme \ref{lemma-quasi-coherent}
pour une caractérisation en termes de restrictions.

\begin{lemma}
\label{lemma-characterize-quasi-coherent}
Soit $p : \mathcal{X} \to (\Sch/S)_{fppf}$ une catégorie
fibrée en groupoïdes. Soit $\mathcal{F}$
un faisceau de $\mathcal{O}_\mathcal{X}$-modules. Alors $\mathcal{F}$
est quasi-cohérent si et seulement si $x^*\mathcal{F}$ est un faisceau
quasi-cohérent sur $(\Sch/U)_{fppf}$ pour tout objet $x$ de
$\mathcal{X}$, où $U = p(x)$.
\end{lemma}

\begin{proof}
D'après le
Lemme \ref{lemma-pullback-quasi-coherent},
la condition est nécessaire. Réciproquement, puisque $x^*\mathcal{F}$
n'est autre que la restriction à $\mathcal{X}_{fppf}/x$, la condition est
suffisante par la définition même d'un faisceau quasi-cohérent
(et par le fait que la quasi-cohérence est une propriété intrinsèque
des faisceaux de modules ; voir
Modules sur les sites, Section \ref{sites-modules-section-intrinsic}).
\end{proof}

\begin{lemma}
\label{lemma-characterize-quasi-coherent-bis}
Soit $p : \mathcal{X} \to (\Sch/S)_{fppf}$ une catégorie
fibrée en groupoïdes. Soit $\mathcal{F}$ un préfaisceau de
modules sur $\mathcal{X}$. Les conditions suivantes sont équivalentes :
\begin{enumerate}
\item $\mathcal{F}$ est un objet de
$\textit{Mod}(\mathcal{X}_{Zar}, \mathcal{O}_\mathcal{X})$
et $\mathcal{F}$ est un module quasi-cohérent sur
$(\mathcal{X}_{Zar}, \mathcal{O}_\mathcal{X})$ au sens de
Modules sur les sites, Définition \ref{sites-modules-definition-site-local},
\item $\mathcal{F}$ est un objet de
$\textit{Mod}(\mathcal{X}_\etale, \mathcal{O}_\mathcal{X})$
et $\mathcal{F}$ est un module quasi-cohérent sur
$(\mathcal{X}_\etale, \mathcal{O}_\mathcal{X})$ au sens de
Modules sur les sites, Définition \ref{sites-modules-definition-site-local}, et
\item $\mathcal{F}$ est un module quasi-cohérent sur $\mathcal{X}$
au sens de la Définition \ref{definition-quasi-coherent}.
\end{enumerate}
\end{lemma}

\begin{proof}
Supposons que l'une des conditions (1), (2) ou (3) soit satisfaite.
Soit $x$ un objet de $\mathcal{X}$ au-dessus du schéma $U$.
Rappelons que $x^*\mathcal{F} = x^{-1}\mathcal{F}$ n'est autre que la
restriction à $\mathcal{X}/x = (\Sch/U)_\tau$, où
$\tau = fppf$, $\tau = \etale$ ou $\tau = Zar$ ; voir la
Section \ref{section-restriction}.
D'après la définition des modules quasi-cohérents sur un site annelé,
cette restriction est quasi-cohérente dès que $\mathcal{F}$ l'est.
D'après Descente, Proposition \ref{descent-proposition-equivalence-quasi-coherent},
$x^*\mathcal{F}$ est le faisceau associé à
un $\mathcal{O}_U$-module quasi-cohérent ; c'est donc
un module quasi-cohérent pour les topologies fppf, étale et de
Zariski ; nous utilisons ici également
Descente, Lemme \ref{descent-lemma-sheaf-condition-holds} et
Définition \ref{descent-definition-structure-sheaf}.
Comme cela vaut pour tout objet $x$ de $\mathcal{X}$,
$\mathcal{F}$ est un faisceau pour chacune des trois topologies.
De plus, la définition de la quasi-cohérence et le fait que
$x$ soit un objet arbitraire de $\mathcal{X}$ montrent directement que $\mathcal{F}$
est quasi-cohérent pour chacune des trois topologies.
\end{proof}







\section{Modules localement quasi-cohérents}
\label{section-locally-quasi-coherent}

\noindent
Bien qu'il existe une variante pour la topologie de Zariski, la
topologie étale semble être la topologie naturelle à employer dans la
définition suivante.

\begin{definition}
\label{definition-locally-quasi-coherent}
Soit $p : \mathcal{X} \to (\Sch/S)_{fppf}$ une catégorie
fibrée en groupoïdes. Soit $\mathcal{F}$
un préfaisceau de $\mathcal{O}_\mathcal{X}$-modules.
On dit que $\mathcal{F}$ est {\it localement quasi-cohérent}\footnote{Cette
notation n'est pas standard.} si
$\mathcal{F}$ est un faisceau pour la topologie étale et si,
pour tout objet $x$ de $\mathcal{X}$, la restriction
$x^*\mathcal{F}|_{U_\etale}$ est un faisceau
quasi-cohérent. Ici $U = p(x)$.
\end{definition}

\noindent
Nous notons $\textit{LQCoh}(\mathcal{O}_\mathcal{X})$ la catégorie des
modules localement quasi-cohérents. Nous avons alors le diagramme suivant
de catégories de modules
$$
\xymatrix{
\QCoh(\mathcal{O}_\mathcal{X}) \ar[r] \ar[d] &
\textit{Mod}(\mathcal{O}_\mathcal{X}) \ar[d] \\
\textit{LQCoh}(\mathcal{O}_\mathcal{X}) \ar[r] &
\textit{Mod}(\mathcal{X}_\etale, \mathcal{O}_\mathcal{X})
}
$$
où les flèches sont des plongements strictement pleins.
Il se trouve que de nombreux résultats sur les faisceaux quasi-cohérents ont un
analogue pour les modules localement quasi-cohérents. De plus, à bien des
égards (comme nous le verrons plus loin), c'est une catégorie naturelle à considérer.
Par exemple, les faisceaux quasi-cohérents sont exactement les
modules localement quasi-cohérents qui sont « cartésiens », c'est-à-dire qui satisfont
la seconde condition du lemme ci-dessous.

\begin{lemma}
\label{lemma-quasi-coherent}
Soit $p : \mathcal{X} \to (\Sch/S)_{fppf}$ une catégorie
fibrée en groupoïdes. Soit $\mathcal{F}$
un préfaisceau de $\mathcal{O}_\mathcal{X}$-modules. Alors $\mathcal{F}$
est quasi-cohérent si et seulement si les deux conditions suivantes sont satisfaites :
\begin{enumerate}
\item $\mathcal{F}$ est localement quasi-cohérent, et
\item pour tout morphisme $\varphi : x \to y$ de $\mathcal{X}$ au-dessus de
$f : U \to V$, le morphisme de comparaison
$c_\varphi : f_{small}^*\mathcal{F}|_{V_\etale} \to
\mathcal{F}|_{U_\etale}$ de
(\ref{equation-comparison-modules}) est un isomorphisme.
\end{enumerate}
\end{lemma}

\begin{proof}
Supposons $\mathcal{F}$ quasi-cohérent. Alors $\mathcal{F}$ est un faisceau
pour la topologie fppf, donc un faisceau pour la topologie étale. De plus,
toute image réciproque de $\mathcal{F}$ sur un topos annelé est quasi-cohérente ; ainsi
les restrictions $x^*\mathcal{F}|_{U_\etale}$ sont quasi-cohérentes.
Cela prouve que $\mathcal{F}$ est localement quasi-cohérent.
Soit $y$ un objet de $\mathcal{X}$ tel que $V = p(y)$.
Nous avons vu que $\mathcal{X}/y = (\Sch/V)_{fppf}$. D'après
Descent, Proposition \ref{descent-proposition-equivalence-quasi-coherent},
il s'ensuit que $y^*\mathcal{F}$ est le module quasi-cohérent
associé à un module quasi-cohérent (usuel) $\mathcal{F}_V$ sur
le schéma $V$. Les morphismes de comparaison
(\ref{equation-comparison-modules}) sont donc bien des isomorphismes.

\medskip\noindent
Réciproquement, supposons que $\mathcal{F}$ satisfasse (1) et (2).
Soit $y$ un objet de $\mathcal{X}$ tel que $V = p(y)$. Notons
$\mathcal{F}_V$ le module quasi-cohérent sur
le schéma $V$ correspondant à la restriction
$y^*\mathcal{F}|_{V_\etale}$, laquelle est quasi-cohérente par
l'hypothèse (1), voir
Descent, Proposition \ref{descent-proposition-equivalence-quasi-coherent}.
La condition (2) signifie alors que les restrictions
$x^*\mathcal{F}|_{U_\etale}$, pour $x$ au-dessus de $y$, sont chacune
isomorphes à l'image réciproque de $\mathcal{F}_V$ (ou plutôt au faisceau étale associé)
par le morphisme de schémas correspondant $U \to V$.
Ainsi $y^*\mathcal{F}$ est le faisceau sur $(\Sch/V)_{fppf}$
associé à $\mathcal{F}_V$. Il est donc quasi-cohérent (encore d'après
Descent, Proposition \ref{descent-proposition-equivalence-quasi-coherent})
et le Lemme \ref{lemma-characterize-quasi-coherent} montre que
$\mathcal{F}$ est quasi-cohérent sur $\mathcal{X}$.
\end{proof}

\begin{lemma}
\label{lemma-pullback-lqc}
Soit $f : \mathcal{X} \to \mathcal{Y}$ un $1$-morphisme de catégories
fibrées en groupoïdes sur $(\Sch/S)_{fppf}$. Le foncteur image réciproque
$f^* = f^{-1} :
\textit{Mod}(\mathcal{Y}_\etale, \mathcal{O}_\mathcal{Y})
\to
\textit{Mod}(\mathcal{X}_\etale, \mathcal{O}_\mathcal{X})$
préserve les faisceaux localement quasi-cohérents.
\end{lemma}

\begin{proof}
Soit $\mathcal{G}$ localement quasi-cohérent sur $\mathcal{Y}$.
Choisissons un objet $x$ de $\mathcal{X}$ au-dessus du schéma $U$.
La restriction $x^*f^*\mathcal{G}|_{U_\etale}$ est égale à
$(f \circ x)^*\mathcal{G}|_{U_\etale}$
et est donc un faisceau quasi-cohérent par hypothèse sur $\mathcal{G}$.
\end{proof}

\begin{lemma}
\label{lemma-lqc-colimits}
Soit $p : \mathcal{X} \to (\Sch/S)_{fppf}$ une catégorie fibrée en
groupoïdes.
\begin{enumerate}
\item La catégorie $\textit{LQCoh}(\mathcal{O}_\mathcal{X})$
admet des colimites, qui coïncident avec les colimites dans la catégorie
$\textit{Mod}(\mathcal{X}_\etale, \mathcal{O}_\mathcal{X})$.
\item La catégorie $\textit{LQCoh}(\mathcal{O}_\mathcal{X})$
est abélienne, les noyaux et conoyaux étant calculés dans
$\textit{Mod}(\mathcal{X}_\etale, \mathcal{O}_\mathcal{X})$,
autrement dit, le foncteur d'inclusion est exact.
\item Étant donnée une suite exacte courte
$0 \to \mathcal{F}_1 \to \mathcal{F}_2 \to \mathcal{F}_3 \to 0$ de
$\textit{Mod}(\mathcal{X}_\etale, \mathcal{O}_\mathcal{X})$
si deux de ces trois modules sont localement quasi-cohérents, le troisième l'est aussi.
\item Étant donnés $\mathcal{F}, \mathcal{G}$ dans
$\textit{LQCoh}(\mathcal{O}_\mathcal{X})$
le produit tensoriel $\mathcal{F} \otimes_{\mathcal{O}_\mathcal{X}} \mathcal{G}$
dans $\textit{Mod}(\mathcal{X}_\etale, \mathcal{O}_\mathcal{X})$
est un objet de $\textit{LQCoh}(\mathcal{O}_\mathcal{X})$.
\item Étant donnés $\mathcal{F}, \mathcal{G}$ dans
$\textit{LQCoh}(\mathcal{O}_\mathcal{X})$
avec $\mathcal{F}$ de présentation finie sur
$\mathcal{X}_\etale$, le faisceau
$\SheafHom_{\mathcal{O}_\mathcal{X}}(\mathcal{F}, \mathcal{G})$
dans $\textit{Mod}(\mathcal{X}_\etale, \mathcal{O}_\mathcal{X})$
est un objet de $\textit{LQCoh}(\mathcal{O}_\mathcal{X})$.
\end{enumerate}
\end{lemma}

\begin{proof}
Dans les arguments ci-dessous, $x$ désigne un objet arbitraire de $\mathcal{X}$
au-dessus du schéma $U$. Pour montrer qu'un objet $\mathcal{H}$
de $\textit{Mod}(\mathcal{X}_\etale, \mathcal{O}_\mathcal{X})$ appartient
à $\textit{LQCoh}(\mathcal{O}_\mathcal{X})$, nous montrerons que la
restriction $x^*\mathcal{H}|_{U_\etale} = \mathcal{H}|_{U_\etale}$
est un objet quasi-cohérent de $\textit{Mod}(U_\etale, \mathcal{O}_U)$.

\medskip\noindent
Preuve de (1). Soit $\mathcal{I} \to \textit{LQCoh}(\mathcal{O}_\mathcal{X})$,
$i \mapsto \mathcal{F}_i$, un diagramme.
Considérons l'objet $\mathcal{F} = \colim_i \mathcal{F}_i$ de
$\textit{Mod}(\mathcal{X}_\etale, \mathcal{O}_\mathcal{X})$.
Le foncteur image réciproque $x^*$ commute à toutes les colimites,
puisqu'il est adjoint à gauche. Ainsi
$x^*\mathcal{F} = \colim_i x^*\mathcal{F}_i$. De même,
$x^*\mathcal{F}|_{U_\etale} =
\colim_i x^*\mathcal{F}_i|_{U_\etale}$.
Or, par hypothèse, chaque $x^*\mathcal{F}_i|_{U_\etale}$
est quasi-cohérent. Par conséquent, $\colim_i x^*\mathcal{F}_i|_{U_\etale}$
est quasi-cohérent d'après
Descent, Lemma \ref{descent-lemma-properties-quasi-coherent}.
Ainsi $x^*\mathcal{F}|_{U_\etale}$ est quasi-cohérent, comme voulu.

\medskip\noindent
Preuve de (2). Il résulte de (1) que les conoyaux existent dans
$\textit{LQCoh}(\mathcal{O}_\mathcal{X})$ et coïncident avec ceux calculés
dans $\textit{Mod}(\mathcal{X}_\etale, \mathcal{O}_\mathcal{X})$.
Soit $\varphi : \mathcal{F} \to \mathcal{G}$ un morphisme de
$\textit{LQCoh}(\mathcal{O}_\mathcal{X})$ et soit
$\mathcal{K} = \Ker(\varphi)$ calculé dans
$\textit{Mod}(\mathcal{X}_\etale, \mathcal{O}_\mathcal{X})$.
Si nous montrons que $\mathcal{K}$ est un module localement quasi-cohérent,
la preuve de (2) sera achevée. Remarquons pour cela que les noyaux
se calculent dans la catégorie des préfaisceaux (sans faisceautisation).
Ainsi $\mathcal{K}|_{U_\etale}$ est le noyau du morphisme
$\mathcal{F}|_{U_\etale} \to \mathcal{G}|_{U_\etale}$,
c'est-à-dire le noyau d'un morphisme de faisceaux quasi-cohérents sur $U_\etale$,
donc est quasi-cohérent d'après
Descent, Lemma \ref{descent-lemma-properties-quasi-coherent}.
Cela prouve (2).

\medskip\noindent
Preuve de (3). Soit
$0 \to \mathcal{F}_1 \to \mathcal{F}_2 \to \mathcal{F}_3 \to 0$
une suite exacte courte de
$\textit{Mod}(\mathcal{X}_\etale, \mathcal{O}_\mathcal{X})$.
Puisque nous employons la topologie étale, la restriction
$0 \to \mathcal{F}_1|_{U_\etale} \to
\mathcal{F}_2|_{U_\etale} \to \mathcal{F}_3|_{U_\etale} \to 0$
est elle aussi une suite exacte courte. L'assertion (3) résulte donc de l'énoncé
correspondant de
Descent, Lemma \ref{descent-lemma-properties-quasi-coherent}.

\medskip\noindent
Preuve de (4). Soient $\mathcal{F}$ et $\mathcal{G}$ dans
$\textit{LQCoh}(\mathcal{O}_\mathcal{X})$. Comme la restriction à $U_\etale$
est donnée par l'image réciproque suivant le morphisme de topos annelés
$U_\etale \to (\Sch/U)_\etale \to \mathcal{X}_\etale$
nous voyons que la restriction du produit tensoriel
$\mathcal{F} \otimes_{\mathcal{O}_\mathcal{X}} \mathcal{G}$
à $U_\etale$ est égale à
$\mathcal{F}|_{U_\etale} \otimes_{\mathcal{O}_U} \mathcal{G}|_{U_\etale}$,
voir Modules on Sites, Lemma \ref{sites-modules-lemma-tensor-product-pullback}.
Comme $\mathcal{F}|_{U_\etale}$ et $\mathcal{G}|_{U_\etale}$
sont quasi-cohérents, leur produit tensoriel l'est aussi, voir
Descent, Lemma \ref{descent-lemma-properties-quasi-coherent}.

\medskip\noindent
Preuve de (5). Soient $\mathcal{F}$ et $\mathcal{G}$ dans
$\textit{LQCoh}(\mathcal{O}_\mathcal{X})$, avec $\mathcal{F}$
de présentation finie. Puisque
$(\Sch/U)_\etale = \mathcal{X}_\etale/x$
est une localisation de $\mathcal{X}_\etale$ en un objet,
nous voyons que la restriction de
$\SheafHom_{\mathcal{O}_\mathcal{X}}(\mathcal{F}, \mathcal{G})$
à $(\Sch/U)_\etale$ est égale à
$$
\mathcal{H} =
\SheafHom_{\mathcal{O}|_{(\Sch/U)_\etale}}(
\mathcal{F}|_{(\Sch/U)_\etale}, \mathcal{G}|_{(\Sch/U)_\etale})
$$
par Modules on Sites, Lemma \ref{sites-modules-lemma-internal-hom-restriction}.
Le morphisme de topos annelés
$(U_\etale, \mathcal{O}_U) \to ((\Sch/U)_\etale, \mathcal{O})$
est plat, car l'image réciproque de $\mathcal{O}$ est $\mathcal{O}_U$.
L'image réciproque de $\mathcal{H}$ par ce morphisme est donc égale à
$\SheafHom_{\mathcal{O}_U}(\mathcal{F}|_{U_\etale}, \mathcal{G}|_{U_\etale})$
par Modules on Sites, Lemma \ref{sites-modules-lemma-pullback-internal-hom}.
Autrement dit, la restriction de
$\SheafHom_{\mathcal{O}_\mathcal{X}}(\mathcal{F}, \mathcal{G})$
à $U_\etale$ est
$\SheafHom_{\mathcal{O}_U}(\mathcal{F}|_{U_\etale}, \mathcal{G}|_{U_\etale})$.
Comme $\mathcal{F}|_{U_\etale}$ et $\mathcal{G}|_{U_\etale}$
sont quasi-cohérents, il en va de même de
$\SheafHom_{\mathcal{O}_U}(\mathcal{F}|_{U_\etale}, \mathcal{G}|_{U_\etale})$,
voir Descent, Lemma \ref{descent-lemma-properties-quasi-coherent}.
On conclut comme précédemment.
\end{proof}

\noindent
Dans le degré de généralité considéré ici, la catégorie des faisceaux quasi-cohérents
n'est pas abélienne. Voir Examples, Section \ref{examples-section-nonabelian-QCoh}.
Voici ce que nous pouvons démontrer sans travail supplémentaire.

\begin{lemma}
\label{lemma-qc-colimits}
Soit $p : \mathcal{X} \to (\Sch/S)_{fppf}$ une catégorie
fibrée en groupoïdes.
\begin{enumerate}
\item La catégorie $\QCoh(\mathcal{O}_\mathcal{X})$
admet des colimites, qui coïncident avec les colimites dans les catégories
$\textit{Mod}(\mathcal{X}_{Zar}, \mathcal{O}_\mathcal{X})$,
$\textit{Mod}(\mathcal{X}_\etale, \mathcal{O}_\mathcal{X})$,
$\textit{Mod}(\mathcal{O}_\mathcal{X})$, et
$\textit{LQCoh}(\mathcal{O}_\mathcal{X})$.
\item Étant donnés $\mathcal{F}, \mathcal{G}$ dans $\QCoh(\mathcal{O}_\mathcal{X})$,
les produits tensoriels $\mathcal{F} \otimes_{\mathcal{O}_\mathcal{X}} \mathcal{G}$
calculés dans $\textit{Mod}(\mathcal{X}_{Zar}, \mathcal{O}_\mathcal{X})$,
$\textit{Mod}(\mathcal{X}_\etale, \mathcal{O}_\mathcal{X})$, or
$\textit{Mod}(\mathcal{O}_\mathcal{X})$ coïncident, et leur valeur commune
est un objet de $\QCoh(\mathcal{O}_\mathcal{X})$.
\item Étant donnés $\mathcal{F}, \mathcal{G}$ dans $\QCoh(\mathcal{O}_\mathcal{X})$,
avec $\mathcal{F}$ localement libre de type fini (pour la topologie fppf, ou, de manière
équivalente, étale, ou encore de Zariski), les Hom internes
$\SheafHom_{\mathcal{O}_\mathcal{X}}(\mathcal{F}, \mathcal{G})$
calculés dans $\textit{Mod}(\mathcal{X}_{Zar}, \mathcal{O}_\mathcal{X})$,
$\textit{Mod}(\mathcal{X}_\etale, \mathcal{O}_\mathcal{X})$, or
$\textit{Mod}(\mathcal{O}_\mathcal{X})$ coïncident, et leur valeur commune
est un objet de $\QCoh(\mathcal{O}_\mathcal{X})$.
\end{enumerate}
\end{lemma}

\begin{proof}
Soit $x$ un objet arbitraire de $\mathcal{X}$ au-dessus du schéma $U$.
Soit $\tau \in \{Zariski, \etale, fppf\}$. Pour montrer qu'un objet $\mathcal{H}$
de $\textit{Mod}(\mathcal{X}_\tau, \mathcal{O}_\mathcal{X})$ appartient
à $\QCoh(\mathcal{O}_\mathcal{X})$, il suffit de montrer que la
restriction $x^*\mathcal{H}$ (Section \ref{section-restriction})
est un objet quasi-cohérent de $\textit{Mod}((\Sch/U)_\tau, \mathcal{O})$.
Voir les Lemmes \ref{lemma-characterize-quasi-coherent} et
\ref{lemma-characterize-quasi-coherent-bis}. Il en va de même pour la propriété d'être
localement libre de type fini. Rappelons que $(\Sch/U)_\tau = \mathcal{X}_\tau/x$
est une localisation de $\mathcal{X}_\tau$ en un objet. La restriction
commute donc aux colimites, aux produits tensoriels et à la formation du Hom interne
(voir Modules on Sites, Lemmes
\ref{sites-modules-lemma-exactness-pushforward-pullback},
\ref{sites-modules-lemma-tensor-product-pullback}, et
\ref{sites-modules-lemma-internal-hom-restriction}).
Le lemme se ramène ainsi à Descent, Lemma \ref{descent-lemma-qc-colimits}.
\end{proof}






\section{Champification et faisceaux}
\label{section-stackification}

\noindent
Il se trouve que la catégorie des faisceaux sur une catégorie fibrée en
groupoïdes ne « connaît » que la champification.

\begin{lemma}
\label{lemma-stackification}
Soit $f : \mathcal{X} \to \mathcal{Y}$ un $1$-morphisme de catégories
fibrées en groupoïdes sur $(\Sch/S)_{fppf}$. Si
$f$ induit une équivalence des champifications, alors le morphisme
de topos
$f : \Sh(\mathcal{X}_{fppf}) \to \Sh(\mathcal{Y}_{fppf})$
est une équivalence.
\end{lemma}

\begin{proof}
Nous pouvons supposer que $\mathcal{Y}$ est la champification de $\mathcal{X}$.
Nous affirmons que $f : \mathcal{X} \to \mathcal{Y}$ est un foncteur
cocontinu spécial, voir
Sites, Definition \ref{sites-definition-special-cocontinuous-functor},
ce qui prouvera le lemme. D'après
Stacks, Lemma \ref{stacks-lemma-topology-inherited-functorial},
le foncteur $f$ est continu et cocontinu. D'après
Stacks, Lemma \ref{stacks-lemma-stackify},
nous voyons que les conditions (3), (4) et (5) de
Sites, Lemma \ref{sites-lemma-equivalence},
sont satisfaites.
\end{proof}

\begin{lemma}
\label{lemma-stackification-quasi-coherent}
Soit $f : \mathcal{X} \to \mathcal{Y}$ un $1$-morphisme de catégories
fibrées en groupoïdes sur $(\Sch/S)_{fppf}$. Si
$f$ induit une équivalence des champifications, alors $f^*$
induit des équivalences
$\textit{Mod}(\mathcal{O}_\mathcal{X}) \to
\textit{Mod}(\mathcal{O}_\mathcal{Y})$
et
$\QCoh(\mathcal{O}_\mathcal{X}) \to
\QCoh(\mathcal{O}_\mathcal{Y})$.
\end{lemma}

\begin{proof}
Nous pouvons supposer que $\mathcal{Y}$ est la champification de $\mathcal{X}$.
La première assertion résulte clairement du
Lemma \ref{lemma-stackification}
et de
$\mathcal{O}_\mathcal{X} = f^{-1}\mathcal{O}_\mathcal{Y}$.
Les images réciproques de faisceaux quasi-cohérents sont quasi-cohérentes, voir
Lemme \ref{lemma-pullback-quasi-coherent}.
Il suffit donc de montrer que, si $f^*\mathcal{G}$ est
quasi-cohérent, alors $\mathcal{G}$ l'est.
Pour cela, soit $y$ un objet de $\mathcal{Y}$.
En explicitant la condition que $\mathcal{Y}$ est la champification
de $\mathcal{X}$, nous voyons qu'il existe un recouvrement fppf $\{y_i \to y\}$
dans $\mathcal{Y}$ tel que $y_i \cong f(x_i)$ pour un
objet $x_i$ de $\mathcal{X}$. Disons que $x_i$ et $y_i$ sont au-dessus du schéma $U_i$.
La quasi-cohérence de $f^*\mathcal{G}$ signifie alors que $x_i^*f^*\mathcal{G}$
est quasi-cohérent. Comme $x_i^*f^*\mathcal{G}$ est isomorphe à
$y_i^*\mathcal{G}$ (en tant que faisceaux sur $(\Sch/U_i)_{fppf}$), nous
voyons que $y_i^*\mathcal{G}$ est quasi-cohérent.
Il résulte de
Modules on Sites, Lemma \ref{sites-modules-lemma-local-final-object}
que la restriction de $\mathcal{G}$ à $\mathcal{Y}/y$ est
quasi-cohérente. Ainsi $\mathcal{G}$ est quasi-cohérent d'après le
Lemma \ref{lemma-characterize-quasi-coherent}.
\end{proof}





\section{Faisceaux quasi-cohérents et présentations}
\label{section-quasi-coherent-presentation}

\noindent
Commençons par comparer les faisceaux quasi-cohérents aux notions définies
précédemment pour les schémas et les espaces algébriques.

\begin{lemma}
\label{lemma-compare-locally-quasi-coherent}
Soit $S$ un schéma. Soit $\mathcal{X} \to (\Sch/S)_{fppf}$ une catégorie
fibrée en groupoïdes représentable par un espace algébrique $F$.
Si $\mathcal{F}$ appartient à $\textit{LQCoh}(\mathcal{O}_\mathcal{X})$,
alors la restriction $\mathcal{F}|_{F_\etale}$ (\ref{equation-restrict})
est quasi-cohérente.
\end{lemma}

\begin{proof}
Soit $U$ un schéma étale sur $F$. Alors
$\mathcal{F}|_{U_\etale} = (\mathcal{F}|_{F_\etale})|_{U_\etale}$.
C'est clair, mais voir aussi la Remarque \ref{remark-compare}.
L'assertion résulte donc des définitions.
\end{proof}

\begin{lemma}
\label{lemma-compare-quasi-coherent}
Soit $S$ un schéma. Soit $\mathcal{X} \to (\Sch/S)_{fppf}$ une catégorie
fibrée en groupoïdes représentable par un espace algébrique $F$.
Le foncteur (\ref{equation-restrict}) définit une équivalence
$$
\QCoh(\mathcal{O}_\mathcal{X}) \to \QCoh(\mathcal{O}_F),\quad
\mathcal{F} \longmapsto \mathcal{F}|_{F_\etale}
$$
dont un quasi-inverse est donné par $\mathcal{G} \mapsto \pi_F^*\mathcal{G}$.
Cette équivalence est compatible aux images réciproques pour les morphismes entre
catégories fibrées en groupoïdes représentables par des espaces algébriques.
\end{lemma}

\begin{proof}
D'après le Lemme \ref{lemma-characterize-quasi-coherent-bis}, nous pouvons travailler avec
la topologie étale. Nous emploierons sans autre mention les notations et résultats du
Lemme \ref{lemma-compare}.
Rappelons que le foncteur de restriction
$\textit{Mod}(\mathcal{X}_\etale, \mathcal{O}_\mathcal{X}) \to
\textit{Mod}(F_\etale, \mathcal{O}_F)$,
$\mathcal{F} \mapsto \mathcal{F}|_{F_\etale}$ est donné par $i_F^*$. D'après le
Lemme \ref{lemma-compare-locally-quasi-coherent}, ou d'après
Modules on Sites, Lemma \ref{sites-modules-lemma-local-pullback}
nous voyons que $\mathcal{F}|_{F_\etale}$ est quasi-cohérent si
$\mathcal{F}$ est quasi-cohérent.
Nous obtenons donc le foncteur indiqué dans l'énoncé du
lemme et, dans l'autre sens, un foncteur $\pi_F^*$.
Comme $\pi_F \circ i_F = \text{id}$, nous avons
$i_F^*\pi_F^*\mathcal{G} = \mathcal{G}$.

\medskip\noindent
Pour $\mathcal{F}$ dans
$\textit{Mod}(\mathcal{X}_\etale, \mathcal{O}_\mathcal{X})$
il existe un morphisme canonique $\pi_F^*(\mathcal{F}|_{F_\etale}) \to \mathcal{F}$,
à savoir le morphisme adjoint à l'identification
$\mathcal{F}|_{F_\etale} = \pi_{F, *}\mathcal{F}$.
Nous montrerons que ce morphisme est un isomorphisme si $\mathcal{F}$ est
un module quasi-cohérent sur $\mathcal{X}$.
Choisissons un schéma $U$ et un morphisme étale surjectif $U \to F$.
Notons $x : U \to \mathcal{X}$ l'objet correspondant de $\mathcal{X}$
au-dessus de $U$. Il suffit de montrer que
$\pi_F^*(\mathcal{F}|_{F_\etale}) \to \mathcal{F}$ est un isomorphisme
après restriction à $\mathcal{X}_\etale/x = (\Sch/U)_\etale$.
Comme $U \to F$ est étale, la Remarque \ref{remark-compare} donne
$$
\pi_F^*(\mathcal{F}|_{F_\etale})|_{\mathcal{X}_\etale/x} =
\pi_U^*(\mathcal{F}|_{U_\etale})
$$
et montre que la restriction du morphisme
$\pi_F^*(\mathcal{F}|_{F_\etale}) \to \mathcal{F}$
à $\mathcal{X}_\etale/x = (\Sch/U)_\etale$
est égale au morphisme correspondant
$\pi_U^*(\mathcal{F}|_{U_\etale}) \to \mathcal{F}|_{(\Sch/U)_\etale}$.
Puisque nous avons démontré le résultat pour les schémas dans
Descent, Section \ref{descent-section-quasi-coherent-sheaves}\footnote{En effet,
si $U$ est un schéma et $\mathcal{F}$ est quasi-cohérent sur
$(\Sch/U)_\etale$, alors $\mathcal{F} = \mathcal{H}^a$ pour un
module quasi-cohérent $\mathcal{H}$ sur le schéma $U$, d'après
Descent, Proposition \ref{descent-proposition-equivalence-quasi-coherent}.
Autrement dit, $\mathcal{F} = (\text{id}_{\etale,Zar})^*\mathcal{H}$
d'après Descent, Remark \ref{descent-remark-change-topologies-ringed-sites},
avec les notations de Descent, Lemma \ref{descent-lemma-compare-sites}.
Or nous avons
$\text{id}_{\etale,Zar} = \pi_U \circ \text{id}_{small,\etale,Zar}$
et nous voyons donc que $\mathcal{F} = \pi_U^*\mathcal{G}$, où
$\mathcal{G} = (\text{id}_{small,\etale,Zar})^*\mathcal{H}$ est
quasi-cohérent. Alors
$\pi_U^*i_U^*\mathcal{F} = \pi_U^*i_U^*\pi_U^*\mathcal{G} =
\pi_U^*\mathcal{G} = \mathcal{F}$, comme voulu.}
nous concluons.

\medskip\noindent
La compatibilité aux images réciproques résulte du fait que
le quasi-inverse est donné par $\pi_F^*$ et du
diagramme commutatif de topos annelés du
Lemme \ref{lemma-compare-morphism}.
\end{proof}

\noindent
Dans
Groupoids in Spaces, Definition
\ref{spaces-groupoids-definition-groupoid-module}
nous avons défini la notion de module quasi-cohérent
sur un groupoïde arbitraire. La proposition (formelle) suivante affirme
que nous pouvons étudier les faisceaux quasi-cohérents sur les champs quotients
en termes de modules quasi-cohérents sur les présentations.

\begin{proposition}
\label{proposition-quasi-coherent}
Soit $(U, R, s, t, c)$ un groupoïde en espaces algébriques sur $S$.
Soit $\mathcal{X} = [U/R]$ le champ quotient.
La catégorie des modules quasi-cohérents sur $\mathcal{X}$
est équivalente à la catégorie des modules quasi-cohérents
sur $(U, R, s, t, c)$.
\end{proposition}

\begin{proof}
Nous allons construire des foncteurs quasi-inverses
$$
\QCoh(\mathcal{O}_\mathcal{X})
\longleftrightarrow
\QCoh(U, R, s, t, c).
$$
où $\QCoh(U, R, s, t, c)$ désigne la catégorie des modules quasi-cohérents
sur le groupoïde $(U, R, s, t, c)$.

\medskip\noindent
Soit $\mathcal{F}$ un objet de $\QCoh(\mathcal{O}_\mathcal{X})$.
Notons $\mathcal{U}$ et $\mathcal{R}$ les catégories fibrées en groupoïdes
correspondant à $U$ et $R$. Notons $x$ l'objet (définissant) de $\mathcal{X}$
au-dessus de $U$. Rappelons que nous avons un diagramme $2$-commutatif
$$
\xymatrix{
\mathcal{R} \ar[r]_s \ar[d]_t & \mathcal{U} \ar[d]^x \\
\mathcal{U} \ar[r]^x & \mathcal{X}
}
$$
Voir Groupoids in Spaces, Lemma
\ref{spaces-groupoids-lemma-quotient-stack-2-arrow}.
D'après le Lemme \ref{lemma-2-morphisms-presheaves},
la $2$-flèche inhérente au diagramme induit un isomorphisme
$\alpha : t^*x^*\mathcal{F} \to s^*x^*\mathcal{F}$
qui satisfait la condition de cocycle sur
$\mathcal{R} \times_{s, \mathcal{U}, t} \mathcal{R}$;
c'est une conséquence de Groupoids in Spaces, Lemma
\ref{spaces-groupoids-lemma-quotient-stack-cocycle}.
Ainsi, si nous posons $\mathcal{G} = x^*\mathcal{F}|_{U_\etale}$,
l'équivalence de catégories du
Lemma \ref{lemma-compare-quasi-coherent}
(employée plusieurs fois de façon compatible aux images réciproques)
donne un isomorphisme
$\alpha : t_{small}^*\mathcal{G} \to s_{small}^*\mathcal{G}$
satisfaisant la condition de cocycle sur $R \times_{s, U, t} R$, c'est-à-dire que
$(\mathcal{G}, \alpha)$ est un objet de $\QCoh(U, R, s, t, c)$.
La règle $\mathcal{F} \mapsto (\mathcal{G}, \alpha)$ est notre
foncteur de gauche à droite.

\medskip\noindent
Construction du foncteur dans l'autre sens.
Soit $(\mathcal{G}, \alpha)$ un objet de $\QCoh(U, R, s, t, c)$.
D'après le Lemme \ref{lemma-stackification-quasi-coherent},
le morphisme de champification $[U/_{\!p}R] \to [U/R]$ (voir
Groupoids in Spaces, Definition
\ref{spaces-groupoids-definition-quotient-stack})
induit une équivalence des catégories de faisceaux quasi-cohérents.
Il suffit donc de construire un module quasi-cohérent $\mathcal{F}$
sur $[U/_{\!p}R]$.

\medskip\noindent
Rappelons qu'un objet $x = (T, u)$ de $[U/_{\!p}R]$
est donné par un schéma $T$ et un morphisme $u : T \to U$. Un morphisme
$(T, u) \to (T', u')$ est donné par une paire $(f, r)$, où $f : T \to T'$
et $r : T \to R$, avec $s \circ r = u$ et $t \circ r = u' \circ f$.
Appelons {\it morphisme spécial} tout morphisme de la forme
$(f, e \circ u' \circ f) : (T, u' \circ f) \to (T', u')$.
La catégorie des $(T, u), munie des morphismes spéciaux, n'est autre que la
catégorie des schémas sur $U$.

\medskip\noindent
Avec ces notations, pour un objet $(T, u)$ de $[U/_{\!p}R]$, posons
$$
\mathcal{F}(T, u) : = \Gamma(T, u_{small}^*\mathcal{G}).
$$
À un morphisme $(f, r) : (T, u) \to (T', u')$ est associé un morphisme
\begin{align*}
\mathcal{F}(T', u') & = \Gamma(T', (u')_{small}^*\mathcal{G}) \\
& \to \Gamma(T, f_{small}^*(u')_{small}^*\mathcal{G}) =
\Gamma(T, (u' \circ f)_{small}^*\mathcal{G}) \\
& = \Gamma(T, (t \circ r)_{small}^*\mathcal{G}) =
\Gamma(T, r_{small}^*t_{small}^*\mathcal{G}) \\
& \to \Gamma(T, r_{small}^*s_{small}^*\mathcal{G}) =
\Gamma(T, (s \circ r)_{small}^*\mathcal{G}) \\
& = \Gamma(T, u_{small}^*\mathcal{G}) \\
& = \mathcal{F}(T, u)
\end{align*}
où la première flèche est l'image réciproque suivant $f$ et la seconde est
$\alpha$. Remarquons que si $(f, r)$ est un morphisme spécial, ce
morphisme est simplement l'image réciproque suivant $f$, car $e_{small}^*\alpha = \text{id}$ par
les axiomes d'un faisceau de modules quasi-cohérents sur un groupoïde.
La condition de cocycle implique que $\mathcal{F}$ est un préfaisceau
de modules (détails omis). La description simple des morphismes de restriction
de $\mathcal{F}$ dans le cas d'un morphisme spécial montre que sa restriction à
$(\Sch/T)_{fppf}$ est quasi-cohérente.
Ainsi $\mathcal{F}$ est un faisceau quasi-cohérent sur $[U/_{\!p}R]$
(Lemme \ref{lemma-characterize-quasi-coherent}).

\medskip\noindent
Nous omettons de vérifier que les foncteurs construits ci-dessus sont
quasi-inverses l'un de l'autre.
\end{proof}

\noindent
Nous terminons cette section par un lemme technique sur les morphismes de source un faisceau
quasi-cohérent. C'est un analogue de
Schemes, Lemma \ref{schemes-lemma-compare-constructions}.
Nous verrons plus loin (Criteria for Representability, Theorem
\ref{criteria-theorem-flat-groupoid-gives-algebraic-stack})
que les hypothèses faites sur le groupoïde impliquent que $\mathcal{X}$ est
un champ algébrique.

\begin{lemma}
\label{lemma-map-from-quasi-coherent}
Soit $(U, R, s, t, c)$ un groupoïde en espaces algébriques sur $S$.
Supposons $s$ et $t$ plats et localement de présentation finie.
Soit $\mathcal{X} = [U/R]$ le champ quotient. Notons
$x$ l'objet de $\mathcal{X}$ au-dessus de $U$.
Soit $\mathcal{F}$ un
$\mathcal{O}_\mathcal{X}$-module quasi-cohérent, et soit $\mathcal{H}$ un objet quelconque
de $\textit{Mod}(\mathcal{O}_\mathcal{X})$.
Le morphisme
$$
\Hom_{\mathcal{O}_\mathcal{X}}(\mathcal{F}, \mathcal{H})
\longrightarrow
\Hom_{\mathcal{O}_U}(x^*\mathcal{F}|_{U_\etale},
x^*\mathcal{H}|_{U_\etale}),
\quad
\phi \longmapsto x^*\phi|_{U_\etale}
$$
est injectif, et son image est exactement formée des
$\varphi : x^*\mathcal{F}|_{U_\etale} \to
x^*\mathcal{H}|_{U_\etale}$ qui donnent lieu à un
diagramme commutatif
$$
\xymatrix{
s_{small}^*(x^*\mathcal{F}|_{U_\etale})
\ar[r] \ar[d]^{s_{small}^*\varphi} &
(x \circ s)^*\mathcal{F}|_{R_\etale} =
(x \circ t)^*\mathcal{F}|_{R_\etale} &
t_{small}^*(x^*\mathcal{F}|_{U_\etale})
\ar[l] \ar[d]_{t_{small}^*\varphi} \\
s_{small}^*(x^*\mathcal{H}|_{U_\etale})
\ar[r] &
(x \circ s)^*\mathcal{H}|_{R_\etale} =
(x \circ t)^*\mathcal{H}|_{R_\etale} &
t_{small}^*(x^*\mathcal{H}|_{U_\etale})
\ar[l]
}
$$
de modules sur $R_\etale$,
où les flèches horizontales sont les morphismes de comparaison
(\ref{equation-comparison-algebraic-spaces-modules}).
\end{lemma}

\begin{proof}
D'après le
Lemma \ref{lemma-stackification-quasi-coherent}
le morphisme de champification $[U/_{\!p}R] \to [U/R]$ (voir
Groupoids in Spaces, Definition
\ref{spaces-groupoids-definition-quotient-stack})
induit une équivalence des catégories de faisceaux quasi-cohérents
et de $\mathcal{O}$-modules fppf.
Il suffit donc de démontrer le lemme pour $\mathcal{X} = [U/_{\!p}R]$.
D'après la Proposition \ref{proposition-quasi-coherent}
et sa preuve, il existe un module quasi-cohérent
$(\mathcal{G}, \alpha)$ sur $(U, R, s, t, c)$ tel que
$\mathcal{F}$ soit donné par la règle
$\mathcal{F}(T, u) = \Gamma(T, u^*\mathcal{G})$.
En particulier, $x^*\mathcal{F}|_{U_\etale} = \mathcal{G}$
et il est clair que le morphisme de l'énoncé du
lemme est injectif. De plus, étant donnés un morphisme
$\varphi : \mathcal{G} \to x^*\mathcal{H}|_{U_\etale}$
et un objet quelconque
$y = (T, u)$ de $[U/_{\!p}R]$, nous pouvons considérer le morphisme
$$
\mathcal{F}(y) = \Gamma(T, u^*\mathcal{G})
\xrightarrow{u_{small}^*\varphi}
\Gamma(T, u_{small}^*x^*\mathcal{H}|_{U_\etale})
\rightarrow
\Gamma(T, y^*\mathcal{H}|_{T_\etale}) = \mathcal{H}(y)
$$
où la seconde flèche est le morphisme de comparaison
(\ref{equation-comparison-modules}) pour le faisceau $\mathcal{H}$.
Cette construction est compatible aux morphismes de restriction des
faisceaux $\mathcal{F}$ et $\mathcal{G}$ pour les morphismes de
$[U/_{\!p}R]$ si la condition de cocycle du
lemme est satisfaite. Preuve omise. Indication : les morphismes de restriction
de $\mathcal{F}$ sont explicités en termes de $(\mathcal{G}, \alpha)$
dans la preuve de la
Proposition \ref{proposition-quasi-coherent}.
\end{proof}







\section{Faisceaux quasi-cohérents sur les champs algébriques}
\label{section-quasi-coherent-algebraic-stacks}

\noindent
Soit $\mathcal{X}$ un champ algébrique sur $S$. D'après
Algebraic Stacks, Lemma \ref{algebraic-lemma-stack-presentation}
il existe une équivalence $[U/R] \to \mathcal{X}$,
où $(U, R, s, t, c)$ est un groupoïde lisse en espaces algébriques.
Alors
$$
\QCoh(\mathcal{O}_\mathcal{X})
\cong
\QCoh(\mathcal{O}_{[U/R]})
\cong
\QCoh(U, R, s, t, c)
$$
où la seconde équivalence est la
Proposition \ref{proposition-quasi-coherent}.
Ainsi, la catégorie des faisceaux quasi-cohérents sur un champ algébrique
est équivalente à la catégorie des modules quasi-cohérents sur un groupoïde
lisse en espaces algébriques. En particulier, d'après
Groupoids in Spaces, Lemma \ref{spaces-groupoids-lemma-abelian}
nous voyons que $\QCoh(\mathcal{O}_\mathcal{X})$ est abélienne !

\medskip\noindent
Notre cadre actuel présente un aspect légèrement déconcertant :
le plongement pleinement fidèle
$$
\QCoh(\mathcal{O}_\mathcal{X})
\longrightarrow
\textit{Mod}(\mathcal{O}_\mathcal{X})
$$
n'est en général {\bf pas} exact. Toutefois, il en va exactement de même
pour les schémas : pour la plupart des schémas $X$, le plongement
$$
\QCoh(\mathcal{O}_X) \cong
\QCoh((\Sch/X)_{fppf}, \mathcal{O}_X) \longrightarrow
\textit{Mod}((\Sch/X)_{fppf}, \mathcal{O}_X)
$$
n'est pas exact, voir
Descent, Lemma \ref{descent-lemma-equivalence-quasi-coherent-limits}.
Incidemment, l'exemple donné dans la preuve de
Descent, Lemma \ref{descent-lemma-equivalence-quasi-coherent-limits}
montre qu'en général le plongement strictement plein
$\QCoh(\mathcal{O}_\mathcal{X}) \to
\textit{LQCoh}(\mathcal{O}_\mathcal{X})$ isn't exact either.

\medskip\noindent
Nous réunissons tous les résultats obtenus jusqu'ici en un seul énoncé.

\begin{lemma}
\label{lemma-quasi-coherent-algebraic-stack}
Soit $\mathcal{X}$ un champ algébrique sur $S$.
\begin{enumerate}
\item Si $[U/R] \to \mathcal{X}$ est une présentation de $\mathcal{X}$,
il existe une équivalence canonique
$\QCoh(\mathcal{O}_\mathcal{X}) \cong
\QCoh(U, R, s, t, c)$.
\item La catégorie $\QCoh(\mathcal{O}_\mathcal{X})$ est abélienne.
\item Le foncteur d'inclusion $\QCoh(\mathcal{O}_\mathcal{X}) \to
\textit{Mod}(\mathcal{O}_\mathcal{X})$ est exact à droite, mais
n'est en général {\bf pas} exact.
\item La catégorie $\QCoh(\mathcal{O}_\mathcal{X})$
admet des colimites, qui coïncident avec les colimites dans la catégorie
$\textit{Mod}(\mathcal{O}_\mathcal{X})$.
\item Étant donnés $\mathcal{F}, \mathcal{G}$ dans
$\QCoh(\mathcal{O}_\mathcal{X})$
le produit tensoriel $\mathcal{F} \otimes_{\mathcal{O}_\mathcal{X}} \mathcal{G}$
dans $\textit{Mod}(\mathcal{O}_\mathcal{X})$
est un objet de $\QCoh(\mathcal{O}_\mathcal{X})$.
\item Étant donnés $\mathcal{F}, \mathcal{G}$ dans
$\QCoh(\mathcal{O}_\mathcal{X})$
avec $\mathcal{F}$ localement libre de type fini, le faisceau
$\SheafHom_{\mathcal{O}_\mathcal{X}}(\mathcal{F}, \mathcal{G})$
dans $\textit{Mod}(\mathcal{O}_\mathcal{X})$
est un objet de $\QCoh(\mathcal{O}_\mathcal{X})$.
\item Étant donnée une suite exacte courte
$0 \to \mathcal{F}_1 \to \mathcal{F}_2 \to \mathcal{F}_3 \to 0$
dans $\textit{Mod}(\mathcal{O}_\mathcal{X})$,
si $\mathcal{F}_1$ et $\mathcal{F}_3$ sont quasi-cohérents, alors
$\mathcal{F}_2$ est quasi-cohérent.
\end{enumerate}
\end{lemma}

\begin{proof}
Les propriétés (4), (5) et (6) ont été démontrées dans le
Lemma \ref{lemma-qc-colimits}.
L'assertion (1) est la
Proposition \ref{proposition-quasi-coherent}.
L'assertion (2) résulte de (1) et de
Groupoids in Spaces, Lemma \ref{spaces-groupoids-lemma-abelian}
comme expliqué ci-dessus. L'exactitude à droite du foncteur d'inclusion
dans (3) résulte de (4) ; comparer avec
Homology, Lemma \ref{homology-lemma-exact-functor}.
Pour la non-exactitude du foncteur d'inclusion dans (3), voir
Descent, Lemma \ref{descent-lemma-equivalence-quasi-coherent-limits}.
Pour (7), remarquons qu'il suffit de vérifier que la restriction
de $\mathcal{F}_2$ au grand site d'un schéma est quasi-cohérente
(Lemme \ref{lemma-characterize-quasi-coherent}) ;
cela résulte donc de l'assertion correspondante de
Descent, Lemma \ref{descent-lemma-equivalence-quasi-coherent-limits}.
\end{proof}

\noindent
Construisons maintenant le cohérateur des modules sur un champ algébrique.

\begin{proposition}
\label{proposition-coherator}
Soit $\mathcal{X}$ un champ algébrique sur $S$.
\begin{enumerate}
\item La catégorie $\QCoh(\mathcal{O}_\mathcal{X})$ est une catégorie
abélienne de Grothendieck. Par conséquent, $\QCoh(\mathcal{O}_\mathcal{X})$
a assez d'injectifs et admet toutes les limites.
\item Le foncteur d'inclusion
$\QCoh(\mathcal{O}_\mathcal{X}) \to
\textit{Mod}(\mathcal{O}_\mathcal{X})$ admet un adjoint à droite\footnote{Ce
foncteur est parfois appelé le {\it cohérateur}.}
$$
Q :
\textit{Mod}(\mathcal{O}_\mathcal{X})
\longrightarrow
\QCoh(\mathcal{O}_\mathcal{X})
$$
tel que, pour tout faisceau quasi-cohérent $\mathcal{F}$, le morphisme d'adjonction
$Q(\mathcal{F}) \to \mathcal{F}$ soit un isomorphisme.
\end{enumerate}
\end{proposition}

\begin{proof}
Cette preuve reprend celle du cas des schémas, voir
Properties, Proposition \ref{properties-proposition-coherator}
et celle du cas des espaces algébriques, voir
Properties of Spaces, Proposition
\ref{spaces-properties-proposition-coherator}.
Nous conseillons au lecteur de lire d'abord l'une de ces preuves.

\medskip\noindent
L'assertion (1) signifie que $\QCoh(\mathcal{O}_\mathcal{X})$ (a) admet toutes les colimites,
(b) que les colimites filtrantes sont exactes, et (c) qu'elle admet un générateur, voir
Injectives, Section \ref{injectives-section-grothendieck-conditions}.
D'après le Lemme \ref{lemma-quasi-coherent-algebraic-stack},
les colimites dans $\QCoh(\mathcal{O}_X)$ existent et coïncident
avec celles de $\textit{Mod}(\mathcal{O}_X)$. D'après
Modules on Sites, Lemma \ref{sites-modules-lemma-limits-colimits}
les colimites filtrantes sont exactes. Ainsi (a) et (b) sont satisfaites.

\medskip\noindent
Choisissons une présentation $\mathcal{X} = [U/R]$ telle que $(U, R, s, t, c)$
soit un groupoïde lisse en espaces algébriques ; en particulier, $s$ et $t$
sont des morphismes plats d'espaces algébriques. D'après le
Lemma \ref{lemma-quasi-coherent-algebraic-stack}
ci-dessus, nous avons
$\QCoh(\mathcal{O}_\mathcal{X}) = \QCoh(U, R, s, t, c)$.
D'après Groupoids in Spaces, Lemma \ref{spaces-groupoids-lemma-set-generators},
il existe un ensemble $T$ et une famille $(\mathcal{F}_t)_{t \in T}$ de
faisceaux quasi-cohérents sur $\mathcal{X}$ tels que tout faisceau quasi-cohérent
sur $\mathcal{X}$ soit la colimite filtrante de ses sous-faisceaux
isomorphes à l'un des $\mathcal{F}_t$.
Ainsi $\bigoplus_t \mathcal{F}_t$ est
un générateur de $\QCoh(\mathcal{O}_X)$, et (c) est satisfait.
Les assertions sur les limites et les injectifs valent dans toute
catégorie abélienne de Grothendieck, voir
Injectives, Theorem
\ref{injectives-theorem-injective-embedding-grothendieck} and
Lemma \ref{injectives-lemma-grothendieck-products}.

\medskip\noindent
Preuve de (2). Pour construire $Q$, nous employons le procédé général suivant.
Étant donné un objet $\mathcal{F}$ de $\textit{Mod}(\mathcal{O}_\mathcal{X})$,
considérons le foncteur
$$
\QCoh(\mathcal{O}_\mathcal{X})^{opp}
\longrightarrow
\textit{Sets},
\quad
\mathcal{G}
\longmapsto
\Hom_\mathcal{X}(\mathcal{G}, \mathcal{F})
$$
Ce foncteur transforme les colimites en limites ;
il est donc représentable, voir
Injectives, Lemma \ref{injectives-lemma-grothendieck-brown}.
Il existe donc un faisceau quasi-cohérent $Q(\mathcal{F})$
et un isomorphisme fonctoriel
$\Hom_\mathcal{X}(\mathcal{G}, \mathcal{F}) =
\Hom_\mathcal{X}(\mathcal{G}, Q(\mathcal{F}))$
pour $\mathcal{G}$ dans $\QCoh(\mathcal{O}_\mathcal{X})$.
D'après le lemme de Yoneda
(Categories, Lemma \ref{categories-lemma-yoneda})
la construction $\mathcal{F} \leadsto Q(\mathcal{F})$ est
fonctorielle en $\mathcal{F}$. Par construction, $Q$ est un adjoint
à droite du foncteur d'inclusion.
Le fait que $Q(\mathcal{F}) \to \mathcal{F}$ soit un isomorphisme
lorsque $\mathcal{F}$ est quasi-cohérent est une conséquence formelle du fait
que le foncteur d'inclusion
$\QCoh(\mathcal{O}_\mathcal{X}) \to
\textit{Mod}(\mathcal{O}_\mathcal{X})$
est pleinement fidèle.
\end{proof}








\section{Cohomologie}
\label{section-cohomology-general}

\noindent
Soit $S$ un schéma et soit $\mathcal{X}$ une catégorie fibrée en groupoïdes
sur $(\Sch/S)_{fppf}$. Pour tout $\tau \in \{Zariski, \etale, smooth,
syntomic, fppf\}$, les catégories $\textit{Ab}(\mathcal{X}_\tau)$ et
$\textit{Mod}(\mathcal{X}_\tau, \mathcal{O}_\mathcal{X})$ ont
assez d'injectifs, voir
Injectives, Theorems \ref{injectives-theorem-sheaves-injectives} and
\ref{injectives-theorem-sheaves-modules-injectives}.
Nous pouvons donc employer le formalisme de
Cohomology on Sites, Section \ref{sites-cohomology-section-cohomology-sheaves}
pour définir les groupes de cohomologie
$$
H^p(\mathcal{X}_\tau, \mathcal{F}) = H^p_\tau(\mathcal{X}, \mathcal{F})
\quad\text{et}\quad
H^p(x, \mathcal{F}) = H^p_\tau(x, \mathcal{F})
$$
pour tout $x \in \Ob(\mathcal{X})$ et tout objet $\mathcal{F}$ de
$\textit{Ab}(\mathcal{X}_\tau)$ or
$\textit{Mod}(\mathcal{X}_\tau, \mathcal{O}_\mathcal{X})$. De plus, si
$f : \mathcal{X} \to \mathcal{Y}$ est un $1$-morphisme de catégories fibrées
en groupoïdes sur $(\Sch/S)_{fppf}$, on obtient les images directes
supérieures $R^if_*\mathcal{F}$ dans $\textit{Ab}(\mathcal{Y}_\tau)$ ou
$\textit{Mod}(\mathcal{Y}_\tau, \mathcal{O}_\mathcal{Y})$.
Bien entendu, comme expliqué dans
Cohomology on Sites, Section \ref{sites-cohomology-section-derived-functors}
il existe aussi des versions dérivées de $H^p(-)$ et de $R^if_*$.

\begin{lemma}
\label{lemma-cohomology-restriction}
Soit $S$ un schéma. Soit $\mathcal{X}$ une catégorie fibrée en groupoïdes
sur $(\Sch/S)_{fppf}$. Soit $\tau \in \{Zariski, \etale, smooth,
syntomic, fppf\}$. Soit $x \in \Ob(\mathcal{X})$ un objet au-dessus
du schéma $U$. Soit $\mathcal{F}$
un objet de $\textit{Ab}(\mathcal{X}_\tau)$ ou de
$\textit{Mod}(\mathcal{X}_\tau, \mathcal{O}_\mathcal{X})$. Alors
$$
H^p_\tau(x, \mathcal{F}) = H^p((\Sch/U)_\tau, x^{-1}\mathcal{F})
$$
et, si $\tau = \etale$, nous avons aussi
$$
H^p_\etale(x, \mathcal{F}) =
H^p(U_\etale, \mathcal{F}|_{U_\etale}).
$$
\end{lemma}

\begin{proof}
La première assertion résulte de
Cohomology on Sites, Lemma \ref{sites-cohomology-lemma-cohomology-of-open}
et de l'équivalence du
Lemma \ref{lemma-localizing-structure-sheaf}.
La seconde assertion résulte de la première, combinée avec
Étale Cohomology, Lemma
\ref{etale-cohomology-lemma-compare-cohomology-big-small}.
\end{proof}














\section{Faisceaux injectifs}
\label{section-lower-shriek}

\noindent
L'image directe d'un faisceau abélien injectif ou d'un module injectif est injective.

\begin{lemma}
\label{lemma-pushforward-injective}
Soit $f : \mathcal{X} \to \mathcal{Y}$ un $1$-morphisme de catégories
fibrées en groupoïdes sur $(\Sch/S)_{fppf}$. Soit
$\tau \in \{Zar, \etale, smooth, syntomic, fppf\}$.
\begin{enumerate}
\item $f_*\mathcal{I}$ est injectif dans $\textit{Ab}(\mathcal{Y}_\tau)$
pour $\mathcal{I}$ injectif dans $\textit{Ab}(\mathcal{X}_\tau)$, et
\item $f_*\mathcal{I}$ est injectif dans
$\textit{Mod}(\mathcal{Y}_\tau, \mathcal{O}_\mathcal{Y})$
pour $\mathcal{I}$ injectif dans
$\textit{Mod}(\mathcal{X}_\tau, \mathcal{O}_\mathcal{X})$.
\end{enumerate}
\end{lemma}

\begin{proof}
Cela résulte formellement du fait que $f^{-1}$ est un adjoint à gauche
exact de $f_*$, voir
Homology, Lemma \ref{homology-lemma-adjoint-preserve-injectives}.
\end{proof}

\noindent
Dans la suite de cette section, nous prouvons que l'image réciproque $f^{-1}$ admet un
adjoint à gauche $f_!$ sur les faisceaux abéliens et les modules. Si $f$ est représentable (par
des schémas ou des espaces algébriques), il se trouvera que $f_!$ est exact
et que $f^{-1}$ préserve les injectifs. Nous démontrons d'abord quelques
lemmes préliminaires sur les produits fibrés et les égalisateurs dans les catégories
fibrées en groupoïdes, ainsi que sur leur comportement à l'égard des morphismes.

\begin{lemma}
\label{lemma-fibre-products}
Soit $p : \mathcal{X} \to (\Sch/S)_{fppf}$
une catégorie fibrée en groupoïdes.
\begin{enumerate}
\item La catégorie $\mathcal{X}$ admet des produits fibrés.
\item Si les préfaisceaux $\mathit{Isom}$ de $\mathcal{X}$
sont représentables par des espaces algébriques, alors $\mathcal{X}$ admet des égalisateurs.
\item Si $\mathcal{X}$ est un champ algébrique (ou, plus généralement,
un champ quotient), alors $\mathcal{X}$ admet des égalisateurs.
\end{enumerate}
\end{lemma}

\begin{proof}
L'assertion (1) résulte de
Categories, Lemma \ref{categories-lemma-fibred-groupoids-fibre-product-goes-up}
puisque $(\Sch/S)_{fppf}$ admet des produits fibrés.

\medskip\noindent
Soient $a, b : x \to y$ des morphismes de $\mathcal{X}$.
Posons $U = p(x)$ et $V = p(y)$. La catégorie des schémas admet des égalisateurs ;
on peut donc prendre pour $W \to U$ l'égalisateur de $p(a)$ et $p(b)$.
Notons $c : z \to x$ un morphisme de $\mathcal{X}$ au-dessus de $W \to U$.
L'égalisateur de $a$ et $b$, s'il existe, est celui de $a \circ c$
et $b \circ c$. Nous pouvons donc supposer que $p(a) = p(b) = f : U \to V$.
Comme $\mathcal{X}$ est fibrée en groupoïdes, il existe un unique automorphisme
$i : x \to x$ dans la catégorie fibre de $\mathcal{X}$ au-dessus de $U$ tel que
$a \circ i = b$. De nouveau, l'égalisateur de $a$ et $b$ est celui de
$\text{id}_x$ et $i$. Rappelons que $\mathit{Isom}_\mathcal{X}(x)$
est le préfaisceau sur $(\Sch/U)_{fppf}$ qui associe à
$T/U$ l'ensemble des automorphismes de $x|_T$ dans la catégorie fibre
de $\mathcal{X}$ au-dessus de $T$, voir
Stacks, Definition \ref{stacks-definition-mor-presheaf}.
Si $\mathit{Isom}_\mathcal{X}(x)$ est représentable par un espace algébrique
$G \to U$, alors $\text{id}_x$ et $i$ définissent des morphismes
$e, i : U \to G$ au-dessus de $U$. Posons $M = U \times_{e, G, i} U$, qui, d'après
Morphisms of Spaces, Lemma \ref{spaces-morphisms-lemma-section-immersion}
est un schéma. Il est alors clair que $x|_M \to x$ est l'égalisateur des
morphismes $\text{id}_x$ et $i$ dans $\mathcal{X}$.
Cela prouve (2).

\medskip\noindent
Si $\mathcal{X} = [U/R]$ pour un groupoïde en espaces algébriques
$(U, R, s, t, c)$ sur $S$, l'hypothèse de (2) est satisfaite d'après
Bootstrap, Lemma \ref{bootstrap-lemma-quotient-stack-isom}.
Si $\mathcal{X}$ est un champ algébrique, nous pouvons choisir une
présentation $[U/R] \cong \mathcal{X}$ d'après
Algebraic Stacks, Lemma \ref{algebraic-lemma-stack-presentation}.
\end{proof}

\begin{lemma}
\label{lemma-fibre-products-morphism}
Soit $f : \mathcal{X} \to \mathcal{Y}$ un $1$-morphisme de catégories
fibrées en groupoïdes sur $(\Sch/S)_{fppf}$.
\begin{enumerate}
\item Le foncteur $f$ transforme les produits fibrés en produits fibrés.
\item Si $f$ est fidèle, alors $f$ transforme les égalisateurs en égalisateurs.
\end{enumerate}
\end{lemma}

\begin{proof}
D'après
Categories, Lemma \ref{categories-lemma-fibred-groupoids-fibre-product-goes-up}
un produit fibré dans $\mathcal{X}$ est tout carré commutatif au-dessus
d'un diagramme de produit fibré dans $(\Sch/S)_{fppf}$. Il en va de même pour
$\mathcal{Y}$. L'assertion (1) est donc claire.

\medskip\noindent
Soit $x \to x'$ l'égalisateur de deux morphismes $a, b : x' \to x''$
dans $\mathcal{X}$. Montrons que $f(x) \to f(x')$ est l'égalisateur
de $f(a)$ et $f(b)$. Soit $y \to f(x')$ un morphisme de $\mathcal{Y}$
égalisant $f(a)$ et $f(b)$. Disons que $x, x', x''$ sont au-dessus des schémas
$U, U', U''$ et que $y$ est au-dessus de $V$. Notons $h : V \to U'$ l'image
de $y \to f(x')$ dans la catégorie des schémas. D'après les axiomes des catégories
fibrées, le morphisme $y \to f(x')$ est isomorphe à $f(h^*x') \to f(x')$.
Comme $f$ est fidèle, nous voyons donc que
$h^*x' \to x'$ égalise $a$ et $b$. Nous obtenons ainsi un unique morphisme
$h^*x' \to x$ dont l'image $y = f(h^*x') \to f(x)$ est le morphisme cherché
dans $\mathcal{Y}$.
\end{proof}

\begin{lemma}
\label{lemma-fibre-products-preserve-properties}
Soient $f : \mathcal{X} \to \mathcal{Y}$ et $g : \mathcal{Z} \to \mathcal{Y}$
des $1$-morphismes fidèles de catégories
fibrées en groupoïdes sur $(\Sch/S)_{fppf}$.
\begin{enumerate}
\item le foncteur $\mathcal{X} \times_\mathcal{Y} \mathcal{Z} \to \mathcal{Y}$
est fidèle, et
\item si $\mathcal{X}$ et $\mathcal{Z}$ admettent des égalisateurs, il en va de même de
$\mathcal{X} \times_\mathcal{Y} \mathcal{Z}$.
\end{enumerate}
\end{lemma}

\begin{proof}
Nous considérons les objets de $\mathcal{X} \times_\mathcal{Y} \mathcal{Z}$ comme des
quadruplets $(U, x, z, \alpha)$, où $\alpha : f(x) \to g(z)$ est un
isomorphisme au-dessus de $U$, voir
Categories, Lemma \ref{categories-lemma-2-product-categories-over-C}.
Un morphisme $(U, x, z, \alpha) \to (U', x', z', \alpha')$ est une
paire de morphismes $a : x \to x'$ et $b : z \to z'$ compatibles
à $\alpha$ et $\alpha'$. Il est donc clair que, si $f$ et
$g$ sont fidèles, le foncteur
$\mathcal{X} \times_\mathcal{Y} \mathcal{Z} \to \mathcal{Y}$.
est fidèle. Supposons maintenant que
$(a, b), (a', b') : (U, x, z, \alpha) \to (U', x', z', \alpha')$
soient deux morphismes du $2$-produit fibré. Considérons l'égalisateur
$x'' \to x$ de $a$ et $a'$, et l'égalisateur $z'' \to z$ de $b$ et $b'$.
Comme $f$ commute aux égalisateurs (d'après le
Lemma \ref{lemma-fibre-products-morphism})
nous voyons que $f(x'') \to f(x)$ est l'égalisateur de $f(a)$ et $f(a')$.
De même, $g(z'') \to g(z)$ est l'égalisateur de $g(b)$ et $g(b')$.
Considérons le diagramme
$$
\xymatrix{
f(x'') \ar[r] \ar@{..>}[d]_{\alpha''}&
f(x) \ar[d]_\alpha
\ar@<0.5ex>[r]^{f(a)}
\ar@<-0.5ex>[r]_{f(a')}
 &
f(x') \ar[d]^{\alpha'} \\
g(z'') \ar[r] &
g(z)
\ar@<0.5ex>[r]^{g(b)}
\ar@<-0.5ex>[r]_{g(b')}
 &
g(z')
}
$$
Il est clair que la flèche pointillée existe et est un isomorphisme.
Toutefois, il n'est pas assuré a priori que l'image de $\alpha''$
dans la catégorie des schémas soit l'identité de sa source. D'autre
part, l'existence de $\alpha''$ signifie que nous pouvons supposer que $x''$
et $z''$ sont définis sur le même schéma et que les morphismes
$x'' \to x$ et $z'' \to z$ ont la même image dans la catégorie des schémas.
En refaisant le diagramme ci-dessus, nous voyons que la flèche pointillée se
projette alors bien sur un morphisme identité, ce qui conclut. Quelques détails sont omis.
\end{proof}

\noindent
Puisque nous travaillons avec de grands sites, nous obtenons le résultat quelque peu
contre-intuitif suivant (qui vaut aussi pour les morphismes entre grands sites
de schémas). Attention : ce résultat est faux si l'on omet l'hypothèse
que $f$ est fidèle.

\begin{lemma}
\label{lemma-pullback-injective}
Soit $f : \mathcal{X} \to \mathcal{Y}$ un $1$-morphisme de catégories
fibrées en groupoïdes sur $(\Sch/S)_{fppf}$. Soit
$\tau \in \{Zar, \etale, smooth, syntomic, fppf\}$.
Le foncteur
$f^{-1} : \textit{Ab}(\mathcal{Y}_\tau) \to \textit{Ab}(\mathcal{X}_\tau)$
admet un adjoint à gauche
$f_! : \textit{Ab}(\mathcal{X}_\tau) \to \textit{Ab}(\mathcal{Y}_\tau)$.
Si $f$ est fidèle et si $\mathcal{X}$ admet des égalisateurs, alors
\begin{enumerate}
\item $f_!$ est exact, et
\item $f^{-1}\mathcal{I}$ est injectif dans $\textit{Ab}(\mathcal{X}_\tau)$
pour $\mathcal{I}$ injectif dans $\textit{Ab}(\mathcal{Y}_\tau)$.
\end{enumerate}
\end{lemma}

\begin{proof}
D'après
Stacks, Lemma \ref{stacks-lemma-topology-inherited-functorial}
le foncteur $f$ est continu et cocontinu. Par conséquent, d'après
Modules on Sites, Lemma \ref{sites-modules-lemma-g-shriek-adjoint}
le foncteur
$f^{-1} : \textit{Ab}(\mathcal{Y}_\tau) \to \textit{Ab}(\mathcal{X}_\tau)$
admet un adjoint à gauche
$f_! : \textit{Ab}(\mathcal{X}_\tau) \to \textit{Ab}(\mathcal{Y}_\tau)$.
Pour (1), nous appliquons
Modules on Sites, Lemma \ref{sites-modules-lemma-exactness-lower-shriek}
et, pour vérifier les hypothèses de ce lemme, nous employons les
Lemmes \ref{lemma-fibre-products} et
\ref{lemma-fibre-products-morphism}
ci-dessus. L'assertion (2) en résulte formellement, voir
Homology, Lemma \ref{homology-lemma-adjoint-preserve-injectives}.
\end{proof}

\begin{lemma}
\label{lemma-pullback-injective-modules}
Soit $f : \mathcal{X} \to \mathcal{Y}$ un $1$-morphisme de catégories
fibrées en groupoïdes sur $(\Sch/S)_{fppf}$. Soit
$\tau \in \{Zar, \etale, smooth, syntomic, fppf\}$.
Le foncteur
$f^* : \textit{Mod}(\mathcal{Y}_\tau, \mathcal{O}_\mathcal{Y}) \to
\textit{Mod}(\mathcal{X}_\tau, \mathcal{O}_\mathcal{X})$
admet un adjoint à gauche
$f_! : \textit{Mod}(\mathcal{X}_\tau, \mathcal{O}_\mathcal{X}) \to
\textit{Mod}(\mathcal{Y}_\tau, \mathcal{O}_\mathcal{Y})$, qui
coïncide avec le foncteur $f_!$ du Lemme \ref{lemma-pullback-injective}
sur les faisceaux abéliens sous-jacents.
Si $f$ est fidèle et si $\mathcal{X}$ admet des égalisateurs, alors
\begin{enumerate}
\item $f_!$ est exact, et
\item $f^{-1}\mathcal{I}$ est injectif dans
$\textit{Mod}(\mathcal{X}_\tau, \mathcal{O}_\mathcal{X})$
pour $\mathcal{I}$ injectif dans
$\textit{Mod}(\mathcal{Y}_\tau, \mathcal{O}_\mathcal{X})$.
\end{enumerate}
\end{lemma}

\begin{proof}
Rappelons que $f$ est un foncteur continu et cocontinu de sites
et que $f^{-1}\mathcal{O}_\mathcal{Y} = \mathcal{O}_\mathcal{X}$. Ainsi
Modules on Sites, Lemma \ref{sites-modules-lemma-lower-shriek-modules}
implique que $f^*$ admet un adjoint à gauche $f_!^{Mod}$.
Soit $x$ un objet de $\mathcal{X}$ au-dessus du schéma $U$.
Alors $f$ induit une équivalence de sites annelés
$$
\mathcal{X}/x \longrightarrow \mathcal{Y}/f(x)
$$
puisque les deux membres sont équivalents à $(\Sch/U)_\tau$, voir
Lemma \ref{lemma-localizing-structure-sheaf}.
Modules on Sites, Remark \ref{sites-modules-remark-when-shriek-equal}
montre que $f_!$ coïncide avec le foncteur sur les faisceaux abéliens.

\medskip\noindent
Supposons maintenant que $\mathcal{X}$ admette des égalisateurs et que $f$ soit fidèle.
Le Lemme \ref{lemma-pullback-injective}
nous dit que $f_!$ est exact. Enfin,
Homology, Lemma \ref{homology-lemma-adjoint-preserve-injectives}
implique l'assertion sur les images réciproques de modules injectifs.
\end{proof}




\section{Le complexe de {\v C}ech}
\label{section-cech}

\noindent
Pour calculer la cohomologie d'un faisceau sur un champ algébrique, nous la comparons
à la cohomologie du faisceau restreint à des recouvrements du
champ algébrique donné.

\medskip\noindent
Dans toute cette section, la situation sera la suivante. On se donne
un $1$-morphisme de catégories fibrées en groupoïdes
\begin{equation}
\label{equation-covering}
\vcenter{
\xymatrix{
\mathcal{U} \ar[rr]_f \ar[rd]_q & &  \mathcal{X} \ar[ld]^p \\
& (\Sch/S)_{fppf}
}
}
\end{equation}
Nous considérerons $\mathcal{U}$ comme un « recouvrement » de $\mathcal{X}$.
Nous voulons donc considérer l'objet simplicial
$$
\xymatrix{
\mathcal{U} \times_\mathcal{X} \mathcal{U} \times_\mathcal{X} \mathcal{U}
\ar@<1ex>[r]
\ar@<0ex>[r]
\ar@<-1ex>[r]
&
\mathcal{U} \times_\mathcal{X} \mathcal{U}
\ar@<0.5ex>[r]
\ar@<-0.5ex>[r]
&
\mathcal{U}
}
$$
dans la catégorie des catégories fibrées en groupoïdes sur
$(\Sch/S)_{fppf}$. Cependant, comme il s'agit d'une $(2, 1)$-catégorie et
non d'une catégorie, précisons explicitement ce que nous entendons. Soit
$\mathcal{U}_n$ la catégorie dont les objets sont les
$(u_0, \ldots, u_n, x, \alpha_0, \ldots, \alpha_n)$
où $\alpha_i : f(u_i) \to x$ est un isomorphisme dans $\mathcal{X}$.
Nous notons $f_n : \mathcal{U}_n \to \mathcal{X}$ le $1$-morphisme
qui associe à $(u_0, \ldots, u_n, x, \alpha_0, \ldots, \alpha_n)$
l'objet $x$. Remarquons que $\mathcal{U}_0 = \mathcal{U}$ et $f_0 = f$.
À un morphisme $\varphi : [m] \to [n]$, nous associons le $1$-morphisme
$\mathcal{U}_\varphi : \mathcal{U}_n \longrightarrow \mathcal{U}_n$
donné sur les objets par
$$
(u_0, \ldots, u_n, x, \alpha_0, \ldots, \alpha_n)
\longmapsto
(u_{\varphi(0)}, \ldots, u_{\varphi(m)}, x,
\alpha_{\varphi(0)}, \ldots, \alpha_{\varphi(m)})
$$
Tous ces $1$-morphismes se composent strictement comme il convient
(sans qu'il faille de $2$-morphismes) et sont tous des $1$-morphismes
sur $\mathcal{X}$. Nous notons $\mathcal{U}_\bullet$ cet objet simplicial.
Si $\mathcal{F}$ est un préfaisceau d'ensembles sur $\mathcal{X}$, nous obtenons un
ensemble cosimplicial
$$
\xymatrix{
\Gamma(\mathcal{U}_0, f_0^{-1}\mathcal{F})
\ar@<0.5ex>[r]
\ar@<-0.5ex>[r]
&
\Gamma(\mathcal{U}_1, f_1^{-1}\mathcal{F})
\ar@<1ex>[r]
\ar@<0ex>[r]
\ar@<-1ex>[r]
&
\Gamma(\mathcal{U}_2, f_2^{-1}\mathcal{F})
}
$$
Ici, les flèches sont les morphismes d'image réciproque suivant les morphismes donnés de
l'objet simplicial.
Si $\mathcal{F}$ est un préfaisceau de groupes abéliens, on obtient un groupe
abélien cosimplicial.

\medskip\noindent
Soit $\mathcal{U} \to \mathcal{X}$ comme ci-dessus et soit $\mathcal{F}$
un préfaisceau abélien sur $\mathcal{X}$.
Le {\it complexe de {\v C}ech} associé à cette situation est noté
$\check{\mathcal{C}}^\bullet(\mathcal{U} \to \mathcal{X}, \mathcal{F})$.
C'est le complexe de cochaînes associé au groupe abélien cosimplicial
ci-dessus, voir
Simplicial, Section \ref{simplicial-section-dold-kan-cosimplicial}.
Ses termes sont
$$
\check{\mathcal{C}}^n(\mathcal{U} \to \mathcal{X}, \mathcal{F})
= \Gamma(\mathcal{U}_n, f_n^{-1}\mathcal{F}).
$$
Les différentielles sont les morphismes
$$
d^n = \sum\nolimits_{i = 0}^{n + 1} (-1)^i \delta^{n + 1}_i :
\Gamma(\mathcal{U}_n, f_n^{-1}\mathcal{F})
\longrightarrow
\Gamma(\mathcal{U}_{n + 1}, f_{n + 1}^{-1}\mathcal{F})
$$
où $\delta^{n + 1}_i$ correspond au morphisme
$[n] \to [n + 1]$ qui omet l'indice $i$. Remarquons que le morphisme
$\Gamma(\mathcal{X}, \mathcal{F}) \to
\Gamma(\mathcal{U}_0, f_0^{-1}\mathcal{F}_0)$
est dans le noyau de la différentielle $d^0$. Nous définissons donc
le {\it complexe de {\v C}ech augmenté} comme le complexe
$$
\ldots \to 0 \to
\Gamma(\mathcal{X}, \mathcal{F}) \to
\Gamma(\mathcal{U}_0, f_0^{-1}\mathcal{F}_0) \to
\Gamma(\mathcal{U}_1, f_1^{-1}\mathcal{F}_1) \to \ldots
$$
où $\Gamma(\mathcal{X}, \mathcal{F})$ est placé en degré $-1$.
Le complexe de {\v C}ech augmenté est acyclique si et seulement si le morphisme canonique
$$
\Gamma(\mathcal{X}, \mathcal{F})[0]
\longrightarrow
\check{\mathcal{C}}^\bullet(\mathcal{U} \to \mathcal{X}, \mathcal{F})
$$
est un quasi-isomorphisme de complexes.

\begin{lemma}
\label{lemma-generalities}
Généralités sur les complexes de {\v C}ech.
\begin{enumerate}
\item Si
$$
\xymatrix{
\mathcal{V} \ar[d]_g \ar[r]_h & \mathcal{U} \ar[d]^f \\
\mathcal{Y} \ar[r]^e & \mathcal{X}
}
$$
est un diagramme $2$-commutatif de catégories fibrées en groupoïdes sur
$(\Sch/S)_{fppf}$, alors il existe un morphisme de complexes de {\v C}ech
$$
\check{\mathcal{C}}^\bullet(\mathcal{U} \to \mathcal{X}, \mathcal{F})
\longrightarrow
\check{\mathcal{C}}^\bullet(\mathcal{V} \to \mathcal{Y}, e^{-1}\mathcal{F})
$$
\item si $h$ et $e$ sont des équivalences, le morphisme de (1) est un isomorphisme,
\item si $f, f' : \mathcal{U} \to \mathcal{X}$ sont $2$-isomorphes, alors
les complexes de {\v C}ech associés sont isomorphes.
\end{enumerate}
\end{lemma}

\begin{proof}
Dans la situation de (1), soit $t : f \circ h \to e \circ g$ un $2$-morphisme.
Le morphisme de complexes est donné en degré $n$ par
l'image réciproque suivant les $1$-morphismes
$\mathcal{V}_n \to \mathcal{U}_n$ définis par la règle
$$
(v_0, \ldots, v_n, y, \beta_0, \ldots, \beta_n)
\longmapsto
(h(v_0), \ldots, h(v_n), e(y),
e(\beta_0) \circ t_{v_0}, \ldots, e(\beta_n) \circ t_{v_n}).
$$
Pour (2), remarquons que l'image réciproque sur les sections globales est un isomorphisme
pour tout préfaisceau d'ensembles lorsqu'elle est prise suivant une équivalence
de catégories. L'assertion (3) résulte de la combinaison de (1) et (2).
\end{proof}

\begin{lemma}
\label{lemma-homotopy}
S'il existe un $1$-morphisme $s : \mathcal{X} \to \mathcal{U}$
tel que $f \circ s$ soit $2$-isomorphe à $\text{id}_\mathcal{X}$,
alors le complexe de {\v C}ech augmenté est homotope à zéro.
\end{lemma}

\begin{proof}
Posons $\mathcal{U}' = \mathcal{U} \times_\mathcal{X} \mathcal{X}$,
égal au produit fibré décrit dans
Categories, Lemma \ref{categories-lemma-2-product-categories-over-C}.
Posons $f' : \mathcal{U}' \to \mathcal{X}$ égal à la seconde projection.
Alors $\mathcal{U} \to \mathcal{U}'$, $u \mapsto (u, f(x), 1)$
est une équivalence sur $\mathcal{X}$ ; nous pouvons donc remplacer
$(\mathcal{U}, f)$ par $(\mathcal{U}', f')$ d'après le
Lemma \ref{lemma-generalities}.
L'avantage est que $f'$ admet maintenant une section $s'$ telle
que $f' \circ s' = \text{id}_\mathcal{X}$ strictement. En effet, si
$t : s \circ f \to \text{id}_\mathcal{X}$ est un $2$-isomorphisme,
nous pouvons poser $s'(x) = (s(x), x, t_x)$. Nous pouvons donc supposer que
$f \circ s = \text{id}_\mathcal{X}$.

\medskip\noindent
Lorsque $f \circ s = \text{id}_\mathcal{X}$, le résultat découle
de principes généraux. Explicitons l'homotopie. Pour
$n \geq 0$, définissons $s_n : \mathcal{U}_n \to \mathcal{U}_{n + 1}$
comme le $1$-morphisme donné sur les objets par la règle
$$
(u_0, \ldots, u_n, x, \alpha_0, \ldots, \alpha_n)
\longmapsto
(u_0, \ldots, u_n, s(x), x,
\alpha_0, \ldots, \alpha_n, \text{id}_x).
$$
Définissons
$$
h^{n + 1} :
\Gamma(\mathcal{U}_{n + 1}, f_{n + 1}^{-1}\mathcal{F})
\longrightarrow
\Gamma(\mathcal{U}_n, f_n^{-1}\mathcal{F})
$$
comme l'image réciproque suivant $s_n$. Posons aussi $s_{-1} = s$ et
$h^0 : \Gamma(\mathcal{U}_0, f_0^{-1}\mathcal{F}) \to
\Gamma(\mathcal{X}, \mathcal{F})$ égal à l'image réciproque suivant $s_{-1}$.
Alors la famille de morphismes $\{h^n\}_{n \geq 0}$ est une homotopie entre
$1$ et $0$ sur le complexe de {\v C}ech augmenté.
\end{proof}









\section{Le complexe de {\v C}ech relatif}
\label{section-sheaf-cech-complex}

\noindent
Soit $f : \mathcal{U} \to \mathcal{X}$ un $1$-morphisme de catégories
fibrées en groupoïdes sur $(\Sch/S)_{fppf}$ comme dans
(\ref{equation-covering}). Considérons l'objet simplicial
associé $\mathcal{U}_\bullet$ et les morphismes
$f_n : \mathcal{U}_n \to \mathcal{X}$. Soit
$\tau \in \{Zar, \etale, smooth, syntomic, fppf\}$.
Enfin, supposons que $\mathcal{F}$ soit un faisceau (d'ensembles)
sur $\mathcal{X}_\tau$. Alors
$$
\xymatrix{
f_{0, *}f_0^{-1}\mathcal{F}
\ar@<0.5ex>[r]
\ar@<-0.5ex>[r]
&
f_{1, *}f_1^{-1}\mathcal{F}
\ar@<1ex>[r]
\ar@<0ex>[r]
\ar@<-1ex>[r]
&
f_{2, *}f_2^{-1}\mathcal{F}
}
$$
est un faisceau cosimplicial sur $\mathcal{X}_\tau$, où nous employons les morphismes d'image réciproque
introduits dans
Sites, Section \ref{sites-section-pullback}.
Si $\mathcal{F}$ est un faisceau abélien, les $f_{n, *}f_n^{-1}\mathcal{F}$
forment un faisceau abélien cosimplicial sur $\mathcal{X}_\tau$.
Le complexe associé (voir
Simplicial, Section \ref{simplicial-section-dold-kan-cosimplicial})
$$
\ldots \to 0 \to
f_{0, *}f_0^{-1}\mathcal{F} \to
f_{1, *}f_1^{-1}\mathcal{F} \to
f_{2, *}f_2^{-1}\mathcal{F} \to \ldots
$$
est appelé le {\it complexe de {\v C}ech relatif} associé à la situation.
Nous noterons ce complexe $\mathcal{K}^\bullet(f, \mathcal{F})$.
Le {\it complexe de {\v C}ech relatif augmenté} est le complexe
$$
\ldots \to 0 \to
\mathcal{F} \to
f_{0, *}f_0^{-1}\mathcal{F} \to
f_{1, *}f_1^{-1}\mathcal{F} \to
f_{2, *}f_2^{-1}\mathcal{F} \to \ldots
$$
où $\mathcal{F}$ est en degré $-1$. Le complexe de {\v C}ech relatif augmenté
est acyclique si et seulement si le morphisme
$\mathcal{F}[0] \to \mathcal{K}^\bullet(f, \mathcal{F})$
est un quasi-isomorphisme de complexes de faisceaux.

\begin{remark}
\label{remark-cech-complex-presheaves}
Nous pouvons aussi définir le complexe $\mathcal{K}^\bullet(f, \mathcal{F})$
si $\mathcal{F}$ est un préfaisceau, mais nous ne pouvons alors employer
Sites, Section \ref{sites-section-pullback}
pour définir les morphismes d'image réciproque. Pour expliquer ceux-ci, supposons donné
un diagramme commutatif
$$
\xymatrix{
\mathcal{V} \ar[rd]_g \ar[rr]_h & &  \mathcal{U} \ar[ld]^f \\
& \mathcal{X}
}
$$
de catégories fibrées en groupoïdes sur $(\Sch/S)_{fppf}$
et un préfaisceau $\mathcal{G}$ sur $\mathcal{U}$.
Nous pouvons alors définir le morphisme d'image réciproque
$f_*\mathcal{G} \to g_*h^{-1}\mathcal{G}$ comme la composée
$$
f_*\mathcal{G} \longrightarrow
f_*h_*h^{-1}\mathcal{G} = g_*h^{-1}\mathcal{G}
$$
où le morphisme provient du morphisme d'adjonction
$\mathcal{G} \to h_*h^{-1}\mathcal{G}$. Cela fonctionne parce que, dans notre situation,
les foncteurs $h_*$ et $h^{-1}$ sont adjoints sur les préfaisceaux (et coïncident avec
leurs homologues sur les faisceaux). Voir les
Sections \ref{section-presheaves} et \ref{section-sheaves}.
\end{remark}

\begin{lemma}
\label{lemma-generalities-sheafified}
Généralités sur les complexes de {\v C}ech relatifs.
\begin{enumerate}
\item Si
$$
\xymatrix{
\mathcal{V} \ar[d]_g \ar[r]_h & \mathcal{U} \ar[d]^f \\
\mathcal{Y} \ar[r]^e & \mathcal{X}
}
$$
est un diagramme $2$-commutatif de catégories fibrées en groupoïdes sur
$(\Sch/S)_{fppf}$, alors il existe un morphisme
$e^{-1}\mathcal{K}^\bullet(f, \mathcal{F}) \to
\mathcal{K}^\bullet(g, e^{-1}\mathcal{F})$.
\item si $h$ et $e$ sont des équivalences, le morphisme de (1) est un isomorphisme,
\item si $f, f' : \mathcal{U} \to \mathcal{X}$ sont $2$-isomorphes, alors
les complexes de {\v C}ech relatifs associés sont isomorphes,
\end{enumerate}
\end{lemma}

\begin{proof}
Littéralement la même que la preuve du
Lemma \ref{lemma-generalities}
en employant les morphismes d'image réciproque de la
Remark \ref{remark-cech-complex-presheaves}.
\end{proof}

\begin{lemma}
\label{lemma-homotopy-sheafified}
S'il existe un $1$-morphisme $s : \mathcal{X} \to \mathcal{U}$
tel que $f \circ s$ soit $2$-isomorphe à $\text{id}_\mathcal{X}$,
alors le complexe de {\v C}ech relatif augmenté est homotope à zéro.
\end{lemma}

\begin{proof}
Littéralement la même que la preuve du
Lemma \ref{lemma-homotopy}.
\end{proof}

\begin{remark}
\label{remark-cech-complex-sections}
« Calculons » la valeur du complexe de {\v C}ech relatif en un
objet $x$ de $\mathcal{X}$. Écrivons $p(x) = U$.
Considérons le diagramme de $2$-produit fibré (qui sert à introduire
la notation $g : \mathcal{V} \to \mathcal{Y}$)
$$
\xymatrix{
\mathcal{V} \ar@{=}[r] \ar[d]_g &
(\Sch/U)_{fppf} \times_{x, \mathcal{X}} \mathcal{U} \ar[r] \ar[d] &
\mathcal{U} \ar[d]^f \\
\mathcal{Y} \ar@{=}[r] &
(\Sch/U)_{fppf} \ar[r]^-x & \mathcal{X}
}
$$
Remarquons que le morphisme $\mathcal{V}_n \to \mathcal{U}_n$ de la preuve du
Lemma \ref{lemma-generalities}
induit une équivalence
$\mathcal{V}_n =
(\Sch/U)_{fppf} \times_{x, \mathcal{X}} \mathcal{U}_n$.
Il résulte donc de
(\ref{equation-pushforward})
que
$$
\Gamma(x, \mathcal{K}^\bullet(f, \mathcal{F})) =
\check{\mathcal{C}}^\bullet(\mathcal{V} \to \mathcal{Y}, x^{-1}\mathcal{F})
$$
En d'autres termes : la valeur du complexe de {\v C}ech relatif en un objet $x$ de
$\mathcal{X}$ est le complexe de {\v C}ech du changement de base de $f$ à
$\mathcal{X}/x \cong (\Sch/U)_{fppf}$. Cela implique par exemple que le
Lemma \ref{lemma-homotopy}
implique le
Lemma \ref{lemma-homotopy-sheafified}
et, plus généralement, que les résultats sur le complexe de {\v C}ech (usuel) impliquent
les résultats correspondants pour le complexe de {\v C}ech relatif.
\end{remark}

\begin{lemma}
\label{lemma-base-change-cech-complex}
Soit
$$
\xymatrix{
\mathcal{V} \ar[d]_g \ar[r]_h & \mathcal{U} \ar[d]^f \\
\mathcal{Y} \ar[r]^e & \mathcal{X}
}
$$
soit un $2$-produit fibré de catégories fibrées en groupoïdes sur
$(\Sch/S)_{fppf}$ et soit $\mathcal{F}$ un préfaisceau abélien
sur $\mathcal{X}$. Alors le morphisme
$e^{-1}\mathcal{K}^\bullet(f, \mathcal{F}) \to
\mathcal{K}^\bullet(g, e^{-1}\mathcal{F})$
du
Lemma \ref{lemma-generalities-sheafified}
est un isomorphisme de complexes de préfaisceaux abéliens.
\end{lemma}

\begin{proof}
Soit $y$ un objet de $\mathcal{Y}$ au-dessus du schéma $T$.
Posons $x = e(y)$. Nous allons montrer que le morphisme induit un isomorphisme
sur les sections au-dessus de $y$. Remarquons que
$$
\Gamma(y, e^{-1}\mathcal{K}^\bullet(f, \mathcal{F})) =
\Gamma(x, \mathcal{K}^\bullet(f, \mathcal{F})) =
\check{\mathcal{C}}^\bullet(
(\Sch/T)_{fppf} \times_{x, \mathcal{X}} \mathcal{U} \to
(\Sch/T)_{fppf}, x^{-1}\mathcal{F})
$$
d'après la
Remarque \ref{remark-cech-complex-sections}. D'autre part,
$$
\Gamma(y, \mathcal{K}^\bullet(g, e^{-1}\mathcal{F})) =
\check{\mathcal{C}}^\bullet(
(\Sch/T)_{fppf} \times_{y, \mathcal{Y}} \mathcal{V} \to
(\Sch/T)_{fppf}, y^{-1}e^{-1}\mathcal{F})
$$
également d'après la
Remark \ref{remark-cech-complex-sections}.
Remarquons que $y^{-1}e^{-1}\mathcal{F} = x^{-1}\mathcal{F}$
et que, puisque le diagramme est $2$-cartésien, le $1$-morphisme
$$
(\Sch/T)_{fppf} \times_{y, \mathcal{Y}} \mathcal{V} \to
(\Sch/T)_{fppf} \times_{x, \mathcal{X}} \mathcal{U}
$$
est une équivalence. Le morphisme sur les sections au-dessus de $y$ est donc un
isomorphisme d'après le
Lemma \ref{lemma-generalities}.
\end{proof}

\noindent
L'exactitude peut se vérifier sur un « recouvrement ».

\begin{lemma}
\label{lemma-check-exactness-covering}
Soit $f : \mathcal{U} \to \mathcal{X}$ un $1$-morphisme de catégories fibrées
en groupoïdes sur $(\Sch/S)_{fppf}$. Soit
$\tau \in \{Zar, \etale, smooth, syntomic, fppf\}$.
Soit
$$
\mathcal{F} \to \mathcal{G} \to \mathcal{H}
$$
un complexe dans $\textit{Ab}(\mathcal{X}_\tau)$. Supposons que
\begin{enumerate}
\item pour tout objet $x$ de $\mathcal{X}$, il existe un recouvrement
$\{x_i \to x\}$ dans $\mathcal{X}_\tau$ tel que chaque $x_i$ soit isomorphe
à $f(u_i)$ pour un objet $u_i$ de $\mathcal{U}$, et
\item $f^{-1}\mathcal{F} \to f^{-1}\mathcal{G} \to f^{-1}\mathcal{H}$ soit exact.
\end{enumerate}
Alors la suite $\mathcal{F} \to \mathcal{G} \to \mathcal{H}$
est exacte.
\end{lemma}

\begin{proof}
Soit $x$ un objet de $\mathcal{X}$ au-dessus du schéma $T$.
Considérons la suite
$x^{-1}\mathcal{F} \to x^{-1}\mathcal{G} \to x^{-1}\mathcal{H}$
de faisceaux abéliens sur $(\Sch/T)_\tau$. Il suffit de montrer
que cette suite est exacte. Par hypothèse, il existe un $\tau$-recouvrement
$\{T_i \to T\}$ tel que $x|_{T_i}$ soit isomorphe à $f(u_i)$ pour
un objet $u_i$ de $\mathcal{U}$ au-dessus de $T_i$ et que, de plus, la suite
$u_i^{-1}f^{-1}\mathcal{F} \to u_i^{-1}f^{-1}\mathcal{G} \to
u_i^{-1}f^{-1}\mathcal{H}$ de faisceaux abéliens sur $(\Sch/T_i)_\tau$
soit exacte. Comme
$u_i^{-1}f^{-1}\mathcal{F} = x^{-1}\mathcal{F}|_{(\Sch/T_i)_\tau}$
nous concluons que la suite
$x^{-1}\mathcal{F} \to x^{-1}\mathcal{G} \to x^{-1}\mathcal{H}$
devient exacte après localisation en chacun des membres d'un recouvrement ;
la suite est donc exacte.
\end{proof}

\begin{proposition}
\label{proposition-exactness-cech-complex}
Soit $f : \mathcal{U} \to \mathcal{X}$ un $1$-morphisme de catégories fibrées
en groupoïdes sur $(\Sch/S)_{fppf}$. Soit
$\tau \in \{Zar, \etale, smooth, syntomic, fppf\}$.
Si
\begin{enumerate}
\item $\mathcal{F}$ est un faisceau abélien sur $\mathcal{X}_\tau$, et
\item pour tout objet $x$ de $\mathcal{X}$, il existe un recouvrement
$\{x_i \to x\}$ dans $\mathcal{X}_\tau$ tel que chaque $x_i$ soit isomorphe
à $f(u_i)$ pour un objet $u_i$ de $\mathcal{U}$,
\end{enumerate}
alors le complexe de {\v C}ech relatif augmenté
$$
\ldots \to 0 \to
\mathcal{F} \to
f_{0, *}f_0^{-1}\mathcal{F} \to
f_{1, *}f_1^{-1}\mathcal{F} \to
f_{2, *}f_2^{-1}\mathcal{F} \to \ldots
$$
est exact dans $\textit{Ab}(\mathcal{X}_\tau)$.
\end{proposition}

\begin{proof}
D'après le
Lemma \ref{lemma-check-exactness-covering}
il suffit de vérifier l'exactitude après image réciproque sur $\mathcal{U}$.
D'après le
Lemma \ref{lemma-base-change-cech-complex}
L'image réciproque du complexe de {\v C}ech relatif augmenté est isomorphe
au complexe de {\v C}ech relatif augmenté du morphisme
$\mathcal{U} \times_\mathcal{X} \mathcal{U} \to \mathcal{U}$
et d'un faisceau abélien sur $\mathcal{U}_\tau$. Comme il existe une section
$\Delta_{\mathcal{U}/\mathcal{X}} : \mathcal{U} \to
\mathcal{U} \times_\mathcal{X} \mathcal{U}$, l'exactitude résulte du
Lemme \ref{lemma-homotopy-sheafified}.
\end{proof}

\noindent
Nous pouvons ainsi construire la suite spectrale de {\v C}ech vers la cohomologie
comme suit. Nous donnons d'abord une version technique précise. Dans la section suivante,
nous donnerons une version qui ne concerne que les champs algébriques.

\begin{lemma}
\label{lemma-cech-to-cohomology}
Soit $f : \mathcal{U} \to \mathcal{X}$ un $1$-morphisme de catégories fibrées
en groupoïdes sur $(\Sch/S)_{fppf}$. Soit
$\tau \in \{Zar, \etale, smooth, syntomic, fppf\}$.
Supposons que
\begin{enumerate}
\item $\mathcal{F}$ soit un faisceau abélien sur $\mathcal{X}_\tau$,
\item pour tout objet $x$ de $\mathcal{X}$, il existe un recouvrement
$\{x_i \to x\}$ dans $\mathcal{X}_\tau$ tel que chaque $x_i$ soit isomorphe
à $f(u_i)$ pour un objet $u_i$ de $\mathcal{U}$,
\item la catégorie $\mathcal{U}$ admette des égalisateurs, et
\item le foncteur $f$ soit fidèle.
\end{enumerate}
Alors il existe une suite spectrale de groupes abéliens du premier quadrant
$$
E_1^{p, q} = H^q((\mathcal{U}_p)_\tau, f_p^{-1}\mathcal{F})
\Rightarrow
H^{p + q}(\mathcal{X}_\tau, \mathcal{F})
$$
qui converge vers la cohomologie de $\mathcal{F}$ pour la topologie $\tau$.
\end{lemma}

\begin{proof}
Avant de commencer la preuve, faisons quelques remarques. D'après le
Lemma \ref{lemma-fibre-products-preserve-properties}
(et par récurrence), toutes les catégories fibrées en groupoïdes $\mathcal{U}_p$
admettent des égalisateurs et tous les morphismes $f_p : \mathcal{U}_p \to \mathcal{X}$
sont fidèles. Soit $\mathcal{I}$ un objet injectif
de $\textit{Ab}(\mathcal{X}_\tau)$. D'après le
Lemma \ref{lemma-pullback-injective}
nous voyons que $f_p^{-1}\mathcal{I}$ est un objet injectif de
$\textit{Ab}((\mathcal{U}_p)_\tau)$.
Ainsi $f_{p, *}f_p^{-1}\mathcal{I}$ est un objet injectif de
$\textit{Ab}(\mathcal{X}_\tau)$ d'après le
Lemma \ref{lemma-pushforward-injective}.
Par conséquent, la
Proposition \ref{proposition-exactness-cech-complex}
montre que le complexe de {\v C}ech relatif augmenté
$$
\ldots \to 0 \to
\mathcal{I} \to
f_{0, *}f_0^{-1}\mathcal{I} \to
f_{1, *}f_1^{-1}\mathcal{I} \to
f_{2, *}f_2^{-1}\mathcal{I} \to \ldots
$$
est un complexe exact dans $\textit{Ab}(\mathcal{X}_\tau)$ dont tous les
termes sont injectifs. La prise des sections globales de ce complexe est exacte,
et nous voyons que le complexe de {\v C}ech
$\check{\mathcal{C}}^\bullet(\mathcal{U} \to \mathcal{X}, \mathcal{I})$
est quasi-isomorphe à $\Gamma(\mathcal{X}_\tau, \mathcal{I})[0]$.

\medskip\noindent
Après ces préliminaires, considérons les deux suites spectrales
associées au complexe double (voir
Homology, Section \ref{homology-section-double-complex})
$$
\check{\mathcal{C}}^\bullet(\mathcal{U} \to \mathcal{X}, \mathcal{I}^\bullet)
$$
où $\mathcal{F} \to \mathcal{I}^\bullet$ est une résolution injective
dans $\textit{Ab}(\mathcal{X}_\tau)$.
La discussion ci-dessus montre que
Homology, Lemma \ref{homology-lemma-double-complex-gives-resolution}
s'applique et montre que
$\Gamma(\mathcal{X}_\tau, \mathcal{I}^\bullet)$
est quasi-isomorphe au complexe total associé au complexe double.
D'après les remarques ci-dessus, le complexe $f_p^{-1}\mathcal{I}^\bullet$ est une
résolution injective de $f_p^{-1}\mathcal{F}$. L'autre suite
spectrale est donc celle indiquée dans le lemme.
\end{proof}

\noindent
Il existe bien sûr aussi une version pour les modules.

\begin{lemma}
\label{lemma-cech-to-cohomology-modules}
Soit $f : \mathcal{U} \to \mathcal{X}$ un $1$-morphisme de catégories fibrées
en groupoïdes sur $(\Sch/S)_{fppf}$. Soit
$\tau \in \{Zar, \etale, smooth, syntomic, fppf\}$.
Supposons que
\begin{enumerate}
\item $\mathcal{F}$ soit un objet de
$\textit{Mod}(\mathcal{X}_\tau, \mathcal{O}_\mathcal{X})$,
\item pour tout objet $x$ de $\mathcal{X}$, il existe un recouvrement
$\{x_i \to x\}$ dans $\mathcal{X}_\tau$ tel que chaque $x_i$ soit isomorphe
à $f(u_i)$ pour un objet $u_i$ de $\mathcal{U}$,
\item la catégorie $\mathcal{U}$ admette des égalisateurs, et
\item le foncteur $f$ soit fidèle.
\end{enumerate}
Alors il existe une suite spectrale du premier quadrant de
$\Gamma(\mathcal{O}_\mathcal{X})$-modules
$$
E_1^{p, q} = H^q((\mathcal{U}_p)_\tau, f_p^*\mathcal{F})
\Rightarrow
H^{p + q}(\mathcal{X}_\tau, \mathcal{F})
$$
qui converge vers la cohomologie de $\mathcal{F}$ pour la topologie $\tau$.
\end{lemma}

\begin{proof}
La preuve de ce lemme est identique à celle du
Lemme \ref{lemma-cech-to-cohomology}
si ce n'est qu'elle emploie une résolution injective dans
$\textit{Mod}(\mathcal{X}_\tau, \mathcal{O}_\mathcal{X})$
et qu'elle emploie le
Lemma \ref{lemma-pullback-injective-modules}
au lieu du
Lemma \ref{lemma-pullback-injective}.
\end{proof}

\noindent
Voici un lemme qui ramène un type plus usuel de recouvrement aux
types de recouvrements rencontrés ci-dessus.

\begin{lemma}
\label{lemma-surjective-flat-locally-finite-presentation}
Soit $f : \mathcal{X} \to \mathcal{Y}$ un $1$-morphisme de
catégories fibrées en groupoïdes sur $(\Sch/S)_{fppf}$.
\begin{enumerate}
\item Supposons $f$ représentable par des espaces algébriques, surjectif,
plat et localement de présentation finie. Alors, pour tout objet $y$ de
$\mathcal{Y}$, il existe un recouvrement fppf $\{y_i \to y\}$ et des objets
$x_i$ de $\mathcal{X}$ tels que $f(x_i) \cong y_i$ dans $\mathcal{Y}$.
\item Supposons $f$ représentable par des espaces algébriques, surjectif
et lisse. Alors, pour tout objet $y$ de
$\mathcal{Y}$, il existe un recouvrement étale $\{y_i \to y\}$ et des objets
$x_i$ de $\mathcal{X}$ tels que $f(x_i) \cong y_i$ dans $\mathcal{Y}$.
\end{enumerate}
\end{lemma}

\begin{proof}
Preuve de (1). Supposons que $y$ soit au-dessus du schéma $V$.
Nous pouvons voir $y$ comme un morphisme $(\Sch/V)_{fppf} \to \mathcal{Y}$.
Par définition, le $2$-produit fibré
$\mathcal{X} \times_\mathcal{Y} (\Sch/V)_{fppf}$
est représenté par un espace algébrique $W$, et le morphisme
$W \to V$ est surjectif, plat et localement de présentation finie.
Choisissons un schéma $U$ et un morphisme étale surjectif $U \to W$.
Alors $U \to V$ est aussi surjectif, plat et localement de présentation finie
(voir Morphisms of Spaces, Lemmes
\ref{spaces-morphisms-lemma-etale-flat},
\ref{spaces-morphisms-lemma-etale-locally-finite-presentation},
\ref{spaces-morphisms-lemma-composition-surjective},
\ref{spaces-morphisms-lemma-composition-finite-presentation}, and
\ref{spaces-morphisms-lemma-composition-flat}).
Ainsi $\{U \to V\}$ est un recouvrement fppf. Notons $x$ l'objet de
$\mathcal{X}$ au-dessus de $U$ correspondant au $1$-morphisme
$(\Sch/U)_{fppf} \to \mathcal{X}$. Alors $\{f(x) \to y\}$ est
le recouvrement fppf cherché de $\mathcal{Y}$.

\medskip\noindent
Preuve de (2). Supposons que $y$ soit au-dessus du schéma $V$.
Nous pouvons voir $y$ comme un morphisme $(\Sch/V)_{fppf} \to \mathcal{Y}$.
Par définition, le $2$-produit fibré
$\mathcal{X} \times_\mathcal{Y} (\Sch/V)_{fppf}$
est représenté par un espace algébrique $W$, et le morphisme
$W \to V$ est surjectif et lisse.
Choisissons un schéma $U$ et un morphisme étale surjectif $U \to W$.
Alors $U \to V$ est aussi surjectif et lisse
(voir Morphisms of Spaces, Lemmes
\ref{spaces-morphisms-lemma-etale-smooth},
\ref{spaces-morphisms-lemma-composition-surjective}, and
\ref{spaces-morphisms-lemma-composition-smooth}).
Ainsi $\{U \to V\}$ est un recouvrement lisse. D'après
More on Morphisms, Lemma \ref{more-morphisms-lemma-etale-dominates-smooth}
il existe un recouvrement étale $\{V_i \to V\}$ tel que
chaque $V_i \to V$ se factorise par $U$. Notons $x_i$ l'objet de
$\mathcal{X}$ au-dessus de $V_i$ correspondant au $1$-morphisme
$$
(\Sch/V_i)_{fppf} \to (\Sch/U)_{fppf} \to \mathcal{X}.
$$
Alors $\{f(x_i) \to y\}$ est
le recouvrement étale cherché de $\mathcal{Y}$.
\end{proof}

\begin{lemma}
\label{lemma-cech-to-cohomology-relative}
Soient $f : \mathcal{U} \to \mathcal{X}$ et
$g : \mathcal{X} \to \mathcal{Y}$
des $1$-morphismes composables de catégories fibrées
en groupoïdes sur $(\Sch/S)_{fppf}$. Soit
$\tau \in \{Zar, \etale, smooth, syntomic, \linebreak[0] fppf\}$.
Supposons que
\begin{enumerate}
\item $\mathcal{F}$ soit un faisceau abélien sur $\mathcal{X}_\tau$,
\item pour tout objet $x$ de $\mathcal{X}$, il existe un recouvrement
$\{x_i \to x\}$ dans $\mathcal{X}_\tau$ tel que chaque $x_i$ soit isomorphe
à $f(u_i)$ pour un objet $u_i$ de $\mathcal{U}$,
\item la catégorie $\mathcal{U}$ admette des égalisateurs, et
\item le foncteur $f$ soit fidèle.
\end{enumerate}
Alors il existe une suite spectrale du premier quadrant de faisceaux abéliens
sur $\mathcal{Y}_\tau$
$$
E_1^{p, q} = R^q(g \circ f_p)_*f_p^{-1}\mathcal{F}
\Rightarrow
R^{p + q}g_*\mathcal{F}
$$
où toutes les images directes supérieures sont calculées pour la topologie $\tau$.
\end{lemma}

\begin{proof}
Remarquons que les hypothèses sur $f : \mathcal{U} \to \mathcal{X}$
et $\mathcal{F}$ sont identiques à celles du
Lemma \ref{lemma-cech-to-cohomology}.
Les remarques préliminaires faites dans la preuve de ce lemme
s'appliquent donc ici aussi. Elles impliquent en particulier que
$$
0 \to g_*\mathcal{I} \to
(g \circ f_0)_*f_0^{-1}\mathcal{I} \to
(g \circ f_1)_*f_1^{-1}\mathcal{I} \to \ldots
$$
est exacte si $\mathcal{I}$ est un objet injectif de
$\textit{Ab}(\mathcal{X}_\tau)$.
Cela étant, considérons les deux suites spectrales de
Homology, Section \ref{homology-section-double-complex}
associées au complexe double $\mathcal{C}^{\bullet, \bullet}$ de termes
$$
\mathcal{C}^{p, q} = (g \circ f_p)_*\mathcal{I}^q
$$
où $\mathcal{F} \to \mathcal{I}^\bullet$ est une résolution injective
dans $\textit{Ab}(\mathcal{X}_\tau)$. La première suite spectrale implique, par
Homology, Lemma \ref{homology-lemma-double-complex-gives-resolution},
que $g_*\mathcal{I}^\bullet$ est quasi-isomorphe au complexe total
associé à $\mathcal{C}^{\bullet, \bullet}$.
Comme $f_p^{-1}\mathcal{I}^\bullet$ est une résolution injective de
$f_p^{-1}\mathcal{F}$ (see
Lemma \ref{lemma-pullback-injective})
la seconde suite spectrale a pour termes
$E_1^{p, q} = R^q(g \circ f_p)_*f_p^{-1}\mathcal{F}$, comme dans l'énoncé
du lemme.
\end{proof}

\begin{lemma}
\label{lemma-cech-to-cohomology-relative-modules}
Soient $f : \mathcal{U} \to \mathcal{X}$ et
$g : \mathcal{X} \to \mathcal{Y}$
des $1$-morphismes composables de catégories fibrées
en groupoïdes sur $(\Sch/S)_{fppf}$. Soit
$\tau \in \{Zar, \etale, smooth, syntomic, \linebreak[0] fppf\}$.
Supposons que
\begin{enumerate}
\item $\mathcal{F}$ soit un objet de
$\textit{Mod}(\mathcal{X}_\tau, \mathcal{O}_\mathcal{X})$,
\item pour tout objet $x$ de $\mathcal{X}$, il existe un recouvrement
$\{x_i \to x\}$ dans $\mathcal{X}_\tau$ tel que chaque $x_i$ soit isomorphe
à $f(u_i)$ pour un objet $u_i$ de $\mathcal{U}$,
\item la catégorie $\mathcal{U}$ admette des égalisateurs, et
\item le foncteur $f$ soit fidèle.
\end{enumerate}
Alors il existe une suite spectrale du premier quadrant dans
$\textit{Mod}(\mathcal{Y}_\tau, \mathcal{O}_\mathcal{Y})$
$$
E_1^{p, q} = R^q(g \circ f_p)_*f_p^{-1}\mathcal{F}
\Rightarrow
R^{p + q}g_*\mathcal{F}
$$
où toutes les images directes supérieures sont calculées pour la topologie $\tau$.
\end{lemma}

\begin{proof}
La preuve est identique à celle du
Lemma \ref{lemma-cech-to-cohomology-relative}
si ce n'est qu'elle emploie une résolution injective dans
$\textit{Mod}(\mathcal{X}_\tau, \mathcal{O}_\mathcal{X})$
et qu'elle emploie le
Lemma \ref{lemma-pullback-injective-modules}
au lieu du
Lemma \ref{lemma-pullback-injective}.
\end{proof}






\section{Cohomologie des champs algébriques}
\label{section-cohomology}

\noindent
Soit $\mathcal{X}$ un champ algébrique sur $S$. Dans les sections précédentes,
nous avons vu comment définir les faisceaux pour les topologies étale, ..., fppf
sur $\mathcal{X}$. En fait, nous avons construit un site
$\mathcal{X}_\tau$ pour chaque $\tau \in \{Zar, \etale, smooth, syntomic,
fppf\}$. Sur ces sites, on dispose de la notion de faisceau abélien $\mathcal{F}$.
Dans le chapitre sur la cohomologie des sites, nous avons expliqué comment définir la
cohomologie. En réunissant ces constructions, définissons les
{\it sections globales dérivées}, ou la {\it cohomologie totale},
$$
R\Gamma_{Zar}(\mathcal{X}, \mathcal{F}),
R\Gamma_\etale(\mathcal{X}, \mathcal{F}), \ldots,
R\Gamma_{fppf}(\mathcal{X}, \mathcal{F})
$$
comme $\Gamma(\mathcal{X}_\tau, \mathcal{I}^\bullet)$, où
$\mathcal{F} \to \mathcal{I}^\bullet$ est une résolution injective
dans $\textit{Ab}(\mathcal{X}_\tau)$. Le $i$-ème groupe de cohomologie de $\mathcal{F}$
est la $i$-ème cohomologie de la cohomologie totale. Nous le noterons
ainsi :
$$
H^i_{Zar}(\mathcal{X}, \mathcal{F}),
H^i_\etale(\mathcal{X}, \mathcal{F}), \ldots,
H^i_{fppf}(\mathcal{X}, \mathcal{F}).
$$
Il se trouvera que $H^i_\etale = H^i_{smooth}$ en vertu de
More on Morphisms, Lemma \ref{more-morphisms-lemma-etale-dominates-smooth}.

\medskip\noindent
Si $\mathcal{F}$ est un préfaisceau de $\mathcal{O}_\mathcal{X}$-modules
qui est un faisceau pour la topologie $\tau$, nous employons des résolutions
injectives dans $\textit{Mod}(\mathcal{X}_\tau, \mathcal{O}_\mathcal{X})$
pour calculer sa cohomologie totale, resp.\ ses groupes de cohomologie ;
le résultat est quasi-isomorphe, resp.\ isomorphe à la cohomologie de
$\mathcal{F}$ considéré comme faisceau de groupes abéliens, d'après le très général
Cohomology on Sites, Lemma
\ref{sites-cohomology-lemma-cohomology-modules-abelian-agree}.

\medskip\noindent
Jusqu'ici, notre seul outil de calcul des groupes de cohomologie est le résultat sur les
complexes de {\v C}ech démontré ci-dessus. Reformulons-le ici dans le langage
des champs algébriques pour les topologies étale et fppf. Soit
$f : \mathcal{U} \to \mathcal{X}$ un $1$-morphisme de champs algébriques.
Rappelons que
$$
f_p : \mathcal{U}_p =
\mathcal{U} \times_\mathcal{X} \ldots \times_\mathcal{X} \mathcal{U}
\longrightarrow
\mathcal{X}
$$
est le morphisme structural, avec $(p + 1)$ facteurs. Rappelons aussi
qu'un faisceau sur $\mathcal{X}$ est un faisceau pour la topologie fppf. Remarquons
que, si $\mathcal{U}$ est un espace algébrique, alors
$f : \mathcal{U} \to \mathcal{X}$ est représentable par des espaces algébriques,
voir
Algebraic Stacks, Lemma \ref{algebraic-lemma-representable-diagonal}.
La proposition s'applique donc en particulier à un recouvrement lisse du
champ algébrique $\mathcal{X}$ par un schéma.

\begin{proposition}
\label{proposition-smooth-covering-compute-cohomology}
\begin{slogan}
La cohomologie d'un champ se calcule sur le nerf de Čech d'une présentation
\end{slogan}
Soit $f : \mathcal{U} \to \mathcal{X}$ un $1$-morphisme de champs algébriques.
\begin{enumerate}
\item Soit $\mathcal{F}$ un faisceau abélien étale sur $\mathcal{X}$.
Supposons $f$ représentable par des espaces algébriques, surjectif et lisse.
Alors il existe une suite spectrale
$$
E_1^{p, q} = H^q_\etale(\mathcal{U}_p, f_p^{-1}\mathcal{F})
\Rightarrow
H^{p + q}_\etale(\mathcal{X}, \mathcal{F})
$$
\item Soit $\mathcal{F}$ un faisceau abélien sur $\mathcal{X}$.
Supposons $f$ représentable par des espaces algébriques, surjectif, plat
et localement de présentation finie. Alors il existe
une suite spectrale
$$
E_1^{p, q} = H^q_{fppf}(\mathcal{U}_p, f_p^{-1}\mathcal{F})
\Rightarrow
H^{p + q}_{fppf}(\mathcal{X}, \mathcal{F})
$$
\end{enumerate}
\end{proposition}

\begin{proof}
Pour le voir, vérifions les hypothèses (1) -- (4) du
Lemma \ref{lemma-cech-to-cohomology}.
Le $1$-morphisme $f$ est fidèle d'après
Algebraic Stacks, Lemma
\ref{algebraic-lemma-characterize-representable-by-algebraic-spaces}.
Cela prouve (4).
L'hypothèse (3) résulte du fait que $\mathcal{U}$ est un champ
algébrique, voir
Lemma \ref{lemma-fibre-products}.
Pour (2), appliquons le
Lemma \ref{lemma-surjective-flat-locally-finite-presentation}.
La condition (1) est satisfaite par hypothèse.
\end{proof}










\section{Images directes supérieures et champs algébriques}
\label{section-higher-direct-images}

\noindent
Soit $g : \mathcal{X} \to \mathcal{Y}$ un $1$-morphisme de champs algébriques
sur $S$. Dans les sections précédentes, nous avons construit un morphisme de topos
annelés $g : \Sh(\mathcal{X}_\tau) \to \Sh(\mathcal{Y}_\tau)$
pour chaque $\tau \in \{Zar, \etale, smooth, syntomic, fppf\}$.
Dans le chapitre sur la cohomologie des sites, nous avons expliqué comment
définir les images directes supérieures. Ainsi l'{\it image directe totale}
$Rg_*\mathcal{F}$ est définie comme $g_*\mathcal{I}^\bullet$, où
$\mathcal{F} \to \mathcal{I}^\bullet$ est une résolution injective dans
$\textit{Ab}(\mathcal{X}_\tau)$. La $i$-ème image directe supérieure
$R^ig_*\mathcal{F}$ est la $i$-ème cohomologie de l'image directe totale.
Important : le choix de la topologie $\tau$ importe ici !

\medskip\noindent
Si $\mathcal{F}$ est un préfaisceau de $\mathcal{O}_\mathcal{X}$-modules
qui est un faisceau pour la topologie $\tau$, nous employons des résolutions
injectives dans $\textit{Mod}(\mathcal{X}_\tau, \mathcal{O}_\mathcal{X})$
pour calculer l'image directe totale et les images directes supérieures.

\medskip\noindent
Jusqu'ici, notre seul outil de calcul des images directes supérieures de $g_*$
est le résultat sur les complexes de {\v C}ech démontré ci-dessus. Il requiert le choix
d'un « recouvrement » $f : \mathcal{U} \to \mathcal{X}$. Si $\mathcal{U}$
est un espace algébrique, alors $f : \mathcal{U} \to \mathcal{X}$
est représentable par des espaces algébriques, voir
Algebraic Stacks, Lemma \ref{algebraic-lemma-representable-diagonal}.
La proposition s'applique donc en particulier à un recouvrement lisse du
champ algébrique $\mathcal{X}$ par un schéma.

\begin{proposition}
\label{proposition-smooth-covering-compute-direct-image}
Soient $f : \mathcal{U} \to \mathcal{X}$ et $g : \mathcal{X} \to \mathcal{Y}$
des $1$-morphismes composables de champs algébriques.
\begin{enumerate}
\item Supposons $f$ représentable par des espaces algébriques, surjectif et
lisse.
\begin{enumerate}
\item Si $\mathcal{F}$ appartient à $\textit{Ab}(\mathcal{X}_\etale)$,
il existe une suite spectrale
$$
E_1^{p, q} = R^q(g \circ f_p)_*f_p^{-1}\mathcal{F}
\Rightarrow
R^{p + q}g_*\mathcal{F}
$$
dans $\textit{Ab}(\mathcal{Y}_\etale)$, les images directes supérieures étant
calculées pour la topologie étale.
\item Si $\mathcal{F}$ appartient à
$\textit{Mod}(\mathcal{X}_\etale, \mathcal{O}_\mathcal{X})$, alors
il existe une suite spectrale
$$
E_1^{p, q} = R^q(g \circ f_p)_*f_p^{-1}\mathcal{F}
\Rightarrow
R^{p + q}g_*\mathcal{F}
$$
dans $\textit{Mod}(\mathcal{Y}_\etale, \mathcal{O}_\mathcal{Y})$.
\end{enumerate}
\item Supposons $f$ représentable par des espaces algébriques, surjectif,
plat et localement de présentation finie.
\begin{enumerate}
\item Si $\mathcal{F}$ appartient à $\textit{Ab}(\mathcal{X})$, il existe
une suite spectrale
$$
E_1^{p, q} = R^q(g \circ f_p)_*f_p^{-1}\mathcal{F}
\Rightarrow
R^{p + q}g_*\mathcal{F}
$$
dans $\textit{Ab}(\mathcal{Y})$, les images directes supérieures étant
calculées pour la topologie fppf.
\item Si $\mathcal{F}$ appartient à $\textit{Mod}(\mathcal{O}_\mathcal{X})$, alors
il existe une suite spectrale
$$
E_1^{p, q} = R^q(g \circ f_p)_*f_p^{-1}\mathcal{F}
\Rightarrow
R^{p + q}g_*\mathcal{F}
$$
dans $\textit{Mod}(\mathcal{O}_\mathcal{Y})$.
\end{enumerate}
\end{enumerate}
\end{proposition}

\begin{proof}
Pour le voir, vérifions les hypothèses (1) -- (4) du
Lemme \ref{lemma-cech-to-cohomology-relative} et du
Lemme \ref{lemma-cech-to-cohomology-relative-modules}.
Le $1$-morphisme $f$ est fidèle d'après
Algebraic Stacks, Lemma
\ref{algebraic-lemma-characterize-representable-by-algebraic-spaces}.
Cela prouve (4).
L'hypothèse (3) résulte du fait que $\mathcal{U}$ est un champ
algébrique, voir
Lemma \ref{lemma-fibre-products}.
Pour (2), appliquons le
Lemma \ref{lemma-surjective-flat-locally-finite-presentation}.
La condition (1) est satisfaite par hypothèse dans les quatre cas.
\end{proof}

\noindent
Voici une description des images directes supérieures pour un
morphisme de champs algébriques.

\begin{lemma}
\label{lemma-pushforward-restriction}
Soit $S$ un schéma. Soit $f : \mathcal{X} \to \mathcal{Y}$ un
$1$-morphisme de champs algébriques\footnote{Ce résultat devrait valoir
pour tout $1$-morphisme de catégories fibrées en groupoïdes sur
$(\Sch/S)_{fppf}$.} sur $S$.
Soit $\tau \in \{Zariski,\linebreak[0] \etale,\linebreak[0]
smooth,\linebreak[0] syntomic,\linebreak[0] fppf\}$.
Soit $\mathcal{F}$
un objet de $\textit{Ab}(\mathcal{X}_\tau)$ ou de
$\textit{Mod}(\mathcal{X}_\tau, \mathcal{O}_\mathcal{X})$.
Alors le faisceau $R^if_*\mathcal{F}$ est le faisceau associé au
préfaisceau
$$
y \longmapsto
H^i_\tau\Big((\Sch/V)_{fppf} \times_{y, \mathcal{Y}} \mathcal{X},
\ \text{pr}^{-1}\mathcal{F}\Big)
$$
Ici, $y$ est un objet de $\mathcal{Y}$ au-dessus du schéma $V$.
\end{lemma}

\begin{proof}
Choisissons une résolution injective $\mathcal{F}[0] \to \mathcal{I}^\bullet$.
D'après la formule d'image directe (\ref{equation-pushforward}),
$R^if_*\mathcal{F}$ est le faisceau associé au préfaisceau qui associe
à $y$ la cohomologie du complexe
$$
\begin{matrix}
\Gamma\Big((\Sch/V)_{fppf} \times_{y, \mathcal{Y}} \mathcal{X},
\ \text{pr}^{-1}\mathcal{I}^{i - 1}\Big) \\
\downarrow \\
\Gamma\Big((\Sch/V)_{fppf} \times_{y, \mathcal{Y}} \mathcal{X},
\ \text{pr}^{-1}\mathcal{I}^i\Big) \\
\downarrow \\
\Gamma\Big((\Sch/V)_{fppf} \times_{y, \mathcal{Y}} \mathcal{X},
\ \text{pr}^{-1}\mathcal{I}^{i + 1}\Big)
\end{matrix}
$$
Comme $\text{pr}^{-1}$ est exact, il suffit de montrer que
$\text{pr}^{-1}$ préserve les injectifs. Cela résulte des
Lemmes \ref{lemma-pullback-injective} et
\ref{lemma-pullback-injective-modules}
ainsi que du fait que $\text{pr}$ est un morphisme représentable de
champs algébriques (de sorte que $\text{pr}$ est fidèle d'après
Algebraic Stacks, Lemma
\ref{algebraic-lemma-characterize-representable-by-algebraic-spaces}
et que
$(\Sch/V)_{fppf} \times_{y, \mathcal{Y}} \mathcal{X}$
admet des égalisateurs d'après le
Lemma \ref{lemma-fibre-products}).
\end{proof}

\noindent
Voici un résultat immédiat de changement de base.

\begin{lemma}
\label{lemma-base-change-higher-direct-images}
Soit $S$ un schéma. Soit
$\tau \in \{Zariski,\linebreak[0] \etale,\linebreak[0]
smooth,\linebreak[0] syntomic,\linebreak[0] fppf\}$. Soit
$$
\xymatrix{
\mathcal{Y}' \times_\mathcal{Y} \mathcal{X} \ar[r]_{g'} \ar[d]_{f'} &
\mathcal{X} \ar[d]^f \\
\mathcal{Y}' \ar[r]^g & \mathcal{Y}
}
$$
un diagramme $2$-cartésien de champs algébriques sur $S$. Alors le morphisme de
changement de base est un isomorphisme
$$
g^{-1}Rf_*\mathcal{F} \longrightarrow Rf'_*(g')^{-1}\mathcal{F}
$$
fonctoriel en $\mathcal{F}$ dans $\textit{Ab}(\mathcal{X}_\tau)$
ou en $\mathcal{F}$ dans $\textit{Mod}(\mathcal{X}_\tau, \mathcal{O}_\mathcal{X})$.
\end{lemma}

\begin{proof}
L'isomorphisme $g^{-1}f_*\mathcal{F} = f'_*(g')^{-1}\mathcal{F}$ est le
Lemme \ref{lemma-base-change} (et vaut pour des préfaisceaux arbitraires).
Pour les images directes totales, il existe un morphisme de changement de base parce que les
morphismes $g$ et $g'$ sont plats, voir
Cohomology on Sites, Section \ref{sites-cohomology-section-base-change-map}.
Pour voir que ce morphisme est un quasi-isomorphisme, nous pouvons employer le fait que, pour
un objet $y'$ de $\mathcal{Y}'$ au-dessus d'un schéma $V$, il existe une équivalence
$$
(\Sch/V)_{fppf} \times_{g(y'), \mathcal{Y}} \mathcal{X}
=
(\Sch/V)_{fppf} \times_{y', \mathcal{Y}'}
(\mathcal{Y}' \times_\mathcal{Y} \mathcal{X})
$$
Nous concluons que le morphisme induit
$g^{-1}R^if_*\mathcal{F} \to R^if'_*(g')^{-1}\mathcal{F}$
est un isomorphisme d'après le
Lemma \ref{lemma-pushforward-restriction}.
\end{proof}












\section{Comparaison}
\label{section-compare}

\noindent
Dans cette section, nous réunissons quelques résultats comparant la cohomologie définie
au moyen des champs et celle définie au moyen des espaces algébriques.

\begin{lemma}
\label{lemma-compare-injectives}
Soit $S$ un schéma. Soit $\mathcal{X}$ un champ algébrique sur $S$
représentable par l'espace algébrique $F$.
\begin{enumerate}
\item Si $\mathcal{I}$ est injectif dans $\textit{Ab}(\mathcal{X}_\etale)$, alors
$\mathcal{I}|_{F_\etale}$ est injectif dans $\textit{Ab}(F_\etale)$,
\item Si $\mathcal{I}^\bullet$ est un complexe K-injectif dans
$\textit{Ab}(\mathcal{X}_\etale)$, then $\mathcal{I}^\bullet|_{F_\etale}$
est un complexe K-injectif dans $\textit{Ab}(F_\etale)$.
\end{enumerate}
Le même énoncé ne vaut pas pour les modules.
\end{lemma}

\begin{proof}
Cela résulte formellement du fait que le foncteur de restriction
$\pi_{F, *} = i_F^{-1}$ (voir Lemme \ref{lemma-compare})
est adjoint à droite du foncteur exact $\pi_F^{-1}$, voir
Homology, Lemma \ref{homology-lemma-adjoint-preserve-injectives} and
Derived Categories, Lemma \ref{derived-lemma-adjoint-preserve-K-injectives}.
Pour voir que le lemme ne vaut pas pour les modules, nous renvoyons le
lecteur à Étale Cohomology, Lemma
\ref{etale-cohomology-lemma-compare-injectives}.
\end{proof}

\begin{lemma}
\label{lemma-compare-morphism-cohomology}
Soit $S$ un schéma. Soit $f : \mathcal{X} \to \mathcal{Y}$ un morphisme
de champs algébriques sur $S$. Supposons $\mathcal{X}$ et $\mathcal{Y}$
représentables par les espaces algébriques $F$ et $G$. Notons $f : F \to G$ le
morphisme induit d'espaces algébriques.
\begin{enumerate}
\item Pour tout $\mathcal{F} \in \textit{Ab}(\mathcal{X}_\etale)$,
nous avons
$$
(Rf_*\mathcal{F})|_{G_\etale} =
Rf_{small, *}(\mathcal{F}|_{F_\etale})
$$
dans $D(G_\etale)$.
\item Pour tout objet $\mathcal{F}$ de
$\textit{Mod}(\mathcal{X}_\etale, \mathcal{O}_\mathcal{X})$
nous avons
$$
(Rf_*\mathcal{F})|_{G_\etale} =
Rf_{small, *}(\mathcal{F}|_{F_\etale})
$$
dans $D(\mathcal{O}_G)$.
\end{enumerate}
\end{lemma}

\begin{proof}
L'assertion (1) résulte immédiatement du
Lemma \ref{lemma-compare-injectives}
et de (\ref{equation-compare-big-small}),
après choix d'une résolution injective de $\mathcal{F}$.

\medskip\noindent
L'assertion (2) se démontre comme suit. Dans le Lemme \ref{lemma-compare-morphism},
nous avons vu que $\pi_G \circ f = f_{small} \circ \pi_F$ en tant que morphismes
de sites annelés. Nous obtenons donc
$R\pi_{G, *} \circ Rf_* = Rf_{small, *} \circ R\pi_{F, *}$
d'après Cohomology on Sites, Lemma
\ref{sites-cohomology-lemma-derived-pushforward-composition}.
Comme les foncteurs de restriction $\pi_{F, *}$ et $\pi_{G, *}$
sont exacts, nous concluons.
\end{proof}

\begin{lemma}
\label{lemma-compare-representable-morphism-cohomology}
Soit $S$ un schéma. Considérons un carré $2$-produit fibré
$$
\xymatrix{
\mathcal{X}' \ar[r]_{g'} \ar[d]_{f'} & \mathcal{X} \ar[d]^f \\
\mathcal{Y}' \ar[r]^g & \mathcal{Y}
}
$$
de champs algébriques sur $S$. Supposons $f$ représentable par des espaces
algébriques et $\mathcal{Y}'$ représentable par un espace algébrique $G'$.
Alors $\mathcal{X}'$ est représentable par un espace algébrique $F'$ et,
en notant $f' : F' \to G'$ le morphisme induit d'espaces algébriques,
nous avons
$$
g^{-1}(Rf_*\mathcal{F})|_{G'_\etale} =
Rf'_{small, *}((g')^{-1}\mathcal{F}|_{F'_\etale})
$$
pour tout $\mathcal{F}$ dans $\textit{Ab}(\mathcal{X}_\etale)$
ou dans
$\textit{Mod}(\mathcal{X}_\etale, \mathcal{O}_\mathcal{X})$
\end{lemma}

\begin{proof}
Cela résulte formellement de la combinaison des
Lemmes \ref{lemma-base-change-higher-direct-images} et
\ref{lemma-compare-morphism-cohomology}.
\end{proof}











\section{Changement de topologie}
\label{section-change-topology}

\noindent
Voici un lemme technique affirmant que la cohomologie fppf
d'un faisceau localement quasi-cohérent est égale à sa cohomologie
étale, pourvu que les morphismes de comparaison soient des isomorphismes
pour les morphismes de $\mathcal{X}$ au-dessus de morphismes plats.

\begin{lemma}
\label{lemma-lqc-flat-base-change-fppf-sheaf}
Soit $S$ un schéma. Soit $\mathcal{X}$ un champ algébrique sur $S$.
Soit $\mathcal{F}$ un préfaisceau de $\mathcal{O}_\mathcal{X}$-modules.
Supposons que
\begin{enumerate}
\item[(a)] $\mathcal{F}$ soit localement quasi-cohérent, et
\item[(b)] pour tout morphisme $\varphi : x \to y$ de $\mathcal{X}$ au-dessus
d'un morphisme de schémas $f : U \to V$ plat et
localement de présentation finie, le morphisme de comparaison
$c_\varphi : f_{small}^*\mathcal{F}|_{V_\etale} \to
\mathcal{F}|_{U_\etale}$ de
(\ref{equation-comparison-modules}) soit un isomorphisme.
\end{enumerate}
Alors $\mathcal{F}$ est un faisceau pour la topologie fppf.
\end{lemma}

\begin{proof}
Soit $\{x_i \to x\}$ un recouvrement fppf de $\mathcal{X}$ au-dessus du
recouvrement fppf $\{f_i : U_i \to U\}$ de schémas sur $S$.
Par hypothèse, la restriction $\mathcal{G} = \mathcal{F}|_{U_\etale}$
est quasi-cohérente et les morphismes de comparaison
$f_{i, small}^*\mathcal{G} \to \mathcal{F}|_{U_{i, \etale}}$
sont des isomorphismes. La condition de faisceau pour $\mathcal{F}$
et le recouvrement $\{x_i \to x\}$ équivaut donc à celle de
$\mathcal{G}^a$ sur $(\Sch/U)_{fppf}$ pour le recouvrement $\{U_i \to U\}$,
laquelle est satisfaite d'après
Descent, Lemma \ref{descent-lemma-sheaf-condition-holds}.
\end{proof}

\begin{lemma}
\label{lemma-compare-fppf-etale}
Soit $S$ un schéma. Soit $\mathcal{X}$ un champ algébrique sur $S$.
Soit $\mathcal{F}$ un $\mathcal{O}_\mathcal{X}$-module préfaisceau tel que
\begin{enumerate}
\item[(a)] $\mathcal{F}$ soit localement quasi-cohérent, et
\item[(b)] pour tout morphisme $\varphi : x \to y$ de $\mathcal{X}$ au-dessus
d'un morphisme de schémas $f : U \to V$ plat et
localement de présentation finie, le morphisme de comparaison
$c_\varphi : f_{small}^*\mathcal{F}|_{V_\etale} \to
\mathcal{F}|_{U_\etale}$ de
(\ref{equation-comparison-modules}) soit un isomorphisme.
\end{enumerate}
Alors $\mathcal{F}$ est un $\mathcal{O}_\mathcal{X}$-module et
les assertions suivantes valent :
\begin{enumerate}
\item Si $\epsilon : \mathcal{X}_{fppf} \to \mathcal{X}_\etale$
est le morphisme de comparaison, alors
$R\epsilon_*\mathcal{F} = \epsilon_*\mathcal{F}$.
\item Les groupes de cohomologie $H^p_{fppf}(\mathcal{X}, \mathcal{F})$ sont égaux
aux groupes de cohomologie calculés pour la topologie étale sur $\mathcal{X}$.
Il en va de même des groupes $H^p_{fppf}(x, \mathcal{F})$ et des
versions dérivées $R\Gamma(\mathcal{X}, \mathcal{F})$ et
$R\Gamma(x, \mathcal{F})$.
\item Si $f : \mathcal{X} \to \mathcal{Y}$ est un $1$-morphisme de
catégories fibrées en groupoïdes sur $(\Sch/S)_{fppf}$, alors
$R^if_*\mathcal{F}$ est égal à la faisceautisation fppf de l'image
directe supérieure calculée en cohomologie étale.
Il en va de même pour l'image réciproque dérivée.
\end{enumerate}
\end{lemma}

\begin{proof}
L'assertion que $\mathcal{F}$ est un $\mathcal{O}_\mathcal{X}$-module
résulte du
Lemma \ref{lemma-lqc-flat-base-change-fppf-sheaf}.
Remarquons que $\epsilon$ est un morphisme de sites donné par le foncteur
identité de $\mathcal{X}$. Le faisceau $R^p\epsilon_*\mathcal{F}$ est donc
le faisceau associé au préfaisceau
$x \mapsto H^p_{fppf}(x, \mathcal{F})$, see
Cohomology on Sites, Lemma \ref{sites-cohomology-lemma-higher-direct-images}.
Pour prouver (1), il suffit de montrer que
$H^p_{fppf}(x, \mathcal{F}) = 0$ pour $p > 0$
chaque fois que $x$ est au-dessus d'un schéma affine $U$. D'après le
Lemma \ref{lemma-cohomology-restriction}
nous avons
$H^p_{fppf}(x, \mathcal{F}) = H^p((\Sch/U)_{fppf}, x^{-1}\mathcal{F})$.
En combinant
Descent, Lemma \ref{descent-lemma-quasi-coherent-and-flat-base-change}
avec Cohomology of Schemes, Lemma
\ref{coherent-lemma-quasi-coherent-affine-cohomology-zero}
nous voyons que ces groupes de cohomologie sont nuls.

\medskip\noindent
Nous avons vu ci-dessus que $\epsilon_*\mathcal{F}$ et $\mathcal{F}$ sont les
faisceaux sur $\mathcal{X}_\etale$ et $\mathcal{X}_{fppf}$
correspondant au même préfaisceau sur $\mathcal{X}$ (et cela vaut plus
généralement pour tout faisceau fppf sur $\mathcal{X}$).
Nous identifions souvent abusivement $\mathcal{F}$ et $\epsilon_*\mathcal{F}$ ;
c'est en ce sens qu'il faut comprendre les assertions (2) et (3) du
lemme. Ainsi (2) résulte formellement de (1) et de la suite spectrale de
Leray, voir
Cohomology on Sites, Lemma \ref{sites-cohomology-lemma-apply-Leray}.

\medskip\noindent
Enfin, prouvons (3). Le faisceau $R^if_*\mathcal{F}$
(resp.\ $Rf_{\etale, *}\mathcal{F}$)
est le faisceau associé au préfaisceau
$$
y \longmapsto
H^i_\tau\Big((\Sch/V)_{fppf} \times_{y, \mathcal{Y}} \mathcal{X},
\ \text{pr}^{-1}\mathcal{F}\Big)
$$
où $\tau$ vaut $fppf$ (resp.\ $\etale$), voir le
Lemma \ref{lemma-pushforward-restriction}.
Remarquons que $\text{pr}^{-1}\mathcal{F}$ satisfait aussi les propriétés (a) et (b)
(d'après les Lemmes \ref{lemma-pullback-lqc} et \ref{lemma-comparison}) ;
ces deux préfaisceaux sont donc égaux par (2).
Cela implique immédiatement (3).
\end{proof}

\noindent
Nous emploierons le lemme suivant pour comparer la cohomologie étale des faisceaux
sur les champs algébriques à la cohomologie du topos lisse-étale.

\begin{lemma}
\label{lemma-cohomology-on-subcategory}
Soit $S$ un schéma. Soit $\mathcal{X}$ un champ algébrique sur $S$.
Soit $\tau = \etale$ (resp.\ $\tau = fppf$). Soit
$\mathcal{X}' \subset \mathcal{X}$ une sous-catégorie pleine possédant les
propriétés suivantes :
\begin{enumerate}
\item si $x \to x'$ est un morphisme de $\mathcal{X}$ au-dessus d'un morphisme
lisse (resp.\ plat et localement de présentation finie) de
schémas et $x' \in \Ob(\mathcal{X}')$, alors $x \in \Ob(\mathcal{X}')$, et
\item il existe un objet $x \in \Ob(\mathcal{X}')$ au-dessus
d'un schéma $U$ tel que le $1$-morphisme associé
$x : (\Sch/U)_{fppf} \to \mathcal{X}$ soit lisse et surjectif.
\end{enumerate}
On obtient un site $\mathcal{X}'_\tau$ en déclarant qu'un recouvrement de $\mathcal{X}'$
est toute famille de morphismes $\{x_i \to x\}$ dans $\mathcal{X}'$ qui est un
recouvrement dans $\mathcal{X}_\tau$. Alors le foncteur d'inclusion
$\mathcal{X}' \to \mathcal{X}_\tau$ est pleinement fidèle, cocontinu et
continu, et définit donc un morphisme de topos
$$
g : \Sh(\mathcal{X}'_\tau) \longrightarrow \Sh(\mathcal{X}_\tau)
$$
et $H^p(\mathcal{X}'_\tau, g^{-1}\mathcal{F}) =
H^p(\mathcal{X}_\tau, \mathcal{F})$ pour tout $p \geq 0$ et tout
$\mathcal{F} \in \textit{Ab}(\mathcal{X}_\tau)$.
\end{lemma}

\begin{proof}
L'hypothèse (1) implique que, si $\{x_i \to x\}$ est un recouvrement
de $\mathcal{X}_\tau$ et $x \in \Ob(\mathcal{X}')$, alors
$x_i \in \Ob(\mathcal{X}')$. Ainsi $\mathcal{X}' \to \mathcal{X}$
est continu et cocontinu, puisque les recouvrements des objets de
$\mathcal{X}'_\tau$ coïncident avec leurs recouvrements comme objets de
$\mathcal{X}_\tau$. Nous obtenons le morphisme $g$, et le foncteur
$g^{-1}$ s'identifie au foncteur de restriction, voir
Sites, Lemma \ref{sites-lemma-when-shriek}.

\medskip\noindent
En particulier, si $\{x_i \to x\}$ est un recouvrement dans $\mathcal{X}'_\tau$,
alors, pour tout faisceau abélien $\mathcal{F}$ sur $\mathcal{X}$,
$$
\check H^p(\{x_i \to x\}, g^{-1}\mathcal{F}) =
\check H^p(\{x_i \to x\}, \mathcal{F})
$$
Ainsi, si $\mathcal{I}$ est un faisceau abélien injectif sur $\mathcal{X}_\tau$,
les groupes de cohomologie de {\v C}ech supérieurs sont nuls
(Cohomology on Sites,
Lemma \ref{sites-cohomology-lemma-injective-trivial-cech}).
Ainsi $H^p(x, g^{-1}\mathcal{I}) = 0$ pour tout objet $x$
de $\mathcal{X}'$
(Cohomology on Sites,
Lemma \ref{sites-cohomology-lemma-cech-vanish-collection}).
Autrement dit, les faisceaux abéliens injectifs sur $\mathcal{X}_\tau$
sont acycliques à droite pour le foncteur $H^0(x, g^{-1}-)$.
Il s'ensuit que $H^p(x, g^{-1}\mathcal{F}) = H^p(x, \mathcal{F})$
pour tout $\mathcal{F} \in \textit{Ab}(\mathcal{X})$ et tout
$x \in \Ob(\mathcal{X}')$.

\medskip\noindent
Choisissons un objet $x \in \mathcal{X}'$ au-dessus d'un schéma $U$
comme dans l'hypothèse (2). En particulier, $\mathcal{X}/x \to \mathcal{X}$
est un morphisme de champs algébriques représentable par des espaces algébriques,
surjectif et lisse. (Remarquons que $\mathcal{X}/x$ est équivalent à
$(\Sch/U)_{fppf}$, see Lemma \ref{lemma-localizing}.)
Le morphisme de faisceaux
$$
h_x \longrightarrow *
$$
dans $\Sh(\mathcal{X}_\tau)$ est surjectif. En effet, pour tout objet $x'$
de $\mathcal{X}$, il existe un $\tau$-recouvrement $\{x'_i \to x'\}$
tel qu'il existe des morphismes $x'_i \to x$, voir
Lemma \ref{lemma-surjective-flat-locally-finite-presentation}.
Comme $g$ est exact, le morphisme de faisceaux
$$
g^{-1}h_x \longrightarrow * = g^{-1}*
$$
dans $\Sh(\mathcal{X}'_\tau)$ est aussi surjectif. Soit $h_{x, n}$ le
produit de $(n + 1)$ facteurs $h_x \times \ldots \times h_x$.
Nous avons alors les suites spectrales
\begin{equation}
\label{equation-spectral-sequence-one}
E_1^{p, q} = H^q(h_{x, p}, \mathcal{F}) \Rightarrow
H^{p + q}(\mathcal{X}_\tau, \mathcal{F})
\end{equation}
et
\begin{equation}
\label{equation-spectral-sequence-two}
E_1^{p, q} = H^q(g^{-1}h_{x, p}, g^{-1}\mathcal{F}) \Rightarrow
H^{p + q}(\mathcal{X}'_\tau, g^{-1}\mathcal{F})
\end{equation}
voir Cohomology on Sites,
Lemma \ref{sites-cohomology-lemma-cech-to-cohomology-sheaf-sets}.

\medskip\noindent
Cas I : $\mathcal{X}$ admet un objet final $x$ qui appartient aussi à
$\mathcal{X}'$. Ce cas résulte immédiatement de la discussion
du deuxième paragraphe ci-dessus.

\medskip\noindent
Cas II : $\mathcal{X}$ est représentable par un espace algébrique $F$.
Dans ce cas, les faisceaux $h_{x, n}$ sont représentables par un
objet $x_n$ de $\mathcal{X}$. (En effet, si $\mathcal{S}_F = \mathcal{X}$
et si $x : U \to F$ est l'objet donné, alors $h_{x, n}$ est représenté
par l'objet $U \times_F \ldots \times_F U \to F$ de $\mathcal{S}_F$.)
Il s'ensuit que $H^q(h_{x, p}, \mathcal{F}) = H^q(x_p, \mathcal{F})$.
Les morphismes $x_n \to x$ sont au-dessus de morphismes lisses de schémas, donc
$x_n \in \mathcal{X}'$ pour tout $n$. Ainsi
$H^q(g^{-1}h_{x, p}, g^{-1}\mathcal{F}) = H^q(x_p, g^{-1}\mathcal{F})$.
Ainsi, dans les deux suites spectrales
(\ref{equation-spectral-sequence-one}) et
(\ref{equation-spectral-sequence-two}) ci-dessus, les termes $E_1^{p, q}$ coïncident
d'après la discussion du deuxième paragraphe. Le lemme en résulte aussi
dans le Cas II.

\medskip\noindent
Cas III : $\mathcal{X}$ est un champ algébrique. Nous affirmons que, dans ce cas,
les groupes de cohomologie $H^q(h_{x, p}, \mathcal{F})$ et
$H^q(g^{-1}h_{x, n}, g^{-1}\mathcal{F})$ coïncident d'après le Cas II.
Une fois ceci prouvé, le résultat suivra comme précédemment.

\medskip\noindent
Considérons en effet la catégorie $\mathcal{X}/h_{x, n}$, voir
Sites, Lemma \ref{sites-lemma-localize-topos-site}.
Comme $h_{x, n}$ est le produit de $(n + 1)$ facteurs $h_x$, un
objet de cette catégorie est un $(n + 2)$-uplet
$(y, s_0, \ldots, s_n)$, où $y$ est un objet de $\mathcal{X}$ et chaque
$s_i : y \to x$ est un morphisme de $\mathcal{X}$.
C'est une catégorie sur $(\Sch/S)_{fppf}$. Il existe une équivalence
$$
\mathcal{X}/h_{x, n}
\longrightarrow
(\Sch/U)_{fppf} \times_\mathcal{X} \ldots \times_\mathcal{X} (\Sch/U)_{fppf}
=: \mathcal{U}_n
$$
sur $(\Sch/S)_{fppf}$. En effet, si $x : (\Sch/U)_{fppf} \to \mathcal{X}$ désigne aussi
le $1$-morphisme associé à $x$ et
$p : \mathcal{X} \to (\Sch/S)_{fppf}$ le foncteur structural,
alors nous pouvons voir $(y, s_0, \ldots, s_n)$ comme
$(y, f_0, \ldots, f_n, \alpha_0, \ldots, \alpha_n)$
où $y$ est un objet de $\mathcal{X}$, $f_i : p(y) \to p(x)$ un
morphisme de schémas et $\alpha_i : y \to x(f_i)$ un isomorphisme.
La catégorie des $(2n+3)$-uplets
$(y, f_0, \ldots, f_n, \alpha_0, \ldots, \alpha_n)$
est une réalisation du produit fibré de $(n + 1)$ facteurs $\mathcal{U}_n$
de champs algébriques affiché ci-dessus, comme discuté dans la
Section \ref{section-cech}.
D'après Cohomology on Sites, Lemma
\ref{sites-cohomology-lemma-cohomology-on-sheaf-sets}
nous avons
$$
H^p(\mathcal{U}_n, \mathcal{F}|_{\mathcal{U}_n}) =
H^p(\mathcal{X}/h_{x, n}, \mathcal{F}|_{\mathcal{X}/h_{x, n}}) =
H^p(h_{x, n}, \mathcal{F}).
$$
Enfin, examinons l'analogue « primé ». La catégorie
$\mathcal{X}'/h_{x, n}$ correspond, par l'équivalence ci-dessus,
à la sous-catégorie pleine $\mathcal{U}'_n \subset \mathcal{U}_n$
formée des uplets
$(y, f_0, \ldots, f_n, \alpha_0, \ldots, \alpha_n)$
avec $y \in \mathcal{X}'$. La propriété (1) de l'énoncé du
lemme vaut donc bien
pour l'inclusion $\mathcal{U}'_n \subset \mathcal{U}_n$.
Pour la propriété (2), choisissons un objet $\xi = (y, s_0, \ldots, s_n)$
au-dessus d'un schéma $W$ tel que $(\Sch/W)_{fppf} \to \mathcal{U}_n$
soit lisse et surjectif (c'est possible puisque $\mathcal{U}_n$ est
un champ algébrique). Alors
$(\Sch/W)_{fppf} \to \mathcal{U}_n \to (\Sch/U)_{fppf}$
est lisse comme composée de changements de base du morphisme
$x : (\Sch/U)_{fppf} \to \mathcal{X}$, voir
Algebraic Stacks, Lemmas
\ref{algebraic-lemma-base-change-representable-transformations-property} and
\ref{algebraic-lemma-composition-representable-transformations-property}.
Ainsi l'axiome (1) pour $\mathcal{X}$ implique que $y$ est un objet de
$\mathcal{X}'$, donc $\xi$ est un objet de $\mathcal{U}'_n$.
En employant encore
$$
H^p(\mathcal{U}'_n, \mathcal{F}|_{\mathcal{U}'_n}) =
H^p(\mathcal{X}'/h_{x, n}, \mathcal{F}|_{\mathcal{X}'/h_{x, n}}) =
H^p(g^{-1}h_{x, n}, g^{-1}\mathcal{F}).
$$
nous pouvons maintenant appliquer le Cas II à
$\mathcal{U}'_n \subset \mathcal{U}_n$
et conclure.
\end{proof}








\section{Restriction aux affines}
\label{section-alternative}

\noindent
Dans cette section, pour une catégorie $\mathcal{X}$ fibrée en groupoïdes sur
$(\Sch/S)_{fppf}$, nous considérons la sous-catégorie pleine $\mathcal{X}_{affine}$
de $\mathcal{X}$ formée des objets $x$ au-dessus de schémas affines $U$.
Nous verrons que, pour toute topologie $\tau$ plus fine que la topologie de Zariski,
les catégories de faisceaux sur $\mathcal{X}$ et sur $\mathcal{X}_{affine, \tau}$
coïncident.

\begin{definition}
\label{definition-alternative}
Soit $p : \mathcal{X} \to (\Sch/S)_{fppf}$ une catégorie fibrée en groupoïdes.
Le {\it site affine associé} est la sous-catégorie pleine
$\mathcal{X}_{affine}$ de $\mathcal{X}$ dont les objets sont les
$x \in \Ob(\mathcal{X})$ au-dessus d'un schéma affine $U$.
La topologie sur $\mathcal{X}_{affine}$ sera la topologie chaotique,
c'est-à-dire que les faisceaux sur $\mathcal{X}_{affine}$ sont les préfaisceaux.
\end{definition}

\noindent
Ainsi le foncteur $p : \mathcal{X} \to (\Sch/S)_{fppf}$ se restreint en un foncteur
$$
p : \mathcal{X}_{affine} \longrightarrow (\textit{Aff}/S)_{fppf}
$$
où la notation du membre de droite est celle introduite dans
Topologies, Definition \ref{topologies-definition-big-small-fppf}.
Il est clair que $\mathcal{X}_{affine}$ est fibrée en groupoïdes sur
$(\textit{Aff}/S)_{fppf}$. Il s'ensuit que $\mathcal{X}_{affine}$
hérite de topologies de Zariski, étale, lisse, syntomique et fppf
à partir de $(\textit{Aff}/S)_{Zar}$, $(\textit{Aff}/S)_\etale$,
$(\textit{Aff}/S)_{smooth}$, $(\textit{Aff}/S)_{syntomic}$, and
$(\textit{Aff}/S)_{fppf}$, see
Stacks, Definition \ref{stacks-definition-topology-inherited}.

\begin{definition}
\label{definition-alternative-inherited-topologies}
Soit $p : \mathcal{X} \to (\Sch/S)_{fppf}$ une catégorie fibrée en groupoïdes.
\begin{enumerate}
\item Le {\it site affine de Zariski associé} $\mathcal{X}_{affine, Zar}$
est la structure de site sur $\mathcal{X}_{affine}$ héritée de
$(\textit{Aff}/S)_{Zar}$.
\item Le {\it site affine étale associé} $\mathcal{X}_{affine, \etale}$
est la structure de site sur $\mathcal{X}_{affine}$ héritée de
$(\textit{Aff}/S)_\etale$.
\item Le {\it site affine lisse associé}
$\mathcal{X}_{affine, smooth}$
est la structure de site sur $\mathcal{X}_{affine}$ héritée de
$(\textit{Aff}/S)_{smooth}$.
\item Le {\it site affine syntomique associé} $\mathcal{X}_{affine, syntomic}$
est la structure de site sur $\mathcal{X}_{affine}$ héritée de
$(\textit{Aff}/S)_{syntomic}$.
\item Le {\it site affine fppf associé} $\mathcal{X}_{affine, fppf}$
est la structure de site sur $\mathcal{X}_{affine}$ héritée de
$(\textit{Aff}/S)_{fppf}$.
\end{enumerate}
\end{definition}

\noindent
Cette définition a un sens d'après la discussion ci-dessus. Pour chaque
$\tau \in \{Zariski, \etale, smooth, syntomic, fppf\}$, une
famille de morphismes $\{x_i \to x\}_{i \in I}$ de but fixé
dans $\mathcal{X}_{affine}$ est un recouvrement de $\mathcal{X}_{affine, \tau}$
si et seulement si la famille $\{p(x_i) \to p(x)\}_{i \in I}$
de morphismes de schémas affines est un $\tau$-recouvrement standard au sens de
Topologies, Definitions
\ref{topologies-definition-standard-Zariski},
\ref{topologies-definition-standard-etale},
\ref{topologies-definition-standard-smooth},
\ref{topologies-definition-standard-syntomic}, and
\ref{topologies-definition-standard-fppf}.

\begin{lemma}
\label{lemma-alternative}
Soit $p : \mathcal{X} \to (\Sch/S)_{fppf}$ une catégorie fibrée en groupoïdes.
Soit $\tau \in \{Zariski, \etale, smooth, syntomic, fppf\}$. Le foncteur
$\mathcal{X}_{affine, \tau} \to \mathcal{X}_\tau$ est un foncteur
cocontinu spécial. Il induit donc une équivalence de topos
de $\Sh(\mathcal{X}_{affine, \tau})$ vers $\Sh(\mathcal{X}_\tau)$.
\end{lemma}

\begin{proof}
Omis. Indication : la preuve est exactement la même que celle de
Topologies, Lemmas
\ref{topologies-lemma-affine-big-site-Zariski},
\ref{topologies-lemma-affine-big-site-etale},
\ref{topologies-lemma-affine-big-site-smooth},
\ref{topologies-lemma-affine-big-site-syntomic}, and
\ref{topologies-lemma-affine-big-site-fppf}.
\end{proof}

\noindent
Soit $p : \mathcal{X} \to (\Sch/S)_{fppf}$ une catégorie fibrée en groupoïdes.
Notons $\mathcal{O}$ la restriction de $\mathcal{O}_\mathcal{X}$
à $\mathcal{X}_{affine}$. Alors $\mathcal{O}$ est un faisceau pour les topologies
de Zariski, étale, lisse, syntomique et fppf sur
$\mathcal{X}_{affine}$. De plus, l'équivalence de topos du
Lemme \ref{lemma-alternative} s'étend en une équivalence
\begin{equation}
\label{equation-alternative-ringed}
(\Sh(\mathcal{X}_{affine, \tau}), \mathcal{O})
\longrightarrow
(\Sh(\mathcal{X}_\tau), \mathcal{O}_\mathcal{X})
\end{equation}
de topos annelés pour $\tau \in \{Zariski, \etale, smooth, syntomic, fppf\}$.









\section{Modules quasi-cohérents et affines}
\label{section-alternative-quasi-coherent}

\noindent
Soit $p : \mathcal{X} \to (\Sch/S)_{fppf}$ une catégorie fibrée en groupoïdes.
Dans la Section \ref{section-alternative}, nous lui avons associé un site annelé
$(\mathcal{X}_{affine}, \mathcal{O})$.

\begin{lemma}
\label{lemma-quasi-coherent-alternative}
Soit $p : \mathcal{X} \to (\Sch/S)_{fppf}$ une catégorie fibrée en groupoïdes.
Soit $\mathcal{F}$ un $\mathcal{O}$-module sur $\mathcal{X}_{affine}$.
Les conditions suivantes sont équivalentes :
\begin{enumerate}
\item pour tout morphisme $x \to x'$ de $\mathcal{X}_{affine}$, le morphisme
$\mathcal{F}(x') \otimes_{\mathcal{O}(x')} \mathcal{O}(x) \to \mathcal{F}(x)$
est un isomorphisme,
\item $\mathcal{F}$ est un module quasi-cohérent sur
$(\mathcal{X}_{affine}, \mathcal{O})$ au sens de
Modules on Sites, Definition \ref{sites-modules-definition-site-local},
\item $\mathcal{F}$ est un faisceau pour la topologie de Zariski sur
$\mathcal{X}_{affine}$ et un module quasi-cohérent sur
$(\mathcal{X}_{affine, Zar}, \mathcal{O})$ au sens de
Modules on Sites, Definition \ref{sites-modules-definition-site-local},
\item comme en (3) pour la topologie étale,
\item comme en (3) pour la topologie lisse,
\item comme en (3) pour la topologie syntomique,
\item comme en (3) pour la topologie fppf, et
\item $\mathcal{F}$ correspond à un module quasi-cohérent
sur $\mathcal{X}$ par l'équivalence (\ref{equation-alternative-ringed}).
\end{enumerate}
\end{lemma}

\begin{proof}
Pour interpréter (2), rappelons que $\mathcal{X}_{affine}$ est le site
obtenu en munissant la catégorie $\mathcal{X}_{affine}$ de la topologie
chaotique (Définition \ref{definition-alternative}) ; un faisceau
de $\mathcal{O}$-modules $\mathcal{F}$ est donc la même chose qu'un
préfaisceau de $\mathcal{O}$-modules. Les conditions (1) et (2)
sont équivalentes d'après Modules on Sites, Lemma
\ref{sites-modules-lemma-chaotic-quasi-coherent}.
Remarquons que, pour $\tau \in \{Zariski, \etale, smooth, syntomic, fppf\}$,
le préfaisceau $\mathcal{F}$ est un $\tau$-faisceau si et seulement si,
pour tout $x \in \Ob(\mathcal{X}_{affine})$, sa restriction
à $\mathcal{X}_{affine}/x$ est un $\tau$-faisceau. Posons $U = p(x)$.
Comme dans la discussion de la Section \ref{section-restriction},
l'objet $x$ de $\mathcal{X}_{affine}$ induit une équivalence
$\mathcal{X}_{affine, \etale}/x \to (\textit{Aff}/U)_\etale$ de sites.
Ainsi, l'équivalence de (1) avec (3) -- (7)
résulte de Descent, Lemma \ref{descent-lemma-quasi-coherent-alternative},
appliqué à chacun de ces sites. L'équivalence de (8) et (7)
résulte immédiatement du fait qu'« être quasi-cohérent » est une
propriété intrinsèque des faisceaux de modules, voir
Modules on Sites, Section \ref{sites-modules-section-intrinsic}
\end{proof}

\begin{lemma}
\label{lemma-locally-quasi-coherent-alternative}
Soit $p : \mathcal{X} \to (\Sch/S)_{fppf}$ une catégorie fibrée en groupoïdes.
Soit $\mathcal{F}$ un $\mathcal{O}$-module sur $\mathcal{X}_{affine}$.
Les conditions suivantes sont équivalentes :
\begin{enumerate}
\item pour tout morphisme $x \to x'$ de $\mathcal{X}_{affine}$ tel que
$p(x) \to p(x')$ soit un morphisme étale (de schémas affines), le morphisme
$\mathcal{F}(x') \otimes_{\mathcal{O}(x')} \mathcal{O}(x) \to \mathcal{F}(x)$
est un isomorphisme,
\item $\mathcal{F}$ est un faisceau pour la topologie étale sur
$\mathcal{X}_{affine}$ et, pour tout objet $x$ de $\mathcal{X}_{affine}$,
la restriction $x^*\mathcal{F}|_{U_{affine, \etale}}$ est quasi-cohérente,
où $U = p(x)$,
\item $\mathcal{F}$ correspond à un module localement quasi-cohérent
sur $\mathcal{X}$ par l'équivalence (\ref{equation-alternative-ringed})
pour la topologie étale.
\end{enumerate}
\end{lemma}

\begin{proof}
Pour interpréter la condition (2), rappelons que $U_{affine, \etale}$
est la sous-catégorie pleine de $U_\etale$ formée des objets affines, voir
Topologies, Definition \ref{topologies-definition-big-small-etale}.
Comme dans la discussion de la Section \ref{section-restriction},
l'objet $x$ de $\mathcal{X}_{affine}$ induit une équivalence
$\mathcal{X}_{affine, \etale}/x \to (\textit{Aff}/U)_\etale$ de sites.
Alors $x^*\mathcal{F}$ est le faisceau de modules sur $(\textit{Aff}/U)_\etale$
correspondant à la restriction
$\mathcal{F}|_{\mathcal{X}_{affine, \etale}/x}$.
Enfin, au moyen du foncteur d'inclusion continu et cocontinu
$U_{affine, \etale} \to (\textit{Aff}/U)_\etale$, nous pouvons encore
restreindre et obtenir $x^*\mathcal{F}|_{U_{affine, \etale}}$.

\medskip\noindent
L'équivalence de (1) et (2) résulte des remarques ci-dessus et de
Descent, Lemma \ref{descent-lemma-quasi-coherent-alternative-small}
appliqué à la restriction de $\mathcal{F}$ à $U_{affine, \etale}$
pour tout objet $x$ de $\mathcal{X}$ au-dessus d'un schéma affine $U$.
L'équivalence de (2) et (3) résulte immédiatement des définitions
et du fait que les modules quasi-cohérents sur $U_{affine, \etale}$
et $U_\etale$ se correspondent (encore d'après
Descent, Lemma \ref{descent-lemma-quasi-coherent-alternative-small}
par exemple).
\end{proof}













\section{Objets quasi-cohérents de la catégorie dérivée}
\label{section-QC}

\noindent
Depuis des décennies, les géomètres algébristes considèrent des invariants de
foncteurs non représentables $X$ (à valeurs dans les ensembles ou les groupoïdes) sur $\Sch/S$.
Par exemple,
avant l'invention de la notion de champ, Mumford a défini \cite{mumford_picard}
le groupoïde de Picard $\Pic(X)$ du foncteur de modules $X$ des courbes elliptiques
comme la $2$-limite $Pic(U)$ sur la catégorie de tous les schémas $U$
munis d'un morphisme vers $X$ (c'est-à-dire d'une famille de courbes elliptiques).
De même, Beilinson-Drinfeld ont défini \cite{BVGD} la catégorie
$\QCoh(X)$ d'un ind-schéma $X = \colim X_i$ comme la $2$-limite
$\lim \QCoh(X_i)$.
Cette stratégie suffit à définir des invariants $1$-catégoriques
comme $\QCoh(-)$, mais non des invariants catégoriques dérivés
(tels que la catégorie dérivée quasi-cohérente), car les $2$-limites
de catégories triangulées se comportent mal.
Avec l'avènement des techniques de catégories supérieures
et de la géométrie algébrique dérivée, ce problème se
résout élégamment : on peut définir la $\infty$-catégorie dérivée
quasi-cohérente $\mathcal{D}_{qc}(X)$ du foncteur $X$
comme la limite $\lim \mathcal{D}_{qc}(U)$, où $U$
parcourt les affines dérivés sur $X$ (voir \cite{lurie-thesis}).

\medskip\noindent
Le but de cette section est d'associer une catégorie triangulée
$\mathit{QC}(X)$ à un foncteur $X$ (à valeurs dans les ensembles ou les groupoïdes)
comme ci-dessus. En fait, la construction vaut pour toute catégorie
$p : \mathcal{X} \to (\Sch/S)_{fppf}$ fibrée en groupoïdes
(et pas seulement pour les catégories scindées). Dans les bons cas, la catégorie $\mathit{QC}(\mathcal{X})$
coïncide avec la catégorie homotopique de
$\mathcal{D}_{qc}(\mathcal{X})$,
mais expliquer cette comparaison dépasse le cadre de ce document.
Les principales propriétés de la construction sont les suivantes :
\begin{enumerate}
\item[(a)] $\mathit{QC}(\mathcal{X})$ est par construction une sous-catégorie pleine de
$D(\mathcal{X}_{affine}, \mathcal{O})$,
\item[(b)] $\mathit{QC}(\mathcal{X})$ coïncide avec $D_\QCoh(\mathcal{O}_X)$
lorsque $\mathcal{X}$ est représentable par l'espace algébrique $X$,
\item[(c)] $\mathit{QC}(\mathcal{X})$ coïncide avec
$D_\QCoh(\mathcal{O}_\mathcal{X})$ lorsque $\mathcal{X}$ est un champ algébrique,
\item[(d)] lorsque $X = \text{Spf}(A)$ est un espace algébrique formel affine
associé à un anneau noethérien $A$ muni de la topologie $I$-adique
pour un idéal $I$, la catégorie triangulée
$\mathit{QC}(X)$ coïncide avec la sous-catégorie pleine
$D_{comp}(A, I) \subset D(A)$ des objets dérivés complets.
\end{enumerate}
Ces résultats sont démontrés dans la
Proposition \ref{proposition-QC-compare},
Derived Categories of Stacks, Proposition
\ref{stacks-perfect-proposition-QC-compare}, and
Proposition \ref{proposition-QC-derived-complete}.

\medskip\noindent
Pour motiver la définition précise de $\mathit{QC}(\mathcal{X})$,
nous renvoyons le lecteur à la caractérisation, dans le
Lemma \ref{lemma-quasi-coherent-alternative},
des modules quasi-cohérents sur $\mathcal{X}$ comme préfaisceaux de
$\mathcal{O}$-modules sur $\mathcal{X}_{affine}$ qui
satisfont une propriété de changement de base.

\begin{definition}
\label{definition-QC}
Soit $p : \mathcal{X} \to (\Sch/S)_{fppf}$ une catégorie fibrée en groupoïdes.
Soit $\mathcal{O}$ le faisceau d'anneaux sur $\mathcal{X}_{affine}$ introduit
dans la Section \ref{section-alternative}. Nous définissons la
{\it catégorie triangulée des objets quasi-cohérents de la catégorie dérivée}
par la formule
$$
\mathit{QC}(\mathcal{X}) = \mathit{QC}(\mathcal{X}_{affine}, \mathcal{O})
$$
où le membre de droite est défini dans
Cohomology on Sites, Definition \ref{sites-cohomology-definition-cartesian}.
\end{definition}

\noindent
Cela a un sens, car $\mathcal{X}_{affine}$ est une catégorie considérée
comme un site en la munissant de la topologie chaotique, et
$\mathcal{O}$ est un faisceau d'anneaux sur cette catégorie, exactement comme requis dans
Cohomology on Sites, Definition \ref{sites-cohomology-definition-cartesian}.

\medskip\noindent
La relation entre cette définition et la catégorie des modules quasi-cohérents
sur $\mathcal{X}$ n'est pas si claire en général ! Supposons par exemple
que $M$ soit un objet de $\mathit{QC}(\mathcal{X})$. Les faisceaux de cohomologie
$H^i(M)$ de $M$ sont alors des (pré)faisceaux de $\mathcal{O}$-modules sur
$\mathcal{X}_{affine}$, mais ils ne sont pas quasi-cohérents en général. Toutefois, le dernier
faisceau de cohomologie non nul est quasi-cohérent.

\begin{lemma}
\label{lemma-QC-bounded-above}
Dans la situation de la Définition \ref{definition-QC}, supposons que
$M$ soit un objet de $\mathit{QC}(\mathcal{X})$ et que $b \in \mathbf{Z}$
vérifie $H^i(M) = 0$ pour tout $i > b$.
Alors $H^b(M)$ est un module quasi-cohérent sur
$(\mathcal{X}_{affine}, \mathcal{O})$, see
Lemma \ref{lemma-quasi-coherent-alternative}.
\end{lemma}

\begin{proof}
Cas particulier de
Cohomology on Sites, Lemma \ref{sites-cohomology-lemma-QC-bounded-above}.
\end{proof}

\begin{lemma}
\label{lemma-QC-compare}
Soit $S$ un schéma. Soit $\mathcal{X} \to (\Sch/S)_{fppf}$ une catégorie
fibrée en groupoïdes. Le morphisme de comparaison
$\epsilon : \mathcal{X}_{affine, \etale} \to \mathcal{X}_{affine}$
satisfait les hypothèses et les conclusions de Cohomology on Sites, Lemma
\ref{sites-cohomology-lemma-cartesion-plus-topology}.
\end{lemma}

\begin{proof}
L'hypothèse (1) vaut par définition de $\mathcal{X}_{affine}$.
Pour la condition (2), nous employons le fait que, pour $x \in \Ob(\mathcal{X})$
au-dessus du schéma affine $U = p(x)$, il existe une équivalence
$\mathcal{X}_{affine, \etale}/x = (\textit{Aff}/U)_\etale$ compatible
aux faisceaux structuraux ; voir la
discussion de la Section \ref{section-restriction}.
Il suffit donc de montrer ceci : pour un schéma affine $U = \Spec(R)$
et un complexe de $R$-modules $M^\bullet$, la cohomologie totale
du complexe de modules sur $(\textit{Aff}/U)_\etale$
associé à $M^\bullet$ est quasi-isomorphe à $M^\bullet$.
Cela résulte de la combinaison de :
Derived Categories of Schemes, Lemma
\ref{perfect-lemma-affine-compare-bounded} (cohomologie totale des complexes
de modules sur les affines pour la topologie de Zariski),
Derived Categories of Spaces, Remark
\ref{spaces-perfect-remark-match-total-direct-images}
(coïncidence de la cohomologie totale pour les petites
topologies de Zariski et étale des
complexes quasi-cohérents de modules), et
Étale Cohomology, Lemma \ref{etale-cohomology-lemma-compare-cohomology}
(pour voir que la cohomologie étale d'un complexe de modules
sur le grand site étale d'un schéma
se calcule après restriction au petit site étale).
\end{proof}

\noindent
Si nous appliquons la définition lorsque notre catégorie fibrée en groupoïdes
$\mathcal{X}$ est représentable par un espace algébrique $X$, nous
retrouvons $D_\QCoh(\mathcal{O}_X)$. Nous énoncerons et démontrerons plus loin le
résultat analogue pour les champs algébriques (insérer ici une référence future).

\begin{proposition}
\label{proposition-QC-compare}
Soit $S$ un schéma. Soit $\mathcal{X} \to (\Sch/S)_{fppf}$ une catégorie
fibrée en groupoïdes. Supposons $\mathcal{X}$ représentable par un espace
algébrique $X$. Alors $\mathit{QC}(\mathcal{X})$ est canoniquement équivalente à
$D_\QCoh(\mathcal{O}_X)$.
\end{proposition}

\begin{proof}
Notons $X_{affine}$ la catégorie des schémas affines étales sur $X$,
munie de la topologie chaotique et de son faisceau structural $\mathcal{O}_X$, voir
Derived Categories of Spaces, Section \ref{spaces-perfect-section-QC}.
Le foncteur $u : X_\etale \to \mathcal{X}_\etale$ du Lemme \ref{lemma-compare}
induit un foncteur $X_{affine} \to \mathcal{X}_{affine}$.
Celui-ci est compatible aux faisceaux structuraux et fournit un foncteur
$$
G :
\mathit{QC}(\mathcal{X}) = \mathit{QC}(\mathcal{X}_{affine}, \mathcal{O})
\longrightarrow
\mathit{QC}(X_{affine}, \mathcal{O}_X)
$$
Voir
Cohomology on Sites, Lemma \ref{sites-cohomology-lemma-cartesian-functorial}.
D'après Derived Categories of Spaces, Lemma
\ref{spaces-perfect-lemma-DQCoh-alternative-small}
la catégorie triangulée $\mathit{QC}(X_{affine}, \mathcal{O}_X)$
est équivalente à $D_\QCoh(\mathcal{O}_X)$. Il suffit donc de
prouver que $G$ est une équivalence.

\medskip\noindent
Considérons les morphismes de comparaison plats
$\epsilon_\mathcal{X} : \mathcal{X}_{affine, \etale} \to \mathcal{X}_{affine}$
et $\epsilon_X : X_{affine, \etale} \to X_{affine}$
de sites annelés. Le Lemme \ref{lemma-QC-compare} et
(la preuve de) Derived Categories of Spaces, Lemma
\ref{spaces-perfect-lemma-DQCoh-alternative-small}
montrent que les foncteurs
$\epsilon_\mathcal{X}^*$ et $\epsilon_X^*$
identifient $\mathit{QC}(\mathcal{X}_{affine}, \mathcal{O})$
et $\mathit{QC}(X_{affine}, \mathcal{O}_X)$
à des sous-catégories
$Q_\mathcal{X} \subset D(\mathcal{X}_{affine, \etale}, \mathcal{O})$
et
$Q_X \subset D(X_{affine, \etale}, \mathcal{O}_X)$.
Sous ces identifications,
le foncteur $G$ du premier paragraphe est induit par le foncteur
$$
Li_X^* = R\pi_{X, *}:
D(\mathcal{X}_{affine, \etale}, \mathcal{O})
\longrightarrow
D(X_{affine, \etale}, \mathcal{O}_X)
$$
où $i_X$ et $\pi_X$ sont les morphismes du Lemme \ref{lemma-compare},
les sites étales étant remplacés par les sites affines correspondants.
Le lecteur peut montrer que ce remplacement est licite soit
en redémontrant directement le lemme pour les sites affines, soit
en employant les équivalences de topos
$\Sh(\mathcal{X}_{affine, \etale}) = \Sh(\mathcal{X}_\etale)$ et
$\Sh(X_{affine, \etale}) = \Sh(X_\etale)$.
Le lemme nous dit aussi que $Li_X^*$ admet un adjoint à gauche
$$
L\pi_X^*:
D(X_{affine, \etale}, \mathcal{O}_X)
\longrightarrow
D(\mathcal{X}_{affine, \etale}, \mathcal{O})
$$
et, de plus, $Li_X^* \circ L\pi_X^* = \text{id}$
puisque $\pi_X \circ i_X$ est l'identité. Il suffit donc
de montrer que (a) $L\pi_X^*$ envoie $Q_X$ dans $Q_\mathcal{X}$
et que (b) le noyau de $Li_X^*$ est nul.
Voir Derived Categories, Lemma
\ref{derived-lemma-fully-faithful-adjoint-kernel-zero}.

\medskip\noindent
Preuve de (a). D'après Derived Categories of Spaces, Lemma
\ref{spaces-perfect-lemma-DQCoh-alternative-small}
nous avons $Q_X = D_\QCoh(X_{affine, \etale}, \mathcal{O}_X)$.
Soit $K$ un objet de $Q_X$. Soit $x$ un objet de
$\mathcal{X}_{affine, \etale}$ au-dessus du schéma
affine $U = p(x)$. Notons $f : U \to X$ le morphisme correspondant
à $x$. Alors
$$
R\Gamma(x, L\pi_X^*K) = R\Gamma(U, Lf^*K)
$$
Cela résulte de la transitivité des images réciproques ; voir la discussion de la
Section \ref{section-restriction-algebraic-spaces}.
Supposons ensuite que $x \to x'$ soit un morphisme de $\mathcal{X}_{affine, \etale}$
au-dessus du morphisme $h : U \to U'$ de schémas affines.
Comme précédemment, notons $f : U \to X$ et $f' : U' \to X$ les morphismes
correspondant à $x$ et $x'$, de sorte que $f = f' \circ h$.
Alors
\begin{align*}
R\Gamma(x, L\pi_X^*K)
& =
R\Gamma(U, Lf^*K) \\
& =
R\Gamma(U, Lh^*L(f')^*K) \\
& =
R\Gamma(U', L(f')^*K) \otimes_{\mathcal{O}(U')}^\mathbf{L} \mathcal{O}(U) \\
& =
R\Gamma(x', L\pi_X^*K) \otimes_{\mathcal{O}(x')}^\mathbf{L} \mathcal{O}(x)
\end{align*}
et nous obtenons donc (a) par la note de l'énoncé de
Cohomology on Sites, Lemma
\ref{sites-cohomology-lemma-cartesion-plus-topology}. La troisième égalité est
Derived Categories of Schemes, Lemma
\ref{perfect-lemma-quasi-coherence-pullback}.

\medskip\noindent
Preuve de (b). Soit $M$ un objet de $Q_\mathcal{X}$
tel que $Li_X^*M = 0$. Soit $x'$ un objet de $\mathcal{X}_{affine, \etale}$
au-dessus du schéma affine $U' = p(x')$, et supposons que le morphisme
correspondant $f' : U' \to X$ soit étale. Alors $f' : U' \to X$ est un objet de
$X_{affine, \etale}$, et la condition $Li_X^*M = 0$ implique
$M|_{U'_\etale} = 0$. En particulier, $R\Gamma(x', M) = 0$.
Cependant, pour tout objet $x$ du site
$\mathcal{X}_{affine, \etale}$, il existe un recouvrement $\{x_i \to x\}$
tel que, pour chaque $i$, il existe un morphisme $x_i \to x'_i$, où
$x'_i$ correspond à un objet de $X_{affine, \etale}$.
Puisque $M$ appartient à $Q_\mathcal{X}$, nous avons
$$
R\Gamma(x_i, M) =
R\Gamma(x_i', M) \otimes_{\mathcal{O}(x_i')}^\mathbf{L} \mathcal{O}(x_i) =
0
$$
et nous concluons que $M$ est nul, comme voulu.
\end{proof}

\noindent
Pour montrer que la construction produit une catégorie intéressante dans un autre
cas, énonçons et démontrons une caractérisation de
$\mathit{QC}(\text{Spf}(A))$ pour le spectre formel d'un
anneau adique noethérien $A$.

\begin{proposition}
\label{proposition-QC-derived-complete}
Soit $S$ un schéma. Soit $X = \text{Spf}(A)$, où $A$ est une
$S$-algèbre topologique adique noethérienne d'idéal de définition $I$, voir
More on Algebra, Definition \ref{more-algebra-definition-topological-ring}
et Formal Spaces, Definition
\ref{formal-spaces-definition-affine-formal-spectrum}.
Soit $p : \mathcal{X} \to (\Sch/S)_{fppf}$
la catégorie fibrée en ensembles associée au foncteur $X$, voir
Categories, Example \ref{categories-example-presheaf}.
Alors $\mathit{QC}(\mathcal{X})$ est canoniquement équivalente à la
catégorie $D_{comp}(A, I)$ des objets de $D(A)$ qui sont
dérivés complets pour $I$.
\end{proposition}

\begin{proof}
Rappelons que $X = \colim \Spec(A/I^n)$ comme faisceau fppf.
Un objet de $\mathcal{X}_{affine}$ est la même chose
qu'un schéma affine $U = \Spec(R)$ muni d'un morphisme $f : U \to X$.
D'après Formal Spaces, Lemma \ref{formal-spaces-lemma-factor-through-thickening},
il existe $n \geq 1$ tel que
$f$ se factorise par le monomorphisme $\Spec(A/I^n) \to X$.
Considérons la sous-catégorie pleine
$\mathcal{C} \subset \mathcal{X}_{affine}$ formée des
objets $\Spec(A/I^n) \to X$.
D'après les remarques précédentes et Differential Graded Sheaves, Lemma
\ref{sdga-lemma-cartesian-cofinal-subcategory}
la restriction à $\mathcal{C}$ est une équivalence exacte
$\mathit{QC}(\mathcal{X}) \to
\mathit{QC}(\mathcal{C}, \mathcal{O}|_\mathcal{C})$.
Pour simplifier, supposons que
$I^n \not = I^{n + 1}$ pour tout $n \geq 1$.
Alors $(\mathcal{C}, \mathcal{O}|_\mathcal{C})$ est isomorphe,
comme site annelé, au site annelé $(\mathbf{N}, (A/I^n))$, voir
Differential Graded Sheaves, Section
\ref{sdga-section-qc-inverse-system}.
Nous concluons donc par Differential Graded Sheaves, Proposition
\ref{sdga-proposition-derived-complete}.
\end{proof}

\noindent
Le lemme suivant servira à comparer $\mathit{QC}(\mathcal{X})$
à $D_\QCoh(\mathcal{O}_\mathcal{X})$ lorsque $\mathcal{X}$ est un champ algébrique.

\begin{lemma}
\label{lemma-QC-compare-fppf}
Soit $S$ un schéma. Soit $\mathcal{X} \to (\Sch/S)_{fppf}$ une catégorie
fibrée en groupoïdes. Le morphisme de comparaison
$\epsilon : \mathcal{X}_{affine, fppf} \to \mathcal{X}_{affine}$
satisfait les hypothèses et les conclusions de Cohomology on Sites, Lemma
\ref{sites-cohomology-lemma-cartesion-plus-topology}.
\end{lemma}

\begin{proof}
La preuve est exactement la même que celle du Lemme \ref{lemma-QC-compare}.
L'hypothèse (1) vaut par définition de $\mathcal{X}_{affine}$.
Pour la condition (2), nous employons le fait que, pour $x \in \Ob(\mathcal{X})$
au-dessus du schéma affine $U = p(x)$, il existe une équivalence
$\mathcal{X}_{affine, \etale}/x = (\textit{Aff}/U)_\etale$ compatible
aux faisceaux structuraux ; voir la
discussion de la Section \ref{section-restriction}.
Il suffit donc de montrer ceci : pour un schéma affine $U = \Spec(R)$
et un complexe de $R$-modules $M^\bullet$, la cohomologie totale
du complexe de modules sur $(\textit{Aff}/U)_{fppf}$
associé à $M^\bullet$ est quasi-isomorphe à $M^\bullet$.
C'est Étale Cohomology, Lemma
\ref{etale-cohomology-lemma-cohomological-descent-complex-modules}.
\end{proof}










\begin{multicols}{2}[\section{Autres chapitres}]
\noindent
Préliminaires
\begin{enumerate}
\item \hyperref[introduction-section-phantom]{Introduction}
\item \hyperref[conventions-section-phantom]{Conventions}
\item \hyperref[sets-section-phantom]{Théorie des ensembles}
\item \hyperref[categories-section-phantom]{Catégories}
\item \hyperref[topology-section-phantom]{Topologie}
\item \hyperref[sheaves-section-phantom]{Faisceaux sur les espaces}
\item \hyperref[sites-section-phantom]{Sites et faisceaux}
\item \hyperref[stacks-section-phantom]{Champs}
\item \hyperref[fields-section-phantom]{Corps}
\item \hyperref[algebra-section-phantom]{Algèbre commutative}
\item \hyperref[brauer-section-phantom]{Groupes de Brauer}
\item \hyperref[homology-section-phantom]{Algèbre homologique}
\item \hyperref[derived-section-phantom]{Catégories dérivées}
\item \hyperref[simplicial-section-phantom]{Méthodes simpliciales}
\item \hyperref[more-algebra-section-phantom]{Compléments d'algèbre}
\item \hyperref[smoothing-section-phantom]{Lissage des morphismes d'anneaux}
\item \hyperref[modules-section-phantom]{Faisceaux de modules}
\item \hyperref[sites-modules-section-phantom]{Modules sur les sites}
\item \hyperref[injectives-section-phantom]{Objets injectifs}
\item \hyperref[cohomology-section-phantom]{Cohomologie des faisceaux}
\item \hyperref[sites-cohomology-section-phantom]{Cohomologie sur les sites}
\item \hyperref[dga-section-phantom]{Algèbre différentielle graduée}
\item \hyperref[dpa-section-phantom]{Algèbre à puissances divisées}
\item \hyperref[sdga-section-phantom]{Faisceaux différentiels gradués}
\item \hyperref[hypercovering-section-phantom]{Hyperrecouvrements}
\end{enumerate}
Schémas
\begin{enumerate}
\setcounter{enumi}{25}
\item \hyperref[schemes-section-phantom]{Schémas}
\item \hyperref[constructions-section-phantom]{Constructions de schémas}
\item \hyperref[properties-section-phantom]{Propriétés des schémas}
\item \hyperref[morphisms-section-phantom]{Morphismes de schémas}
\item \hyperref[coherent-section-phantom]{Cohomologie des schémas}
\item \hyperref[divisors-section-phantom]{Diviseurs}
\item \hyperref[limits-section-phantom]{Limites de schémas}
\item \hyperref[varieties-section-phantom]{Variétés}
\item \hyperref[topologies-section-phantom]{Topologies sur les schémas}
\item \hyperref[descent-section-phantom]{Descente}
\item \hyperref[perfect-section-phantom]{Catégories dérivées de schémas}
\item \hyperref[more-morphisms-section-phantom]{Compléments sur les morphismes}
\item \hyperref[flat-section-phantom]{Compléments sur la platitude}
\item \hyperref[groupoids-section-phantom]{Schémas en groupoïdes}
\item \hyperref[more-groupoids-section-phantom]{Compléments sur les schémas en groupoïdes}
\item \hyperref[etale-section-phantom]{Morphismes étales de schémas}
\end{enumerate}
Thèmes de théorie des schémas
\begin{enumerate}
\setcounter{enumi}{41}
\item \hyperref[chow-section-phantom]{Homologie de Chow}
\item \hyperref[intersection-section-phantom]{Théorie de l'intersection}
\item \hyperref[pic-section-phantom]{Schémas de Picard des courbes}
\item \hyperref[weil-section-phantom]{Théories cohomologiques de Weil}
\item \hyperref[adequate-section-phantom]{Modules adéquats}
\item \hyperref[dualizing-section-phantom]{Complexes dualisants}
\item \hyperref[duality-section-phantom]{Dualité pour les schémas}
\item \hyperref[discriminant-section-phantom]{Discriminants et différentes}
\item \hyperref[derham-section-phantom]{Cohomologie de de Rham}
\item \hyperref[local-cohomology-section-phantom]{Cohomologie locale}
\item \hyperref[algebraization-section-phantom]{Géométrie algébrique et formelle}
\item \hyperref[curves-section-phantom]{Courbes algébriques}
\item \hyperref[resolve-section-phantom]{Résolution des surfaces}
\item \hyperref[models-section-phantom]{Réduction semi-stable}
\item \hyperref[functors-section-phantom]{Foncteurs et morphismes}
\item \hyperref[equiv-section-phantom]{Catégories dérivées de variétés}
\item \hyperref[pione-section-phantom]{Groupes fondamentaux de schémas}
\item \hyperref[etale-cohomology-section-phantom]{Cohomologie étale}
\item \hyperref[crystalline-section-phantom]{Cohomologie cristalline}
\item \hyperref[proetale-section-phantom]{Cohomologie pro-étale}
\item \hyperref[relative-cycles-section-phantom]{Cycles relatifs}
\item \hyperref[more-etale-section-phantom]{Compléments de cohomologie étale}
\item \hyperref[trace-section-phantom]{La formule des traces}
\end{enumerate}
Espaces algébriques
\begin{enumerate}
\setcounter{enumi}{64}
\item \hyperref[spaces-section-phantom]{Espaces algébriques}
\item \hyperref[spaces-properties-section-phantom]{Propriétés des espaces algébriques}
\item \hyperref[spaces-morphisms-section-phantom]{Morphismes d'espaces algébriques}
\item \hyperref[decent-spaces-section-phantom]{Espaces algébriques décents}
\item \hyperref[spaces-cohomology-section-phantom]{Cohomologie des espaces algébriques}
\item \hyperref[spaces-limits-section-phantom]{Limites d'espaces algébriques}
\item \hyperref[spaces-divisors-section-phantom]{Diviseurs sur les espaces algébriques}
\item \hyperref[spaces-over-fields-section-phantom]{Espaces algébriques sur des corps}
\item \hyperref[spaces-topologies-section-phantom]{Topologies sur les espaces algébriques}
\item \hyperref[spaces-descent-section-phantom]{Descente et espaces algébriques}
\item \hyperref[spaces-perfect-section-phantom]{Catégories dérivées d'espaces}
\item \hyperref[spaces-more-morphisms-section-phantom]{Compléments sur les morphismes d'espaces}
\item \hyperref[spaces-flat-section-phantom]{Platitude sur les espaces algébriques}
\item \hyperref[spaces-groupoids-section-phantom]{Groupoïdes dans les espaces algébriques}
\item \hyperref[spaces-more-groupoids-section-phantom]{Compléments sur les groupoïdes dans les espaces}
\item \hyperref[bootstrap-section-phantom]{Amorçage}
\item \hyperref[spaces-pushouts-section-phantom]{Sommes amalgamées d'espaces algébriques}
\end{enumerate}
Thèmes de géométrie
\begin{enumerate}
\setcounter{enumi}{81}
\item \hyperref[spaces-chow-section-phantom]{Groupes de Chow des espaces}
\item \hyperref[groupoids-quotients-section-phantom]{Quotients de groupoïdes}
\item \hyperref[spaces-more-cohomology-section-phantom]{Compléments sur la cohomologie des espaces}
\item \hyperref[spaces-simplicial-section-phantom]{Espaces simpliciaux}
\item \hyperref[spaces-duality-section-phantom]{Dualité pour les espaces}
\item \hyperref[formal-spaces-section-phantom]{Espaces algébriques formels}
\item \hyperref[restricted-section-phantom]{Algébrisation des espaces formels}
\item \hyperref[spaces-resolve-section-phantom]{Résolution des surfaces revisitée}
\end{enumerate}
Théorie des déformations
\begin{enumerate}
\setcounter{enumi}{89}
\item \hyperref[formal-defos-section-phantom]{Théorie formelle des déformations}
\item \hyperref[defos-section-phantom]{Théorie des déformations}
\item \hyperref[cotangent-section-phantom]{Le complexe cotangent}
\item \hyperref[examples-defos-section-phantom]{Problèmes de déformation}
\end{enumerate}
Champs algébriques
\begin{enumerate}
\setcounter{enumi}{93}
\item \hyperref[algebraic-section-phantom]{Champs algébriques}
\item \hyperref[examples-stacks-section-phantom]{Exemples de champs}
\item \hyperref[stacks-sheaves-section-phantom]{Faisceaux sur les champs algébriques}
\item \hyperref[criteria-section-phantom]{Critères de représentabilité}
\item \hyperref[artin-section-phantom]{Axiomes d'Artin}
\item \hyperref[quot-section-phantom]{Espaces de Quot et de Hilbert}
\item \hyperref[stacks-properties-section-phantom]{Propriétés des champs algébriques}
\item \hyperref[stacks-morphisms-section-phantom]{Morphismes de champs algébriques}
\item \hyperref[stacks-limits-section-phantom]{Limites de champs algébriques}
\item \hyperref[stacks-cohomology-section-phantom]{Cohomologie des champs algébriques}
\item \hyperref[stacks-perfect-section-phantom]{Catégories dérivées de champs}
\item \hyperref[stacks-introduction-section-phantom]{Introduction aux champs algébriques}
\item \hyperref[stacks-more-morphisms-section-phantom]{Compléments sur les morphismes de champs}
\item \hyperref[stacks-geometry-section-phantom]{La géométrie des champs}
\end{enumerate}
Thèmes de théorie des espaces de modules
\begin{enumerate}
\setcounter{enumi}{107}
\item \hyperref[moduli-section-phantom]{Champs de modules}
\item \hyperref[moduli-curves-section-phantom]{Espaces de modules de courbes}
\end{enumerate}
Divers
\begin{enumerate}
\setcounter{enumi}{109}
\item \hyperref[examples-section-phantom]{Exemples}
\item \hyperref[exercises-section-phantom]{Exercices}
\item \hyperref[guide-section-phantom]{Guide bibliographique}
\item \hyperref[desirables-section-phantom]{Desiderata}
\item \hyperref[coding-section-phantom]{Style de programmation}
\item \hyperref[obsolete-section-phantom]{Obsolète}
\item \hyperref[fdl-section-phantom]{Licence de documentation libre GNU}
\item \hyperref[index-section-phantom]{Index généré automatiquement}
\end{enumerate}
\end{multicols}


\bibliography{my}
\bibliographystyle{amsalpha}

\end{document}
